\section{Integral interpolation on a Frobenius cycle}
\label{sec:interpolation}
\hypertarget{source-interpolation}{}

Section~\ref{sec:word-division} requires bounds on the directions lost in integral local tests. On a Frobenius cycle the losses come both from specializing a weight and from evaluating at neighboring lattices. We first bound the exceptional
evaluations over an interval of integer weights by bounding the
average number of Smith exponents at or above a given cutoff. We then examine simultaneous
evaluation on the neighboring lattices. These are the two
rank tests for the exceptional directions mentioned in
Section~\ref{sec:family-window}. For the first test, fix a prime \(p\) and a square matrix polynomial \(A(T)\) over a valued extension of \(\mathbb Q_p\). Lemma~\ref{lem:smith-average} varies the integral
argument \(t_j\) and counts, with multiplicity, the invariant factors
of \(A(t_j)\) whose normalized \(p\)-adic order reaches a prescribed threshold \(E\ge1\); a zero
diagonal entry has order \(+\infty\).

\subsection{Counting Smith factors of large valuation}

At an individual weight even the determinant may vanish. Averaging the number of large Smith factors over the sampled weights gives a rank estimate that allows these singular specializations.

\begin{lemma}[Average number of large Smith exponents]\label{lem:smith-average}
Let \(A(T)\) be an \(l\)-by-\(l\) matrix over a valued extension
of \(\mathbb Q_p\), with the valuation normalized by \(v(p)=1\).
Suppose its entries are linear combinations
of \(\binom{T}{i}\), \(0\le i\le\delta\), with integral
coefficients. Suppose its determinant is nonzero and write
\[
 \det A(T)=C\prod_{i=1}^d(T-\zeta_i),\qquad
 v(C)+d/(p-1)\le l L .
\]
Measure the roots in a splitting field with its extended valuation.
Take \(H\ge1\) points \(t_j=t_0+j b\), \(0\le j<H\), where
\(t_0,b\in\mathbb Z\) and \(p\nmid b\). Choose an integer
\(h\ge0\) such that \(H\le p^h\) and \(\delta<p^h\).
For \(E\ge1\) let \(b_j\) be the number of Smith diagonal
entries of \(A(t_j)\) of valuation at least \(E\), including
zeros. Then
\begin{equation}\label{eq:smith-average}
 \frac1H\sum_j b_j\le
 l\left(\frac{\delta}{H}+
        \frac{L+\delta h/H}{E}\right).
\end{equation}

\end{lemma}

\begin{proof}
\proofstep{Integral translation of the binomial basis.}
For a polynomial expanded at \(t\), write
\(v_u(\sum a_i(T-t)^i)=\min_i(v(a_i)+iu)\).
It is multiplicative. In each factor take the smallest index
attaining the minimum. At the sum of the two smallest indices
the product has exactly one least-valuation contribution.
For any \(t\in\mathbb Z\), the coefficients in \(Z\) of
\(A(t+p^hZ)\) are integral. Indeed the binomial addition
identity reduces this to integrality of \(\binom ti\) and
of the polynomials
\[
 \binom{p^hZ}{i}
 =(-1)^{i-1}\frac{p^hZ}{i}
                \prod_{j=1}^{i-1}(1-p^hZ/j)
                  \qquad(1\le i\le\delta).
\]
Each denominator has valuation at most \(h\).
\proofstep{A determinant valuation at each evaluation.}
A finite matrix over the valuation ring admits the diagonal reduction used here even if the value group is nondiscrete. If the matrix is nonzero, choose a nonzero entry of least valuation among its entries. Permute it to the upper left. It divides every entry, so integral row and column additions clear its column and row. The new smaller matrix still has every entry divisible by the first pivot. Iteration yields
\[
  P A(t_j)Q=\operatorname{diag}(a_1,\ldots,a_r,0,\ldots,0),
  \qquad a_1\mid\cdots\mid a_r ,
\]
where \(P,Q\) have unit determinants. Use the same constant \(P,Q\) on \(A(t_j+p^hZ)\).
These preserve coefficient integrality and determinant valuation.


Every column corresponding to one of the \(b_j\) entries has
constant term of valuation at least \(E\). In the expansion in
\(Z\) its nonconstant terms have weighted valuation at least
\(E\) at weight \(E\). All other columns have nonnegative
weighted valuation. The determinant expansion now gives
\[
 E b_j \le v_{h+E}(\det A)
       =v(C)+\sum_{i=1}^d\min(h+E,v(t_j-\zeta_i)).
\]
This inequality also applies when \(t_j=\zeta_i\).

\proofstep{Counting close evaluations.}
For each root and each real \(0<u\le h\), the indices with
\(v(t_j-\zeta_i)\ge u\) lie in one residue class modulo
\(p^{\lceil u\rceil}\), since differences of two points have
valuation \(v(j-j')\). There are at most
\(Hp^{-\lceil u\rceil}+1\) such indices. For \(h<u\le h+E\)
there is at most one. Integrating these counts on the two
ranges shows that the sum over \(j\) of each truncated root
valuation is at most \(H/(p-1)+h+E\). Negative valuations only
decrease this sum. Thus
\[
 E\sum_j b_j\le
 H\big(v(C)+d/(p-1)\big)+d(h+E).
\]
The degree bound \(d\le l\delta\) proves
\eqref{eq:smith-average}.
\end{proof}

In the segment Gram calculation the
parameter \(T\) is the degree of the appended covector. If the
binomial degree of its Gram matrix is \(2S\), the length
is \(H\), \(h\le\log_p(8a+r)+1\), and the leading coefficient
bound gives \(L\le2\kappa\), the resulting discarded fraction
is bounded by
\[
 \frac{2S}{H}+
 \frac{2\kappa+2S(\log_p(8a+r)+1)/H}{E}.
\] The following calculation identifies the polynomial and
its leading coefficient for the segment form.

\subsection{The Gram matrix when a weight varies}
\hypertarget{source-tied-gram}{}

To apply Lemma~\ref{lem:smith-average}, we need a matrix polynomial with a controlled determinant. A decomposition by the degree of the final vector produces it from the Gram form.

Let \(f\) be a polynomial in the \(b\) vector or covector variables of a fixed initial segment, each having \(r\) coordinates. Their degrees are \(o_i\), the
total degree is \(S\), and the last letter \(w\) is a vector of
degree \(o\). The others may be vectors or covectors. The Fock
pairing is bilinear, with orthogonal monomials and squared
monomial norms \(\prod_\alpha m_\alpha!\). Choose a coordinate
\(\rho\) and add a covector \(\eta\) of variable degree \(n\ge S\).
Define the substitution \(T_\eta\) on each vector \(v\) and covector \(u\). For every \(i\ne\rho\) its coordinate rules are
\[
 \begin{aligned} (T_\eta v)_\rho&=v\cdot\eta, & (T_\eta v)_i&=\eta_\rho v_i,\\ (T_\eta u)_\rho&=\eta_\rho u_\rho, & (T_\eta u)_i&=\eta_\rho u_i-u_\rho\eta_i. \end{aligned}
\]
Set \(f_n=\eta_\rho^{n-S}T_\eta f\). Evaluating at the \(\rho\)-th coordinate vector recovers
\(f\). In particular the substitution is an integral injection.

Let \(U\) multiply by \(w\cdot\eta\), let
\(D=\sum_i\partial_{w_i}\partial_{\eta_i}\), and let \(z\)
be the nonnegative integer decoration of this adjacent pair.
Write \(B=o+n+r+z-1\) and define
\[
 \Phi_z=\sum_{j=0}^o\frac{(-1)^j U^jD^j}
                           {(j!)^2\binom Bj},\qquad
 \Gamma_n(f,g)=\binom Bo\,
       \frac{\langle\Phi_z f_n,g_n\rangle}
            {(n-S)!\,S!\prod_i o_i!}.
\]
Multiplication \(U\) and contraction \(D\) are adjoint for the Fock pairing. The scalar
in this formula is used before the bound on small Smith factors.

\begin{lemma}[Polynomial Gram pairing at a tied pair]
\label{lem:tied-gram-polynomial}
For fixed inputs and fixed decoration \(z\), \(\Gamma_n\) is polynomial
in \(n\) of degree at most \(S+o\). For integral inputs its entries
after multiplication by \(p^u\) lie in the integral binomial basis, with
\[
 u=\left\lceil ((b+3)r+1)
                    (1+\log_p(\max(1,3S)))\right\rceil .
\]
If a basis \((f_i)\) is adapted to the filtration by order in \(w_\rho\),
write \(e_i\) for the order of \(f_i\), \(l\) for its size and
\(d_0=\sum_i(S+o-e_i)\). Suppose each normalized graded determinant
has valuation at most \(v_*\) times its rank. For \(S,z\le2a\) the
determinant of \((p^u\Gamma_n(f_i,f_j))\) has degree \(d_0\) and leading
coefficient \(C\) with
\[
 v(C)+d_0/(p-1)\le
     l\big(u+v_*+3(1+\log_p(8a+r))\big).
\]
\end{lemma}

\begin{proof}
\proofstep{Degree in the varied weight.}
Each entry is polynomial in \(n\) of degree at most \(S+o\).
In the \(j\)-th term move \(U^j\) by adjunction.
The ratio \(\binom Bo/\binom Bj=\binom{B-j}{o-j}/\binom oj\)
is polynomial of degree \(o-j\). To see the remaining polynomial,
put \(m=n-S\). Before applying \(D^j\), the two occurrences of
\(\eta_\rho\) in a monomial pairing have powers \(m+t,m+t'\),
where \(0\le t,t'\le S\). If the differentiation orders there
are \(a,a'\), the residual powers must agree. Their contribution,
including the Fock factorial and division by \(m!\), is
\[
   \frac{(m+t)!}{m!}(m+t')_{\underline{a'}},
       \qquad t-a=t'-a' .
\]
All remaining factors are independent of \(m\). This polynomial
has degree \(\le S+j\), and its falling factorial also gives the
correct zero when the differentiation order is too large.

\proofstep{Integrality in the binomial basis.}
For integral inputs we can clear the binomial-basis denominators
with logarithmic loss. Evaluate at the consecutive integers
\(S\le n\le S+S+o\le3S\). In the expansion of \(D^j/j!\)
let \(h_i\) contractions take place at coordinate \(i\).
On a monomial with exponents \(u_i,v_i\) in \(w,\eta\), the multi-index \(h\) contributes the coefficient
\[
 \prod_i\frac{u_i^{\underline{h_i}}v_i^{\underline{h_i}}}{h_i!} =\prod_i h_i!\binom{u_i}{h_i}\binom{v_i}{h_i},\qquad \sum_i h_i=j.
\]
Use this on both sides, with multiplicities \(h_i,h'_i\).

If the resulting monomial has exponent vector \(d\), the
numerator contains \(\prod_\alpha d_\alpha!\prod_i h_i!h'_i!\).
The sum of the factorial arguments is exactly \(S+n\), the
same as for \((n-S)! S!\prod_i o_i!\). Legendre's formula
gives
\[
  0\le t/(p-1)-v(t!)\le 1+\log_p(\max(1,t)).
\]
There are at most \((b+3)r\) numerator factorials. The additional
denominator \(\binom oj\) costs at most \(1+\log_p(\max(1,o))\)
by counting carries in base \(p\). The other factor
\(\binom{B-j}{o-j}\) is integral for every \(z\). Thus the
values become integral after multiplication by \(p^u\) for
\[
 u=\left\lceil\big((b+3)r+1\big)
                \big(1+\log_p(\max(1,3S))\big)\right\rceil .
\]
Taking forward differences at \(S\) recovers the integral
coefficients in the basis \(\binom{n-S}{j}\). Translation
through an integer preserves the integral binomial basis,
so the hypothesis on the entries in \eqref{eq:smith-average}
holds with \(\delta=2S\).

\proofstep{The leading harmonic pairing.}
We compute the leading coefficient to keep the dependence on
\(S\) out of the final valuation allowance. Write \(e(f)\) for
the least exponent of \(w_\rho\) in a nonzero \(f\), and \(I_e f\)
for its terms of exponent exactly \(e\). For
\(e=e(f),e'=e(g)\) one has
\begin{align}
 \deg\Gamma_n(f,g)&\le S+o-\max(e,e'),\label{eq:tie-degree}\\
 {}[n^{S+o-e}]\Gamma_n(f,g)
   &=\frac{e!\binom{z+e}{e}}{o!\,S!\prod_i o_i!}
                   \langle I_e f,I_e g\rangle
                      \qquad(e=e').\label{eq:tie-top}
\end{align}
It suffices to prove these rational identities on monomials,
so we can use the positive real Fock metric to compute the
leading coefficient there. On bidegree \((o,n)\) the orthogonal
decomposition by powers of \(U\) has summands \(U^j\ker D\)
with harmonic factors of degrees \((o-j,n-j)\). Indeed
\((\ker D)^\perp=\operatorname{im}U\), and iteration gives the
summands. The commutator \([D,U]\) is the sum of the two Euler
operators plus \(r\), which proves orthogonality by moving powers
of \(U\) across the pairing.

It gives the eigenvalue of
\(\Phi_z\) on the \(j\)-th summand as
\[
 \sum_{l=0}^j (-1)^l\binom jl
       \binom{o+n+r-j-1}{l}/\binom Bl
      =\binom{z+j}{j}/\binom Bj .
\]
The equality is inclusion--exclusion for a sample of \(j\)
objects from \(B\) objects, of which \(o+n+r-j-1\) are forbidden.
Thus the eigenvalue is \(O(n^{-j})\) for fixed prefix and \(z\).
The image \(f_n\) belongs to \(U^e\) times the smaller bidegree,
since the substitution sends \(w_\rho\) to \(w\cdot\eta\).
Also
\[
 \|f_n\|^2/(n-S)!=O(n^S)
\]
by expanding its monomials.
Only levels \(j\ge\max(e,e')\) can pair, proving
\eqref{eq:tie-degree}. The real limits here extract coefficients
at fixed prefix. The uniform estimates used as the prefix grows
are the ensuing factorial valuations.

For \(e=e'\), write \(f_n=U^e h_n\). The polynomial \(h_n\)
is \(\eta_\rho^{n-S}\) times the substitution on
\(h=f/w_\rho^e\), and has bidegree \((o-e,n-e)\).
Use \(g_n=U^e h'_n\) in the same way. After division by \((n-S)!\), a level \(j>e\) contributes at most
\[
 O(n^{S})\,O(n^{-j})=O(n^{S-j})=O(n^{S-e-1}).
\]
Thus only level \(e\) can survive at order \(n^{S-e}\).

Let \(P\) be the harmonic projector on this smaller bidegree.
Its formula is the sum defining \(\Phi_z\), with degrees
\(o-e,n-e\) and denominator parameter \(o+n-2e+r-2\).
The same eigenvalue identity gives one on level zero and zero
on the higher levels (use the binomial identity at \(z=-1\)).
Only large \(n\) are needed in this limit. On harmonic vectors
\[
 D^eU^e=e!\,(o+n-2e+r)^{\overline e}.
\]
Multiplying by the level-\(e\) eigenvalue therefore tends to
\((e!)^2\binom{z+e}{e}\).

The remaining harmonic pairing has limit
\[
  \lim_{n\to\infty}\frac{\langle P h_n,h'_n\rangle}
                            {(n-S)!\,n^{S-e}}
       =\langle h|_{w_\rho=0},h'|_{w_\rho=0}\rangle .
\]
In the \(l\)-th term of \(P\), move the \(l\) multiplications
across as derivatives. The term before the projector coefficient,
divided by \((n-S)! n^{S-e+l}\), tends to
\(\langle \partial_{w_\rho}^l h,\partial_{w_\rho}^l h'\rangle\).
To check the powers, let \(d,d'\) count occurrences of
\(\eta_i\), \(i\ne\rho\), and \(a,a'\) the respective numbers
of derivatives falling there. A contributing pair has
\(d-a=d'-a'\). After factorial division its degree is at most
\[
 S-e+l-\underbrace{(d-a)}_{\ge0}-a-a'\le S-e+l.
\]
Equality forces \(a=a'=0\) and \(d-a=d'-a'=0\), hence \(d=d'=0\).

The coefficient of the pure
\(\eta_\rho\)-power in the substituted \(h\) is \(h\) itself.
This proves the limit term by term. The leading projector
coefficient is \((-1)^l/(l! n^l)\). On the terms of \(w_\rho\)-degree \(t\) in each argument,
moving the derivatives across the Fock form and the factor
\(1/l!\) give the scalar
\[
       \sum_{l=0}^t(-1)^l\binom tl
           =\begin{cases}1,&t=0,\\0,&t>0.\end{cases}
\]
Distinct degrees are orthogonal. Thus the limiting operator
keeps only the restrictions \(w_\rho=0\). Restoring the factor
\(w_\rho^e\) in the initial terms multiplies this pairing by
\(e!\). Finally \(\binom Bo/n^o\to1/o!\),
which proves \eqref{eq:tie-top}.

\proofstep{The leading determinant in a filtration basis.}
Choose a basis \(f_i\) adapted to the integral filtration by
vanishing order, so that the \(I_e f_i\) of order \(e\) form
a basis for the image at graded order \(e\). Such a basis
exists over a DVR by successive lifting from the free graded
images of the coefficient map. In the determinant of
\((p^u\Gamma_n(f_i,f_j))\) the degree is at most
\(d_0=\sum_i(S+o-e_i)\). For a permutation \(\pi\) the excess in vanishing orders is
\[
 \sum_i\max(e_i,e_{\pi(i)})-\sum_i e_i =\tfrac12\sum_i|e_i-e_{\pi(i)}|.
\]
Thus its determinant term can attain degree \(d_0\) only if \(e_i=e_{\pi(i)}\) throughout, by
\eqref{eq:tie-degree} and \(\sum_i e_i=\sum_i e_{\pi(i)}\).

Consequently its top coefficient is the product of the
determinants of the graded Fock forms normalized by \(\prod o_i!\),
times
\[
                  \prod_i\frac{p^u e_i!\binom{z+e_i}{e_i}}
                                         {o!\,S!}.
\]
If each normalized graded determinant has valuation at most
\(v_*\) times its rank, the top coefficient is nonzero. For
\(S,z\le2a\), Legendre's formula and the carry bound give
\[
  v(C)+d_0/(p-1)\le l\big(u+v_*+3(1+\log_p(8a+r))\big),
\]
where \(l\) is the size of the matrix. Indeed the three
linear factorial terms cancel \((S+o-e_i)/(p-1)\) in every
row, and the digit error from \(e_i!\) has favorable sign.
This supplies the polynomial and the precise leading-coefficient
input to \eqref{eq:smith-average} from a bound on the graded
form.
\end{proof}

\subsection{Denominators of the projections on types}
\hypertarget{source-graded-projectors}{}

The preceding Gram calculation gives a rational decomposition. Its use in integral interpolation also requires the denominators of the projections to the graded types. In the
first, let \(E_{ij}\) be the matrix units acting on a
finite-dimensional rational \(\mathrm{GL}_r\)-module in
characteristic zero. Suppose an integral lattice respects the
weight decomposition and is stable under
\(E_{ij}^{a}/a!\) for all \(a\) and \(i\ne j\). We use complete reducibility, highest weights and the Pieri
and branching rules for split general linear groups in
characteristic zero. Their Schur-module constructions
commute with extension of the ground field to
\(\mathbb Q_p\). \cite{fulton-harris}
The polynomial
lattices in the letters have this property by taking the
coefficients of the actions of elementary unipotent matrices.
The lattice of homomorphisms between two such modules has it
as well, by expanding the left and inverse right unipotent
actions.

\begin{lemma}[Denominators for invariant and isotypic projections]
\label{lem:integral-type-projection}
If the weights in the module are in \([-l,l]\) in each coordinate and
\(l+r-1\ge1\), write \(P\) for projection onto the invariants and \(M_0\) for the zero-weight lattice. The denominator bound is
\[
 p^{\binom r2\lfloor\log_p(l+r-1)\rfloor}\,P(M_0)\subseteq M_0.
\]
 The same multiplier
works for an isotypic projector when \(l\) bounds the coordinate weights
of the corresponding Hom representation.
\end{lemma}

\begin{proof}
\proofstep{Projecting onto invariants by the integral trace.}
On the zero weight space one can project onto the invariant
vectors with denominator at most
\[
                  p^{\binom r2\lfloor\log_p(l+r-1)\rfloor},
\]
when all coordinate weights in the module lie in \([-l,l]\).
Take \(l+r-1\ge1\). To see this, inductively suppose that
the input \(v\) is invariant in a set \(I\) of \(s\) coordinates.
Adjoin coordinate \(z\). Put \(\Omega_0=1\),
\(\Omega_{q+1}=\sum_{i\in I}E_{zi}\Omega_q E_{iz}\),
and apply
\[
   P_{I,z}=\sum_{q=0}^l c_q\Omega_q,\qquad
   c_q=\frac{(-1)^q s!}{(q+s)!q!}.
\]
Each \(\Omega_q\) commutes with \(\mathfrak{gl}_I\) by contracting
the two indices in the recursion, and preserves weights.
For \(j\in I\), commuting the matrix units gives
\[
 E_{jz}\Omega_q v=\Omega_q E_{jz}v+
                       q(q+s)\Omega_{q-1}E_{jz}v
                 \qquad(q\ge1).
\]
For example expand \(\Omega_q\) as the sum of all products of
\(q\) factors \(E_{zi}\) followed in reverse order by the
corresponding \(E_{iz}\).

In moving \(E_{jz}\) through the
first \(q\) factors each commutator is
\(E_{ji}-\delta_{ij}E_{zz}\). Moving these through the remaining
factors and using \(E_{ii'}v=E_{zz}v=0\) for \(i,i'\in I\)
gives the stated recurrence. Substituting the coefficient at the next index gives, for \(0\le q<l\),
\[
 (q+1)(q+s+1)c_{q+1}
 =(q+1)(q+s+1)\frac{(-1)^{q+1}s!}{(q+s+1)!(q+1)!}=-c_q .
\]
Thus the two contributions at each interior index cancel.
Its endpoint \(\Omega_l E_{jz}v\) vanishes, because the inner
\(l+1\) raising factors move the \(z\)-weight below \(-l\).
Thus every \(E_{jz}\) kills the result. On its zero weight
the opposite \(E_{zj}\) also kills it, by the finite-dimensional
\(\mathfrak{sl}_2\) string for \((j,z)\). Vectors already
invariant are fixed. All terms beyond the identity are in the
span of images of off-diagonal generators.

For a multi-index \(a=(a_i)_{i\in I}\) with \(\sum a_i=q\), expand
\(c_q\Omega_q\) using commutation within each of the two groups
of root operators. The coefficient of the product of their
\emph{divided} powers has absolute rational value
\(s!\prod_i a_i!/(s+q)!\).
Number the \(s\) indices in order. Cancelling successive factorials gives
\[
 \frac{(s+q)!}{s!\prod_i a_i!}=\prod_{j=1}^s\binom{s+a_1+\cdots+a_j}{a_j}.
\]
Each has valuation at most \(\lfloor\log_p(l+s)\rfloor\)
by the carry bound. Summing the denominator exponents for
\(s=0,\ldots,r-1\) proves the claim. The construction commutes
with intertwiners, since it uses only weight projection and the
matrix units. On an endomorphism space it preserves the trace.
Projection onto weight zero retains all diagonal entries, and
every subsequent correction is a sum of commutators.

\proofstep{Denominators on a highest-weight summand.}
In particular take an irreducible summand \(U\) and the module
\(\operatorname{Hom}(U,V)\), with action
\(E\phi=E_V\phi-\phi E_U\). Use the coordinate weight
bound \(l\) on differences of weights. Denote the invariant
projection just proved, preceded by weight-zero projection,
by \(P\). With an integral basis \(u_j\) of a divided-power
stable weight lattice in \(U\) the operator
\[
     Q_U(v)=(\dim U)\sum_j P(u_j^*\otimes v)(u_j)
\]
is the isotypic projector. Such a lattice is obtained by
intersecting a copy of the summand defined over \(\mathbb Q_p\) with the ambient integral
lattice. Its preservation under the divided powers follows by
intersection. On an inequivalent irreducible the invariant Hom
space is zero.

For an embedding \(i:U\to V\), naturality and
Schur's lemma give
\[
 P(u_j^*\otimes i(u))=i\circ(u_j^*(u)/\dim U)\operatorname{id}_U,
\]
since the invariant projection preserves trace on endomorphisms.
This gives \(Q_U(i(u))=i(u)\). Hence the same denominator bound
works for a sum of any chosen types. For letters of fixed total
degree \(S\), the span of coordinate weights in the full module
is at most \(S\): a vector of degree \(d\) contributes in \([0,d]\), and a covector of degree \(d\) in \([-d,0]\). The width of the total interval is therefore bounded by
\[
 \sum_{\rm vectors} d+\sum_{\rm covectors} d=S.
\]
Thus one may use \(l=S\) also for the relevant Hom spaces.
In the Fock metric \(E_{ij}\) is adjoint to \(E_{ji}\), including
the contragredient action on covectors, so distinct types are
orthogonal.

For the second estimate, fix an order \(e\) and let \(M_e\) be
the lattice of initial terms \(I_e f\). Here the integral \(f\)
are in a fixed full weight, highest for the coordinates other
than \(\rho\), and satisfy the prefix bounds on first coordinates
of the generated modules. Put \(H_e=\mathbb Q_pM_e\) and
\[
 L_e=H_e\cap\mathbb Z_p[\mathbf x].
\]
Suppose a self-adjoint
projection onto \(H_e\) has denominator at most \(p^d\) on
the monomials of the required multidegree and that
\(p^c L_e\subseteq M_e\). The graded determinant allowance
is then \(v_*=d+2c\). Indeed a primitive \(f\in M_e\) has a
coefficient of valuation \(\le c\). Otherwise
\[
 f\in p^{c+1}L_e\quad\Longrightarrow\quad f/p\in p^cL_e\subseteq M_e,
\]
contradicting primitivity.

If its exponent is \(\alpha\),
the normalized Fock pairing with \(\mathbf x^\alpha\) has valuation
at most \(c\), because the reciprocal of the monomial norm is a
product of integral multinomial coefficients. Multiply the
projected monomial by \(p^{d+c}\) to obtain an integral partner
in \(M_e\) with pairing of valuation \(\le d+2c\).
Using the row corresponding to a deepest invariant factor in
a Smith decomposition, this primitive-vector test bounds every
invariant factor by \(d+2c\), and thus the determinant valuation
by \((d+2c)\operatorname{rank}M_e\). The zero-rank convention has
determinant one.

\end{proof}

\subsection{Lifting integral initial terms}

Clearing a projector controls its image. We also need to lift prescribed leading terms into that image while keeping the prefix bounds.
Assume \(r\ge2\), take \(\rho\) last, and let \(K\) be the number
of letters preceding the final vector \(w\). A cutoff on a prefix
means that every irreducible type of that prefix occurring in
the vector has first highest-weight coordinate at most the specified
integer. If \(C\) is the cutoff and \(G\) the prefix group, the test on a nonzero polynomial \(g\) is
\[
 \max\{\mu_1:(xg)_\mu\ne0\text{ for some }x\in G\}\le C,
\]
where the subscript denotes a weight component. An infinite cutoff imposes
no condition. If \(f=w_\rho^e h\), a prefix containing \(w\)
satisfies its cutoff precisely when \(h\) satisfies the cutoff
lowered by \(e\). Indeed for a generic group translate the top
first coordinate in the translated linear form \(w_\rho\) is one.

Let \(m_1(f)\) be the largest first coordinate weight among
prefix translates of \(f\), and similarly for \(h\). On a common
Zariski-open set the highest first weight of the translated
\(w_\rho\) is one, and the component of \(h\) realizing \(m_1(h)\)
is nonzero. Their top components multiply to a nonzero polynomial,
so
\[
                  m_1(w_\rho^e h)=e+m_1(h).
\] For shorter prefixes the
factor is fixed. The group on the coordinates other than \(\rho\)
also fixes it. We may therefore divide out this power, subtract
\(e\) from the last degree and the \(\rho\)-weight, and take
restriction \(R\) at \(w_\rho=0\). Multiplication back by
\(w_\rho^e\) recovers the graded term. It scales the Fock
form on the restricted monomials by a common factor.

\begin{lemma}[Saturation of the restricted coefficient map]
\label{lem:graded-coefficient-lift}
For \(r\ge2\), with the prefix-first-weight cutoffs and
restriction at \(w_\rho=0\), the image of restriction
on the integral polynomial lattice contains the integral restricted
lattice after the multiplier in \eqref{eq:res-sat} before the highest
weight equations. After the weight and type projectors let \(d_{\rm gr}\) be the normalized valuation per rank of the graded determinant. With \(K\) the number of earlier letters,
\[
 d_{\rm gr}\le(K+3)r^2(1+\log_p(8a+r))\qquad(S\le2a).
\]

\end{lemma}

\begin{proof}
\proofstep{The lowered last degree.}
Keep the shortened prefix cutoffs fixed and vary the reduced last
degree \(t\). Write \(Y_t\) for the space satisfying these cutoffs
before restriction, with no highest-vector or weight requirements.
In the restricted space define \(X_t\) by the same earlier prefix
conditions. If there is a finite full-prefix cutoff, let \(C\) be
the least one. Impose in addition the bound \(C\) on the first
\(K\) letters under \(\mathrm{GL}_r\) and on the entire restricted
word under \(\mathrm{GL}_{r-1}\).

We have \(RY_t\subset X_t\).
Indeed, taking a prefix type and tensoring it by the last symmetric
power adds a horizontal strip, by the Pieri rule (twist by a
determinant to reduce integral highest weights to partitions).
If \(\alpha\) is the earlier type and \(\lambda\) a resulting type, the strip condition gives
\[
 \lambda_1\ge\alpha_1\ge\lambda_2.
\]
In particular a cutoff on \(\lambda_1\) also bounds \(\alpha_1\). Restriction commutes with the earlier
group, and with \(\mathrm{GL}_{r-1}\) on the whole word.
Interlacing bounds the first coordinate of the latter subgroup
types by that of the full type.

Let \(S\) be the total degree before the shortening, and set
\[
 A_0=\binom r2\lfloor\log_p(S+r-1)\rfloor,\qquad
 A_1=\binom r2\lfloor\log_p(4S+r-1)\rfloor .
\]
We claim that at the required last degree \(t\)
\begin{equation}\label{eq:res-sat}
  p^{2A_1}(X_t\cap\mathbb Z_p[\mathbf x])
                  \subset R(Y_t\cap\mathbb Z_p[\mathbf x]).
\end{equation}
With no full cutoff, or with \(C\) at least the total reduced
degree at this \(t\), the polynomial independent of \(w_\rho\)
lifts itself. For the other case use an auxiliary Lie algebra action.
Denote the surviving coordinates by \(I\), put \(E_w=\sum w_i
\partial_{w_i}\), and write \(B_{ab}\) for the matrix units on
the first \(K\) letters. Use an extra index \(*\) preceding all
of \(I\). The following operators act on \(\bigoplus_t X_t\):
\begin{align*}
 G_{ab}&=B_{ab}+w_a\partial_{w_b},&
 G_{**}&=C-E_w,\\
 G_{*a}&=\partial_{w_a},&
 G_{a*}&=w_a(C-E_w)-\sum_{b\in I} w_b B_{ab}.
\end{align*}
The relations \([\partial_i,w_j]=\delta_{ij}\) and the matrix-unit brackets for \(B\) give
\[
 \begin{aligned} {}[G_{ab},G_{c*}]&=\delta_{bc}G_{a*},\\ {}[G_{*b},G_{a*}]&=\delta_{ab}G_{**}-G_{ab},\\ {}[G_{a*},G_{b*}]&=0.\end{aligned}
\]
Together with the degrees these give the
\(\mathfrak{gl}_r\) brackets.

\proofstep{Preservation of the endpoint cutoff.}
We check the cutoff at the endpoint, where raising could otherwise
leave the space. The same formulas apply on all coordinates before
restriction, with the sum over all \(b\); restriction intertwines
them for \(a,b\in I\), since every omitted summand has a factor
\(w_\rho\). At an earlier prefix each \(B_{ab}\) splits into an
operator of the prefix group and one on disjoint letters.
Both preserve its type condition. At the full prefix the raises
transform as a vector under the coordinate group and the derivatives
as a covector. A derivative never increases first coordinate
weight. A raise with index different from the first has first
weight zero.

For the first index, denoted by 1, the identity
\[
                   G_{1*}=C w_1-\sum_b w_bG_{1b}
\]
shows that it kills the component of first weight \(C\).
On this component \(G_{11}=C\) and each \(G_{1b}\) with
\(b\ne1\) vanishes because it raises the first weight beyond
the cutoff. On lower weights the result stays below or at \(C\).
Applying this argument after every group translate proves the
stable type bound for all the raises, upstairs and downstairs.

The divided powers also preserve integral coefficients. This
is immediate for differentiation. For the raise with index \(a\),
let \(w\) be the column of variables, and let \(\mathcal B\)
be the representation on the earlier letters. On a polynomial
\(h\) valued in that representation the generating series is
\[
 (1+u w_a)^C\,
 \mathcal B\!\left(I-\frac{u e_a w^{\mathsf T}}{1+u w_a}\right)
                 h\!\left(\frac{w}{1+u w_a}\right).
\]
Let \(T(u)\) denote this substitution. The rank-one matrix product
gives \(T(u+v)=T(u)T(v)\), and \(T(0)=1,\ T'(0)=G_{a*}\).

Differentiate in \(v\) at zero and compare coefficients in \(u\):
\[
 T'(u)=T(u)G_{a*},\qquad
 (j+1)[u^{j+1}]T=([u^j]T)G_{a*},\qquad
 [u^j]T=G_{a*}^j/j! .
\] The displayed matrix and its inverse
\(I+u e_aw^{\mathsf T}\) have integral series entries, so this
also works for covectors and negative \(C\). For the operators on
\(I\) the ordinary elementary matrices give integral divided
powers. Thus the integral lattice in the direct sum of the \(X_t\),
and the sum of the images of the integral \(Y_t\), are stable
under divided powers and their integral weight decompositions.

\proofstep{Lifting the initial forms over the coefficient ring.}
For completeness this is a finite-dimensional calculation. Let
\(A\) be the total earlier degree. In a nonzero \(X_t\) a generic
\(\mathrm{GL}_{r-1}\) translate has a nonzero coefficient of the
\(t\)-th power of the first \(w\)-coordinate. Indeed a nonzero
homogeneous vector-valued polynomial is nonzero at some point,
and the simultaneous change of the earlier letters is invertible.
Its earlier letters have first weight at least \(-A\). The cutoff at \(C\) therefore gives
\[
 t-A\le C,\qquad t\le C+A,\qquad C-t\ge -A.
\]
All coordinate weights lie between \(-A\) and \(C<S\). In particular the coordinate weight differences
in the Hom construction are bounded by \(4S\). Highest vectors
for this auxiliary ordering are killed by every \(G_{*a}\).
In characteristic zero they are independent of \(w\). At degree
zero the early cutoff \(C\) is already the full cutoff, so all
integral highest vectors lift identically.

To see the index conclusion, first apply the type projectors to
an integral element of the direct sum, multiplying by \(p^{A_1}\)
to make each component integral. For a component of type \(U\)
use in \(U\) the lattice generated by divided powers from an
integral highest vector. This is a lattice: it spans by highest
weight theory and lies in a fixed stable lattice after scaling
the highest vector. Let \(L_{\rm im}\) be the integral image of the lift
and
\[
 T_j=p^{A_1}P(u_j^*\otimes v).
\]
These integral
intertwiners send the highest generator into \(L_{\rm im}\).
Every \(u_j\) is generated from that highest vector by
divided powers, which also preserve \(L_{\rm im}\).
Hence \(T_j(u_j)\in L_{\rm im}\) for every \(j\), and
the isotypic projector formula reads
\[
       p^{A_1}v=(\dim U)\sum_j T_j(u_j)\in L_{\rm im}
\]
on the component of type \(U\). Sum over the components and retain the required
homogeneous last degree of a lift to obtain \eqref{eq:res-sat}.

\proofstep{Imposing the highest-weight equations.}
It remains to impose the weight and highest-vector conditions.
Full weight projection is integral. Within that weight the
highest vectors for \(I\) are precisely the copies of the
isotypic summand with highest weight equal to the specified
\(I\)-weight, intersected with the weight space. This projector
costs at most \(A_0\), commutes with restriction and preserves
every upstairs prefix condition. Indeed those conditions are
stable under the total action on \(I\) and the full torus.
Thus the saturation exponent in the primitive-vector test is
at most
\[
 c=2A_1+A_0.
\]
For its self-adjoint projector first
retain monomials of the required \(w_\rho\)-order and weight.
On the restricted polynomials impose the early prefix type
projectors including the cutoff at \(K\). They commute: a
projector on an earlier prefix commutes with its group action
and with all operators on later letters, hence with the action
on the longer prefix.

The prefix projectors also commute with the total \(I\)-type projector. That projector selects the required
highest type, giving zero if it violates the full cutoff.
Repeated bounds on the same prefix take their minimum. The
zero-letter and one-letter prefix tests are scalar on a fixed
multidegree. For \(K\ge1\) at most \(K+1\) projectors with
denominator \(A_0\) suffice. These self-adjoint projections have
the required rational image by \eqref{eq:res-sat}; putting back
\(w_\rho^e\) preserves self-adjointness. We obtain
\[
             v_*\le (K+3)A_0+4A_1
             \le (K+3)r^2(1+\log_p(8a+r))
                        \qquad(S\le2a,\ K\ge1).
\]
The last inequality uses \(K+7\le2(K+3)\). For \(K=0\) the
original single symmetric power is irreducible, so the cutoff
is scalar and every allowed restricted polynomial lifts.
The weight and highest tests alone suffice and obey the same
final bound. For \(e\) greater than the last degree the graded
piece is zero.
\end{proof}

\subsection{The estimate on a segment with several pairs}
\hypertarget{source-tied-segments}{}

A segment may contain several adjacent pairs with linked variables. The one-pair Gram calculation must then be iterated on the successive spaces of restricted polynomials.

Consider \(N\) ordered vector or covector variables. Write \(\mathbf n=(n_i)\) for the degrees of the covectors that vary, and let \(Y(\mathbf n)\) be the space defined by the prescribed upper bounds on the first entry of the highest weight at each designated prefix. A tie is an adjacent vector--covector pair with operator \(\Phi_{z_i}\) as in Lemma~\ref{lem:tied-gram-polynomial}. The first positions
are strictly increasing, hence the pairs are disjoint. The end
of the vector of each pair is among the specified prefixes.
The covector degrees \(n_i\) range independently over progressions
of lengths \(H_i\ge1\) and positive steps prime to \(p\).
Denote progression \(i\) by \(I_i\). Its upper endpoint and the parameter set are
\[
 n_i^+=\max I_i,\qquad \mathbf n\in\prod_i I_i.
\]

Write \(S_i^+\) for the total degree through the vector of pair
\(i\) at the upper corner of the box. Suppose \(S_i^+\le\min n_i\),
all degrees and \(z_i\) are at most \(2a\), \(r\ge2\), \(p\) is odd,
\(a\ge1\), and
\[
 2aN+r\le p^2(8a+r)^2.
\]
Put
\[
 \Lambda=1+\log_p(8a+r),\qquad
 \kappa\ge\tfrac32(N+2)r^2\Lambda,\qquad E\ge1 .
\]
On the coefficient lattice \(Y_{\mathbb Z_p}(\mathbf n)\) use the
Fock form divided by the product of the letter factorials, with the
product of the tie operators applied in the first argument.
Denote by \(b(\mathbf n,D)\) its rank minus the largest size of
two systems of integral vectors whose cross matrix is \(p^D I\).
\begin{proposition}[Mean rank loss for tied pairs]
\label{prop:mean-tied-rank}
For \(t\) pairs and \(D=\lfloor t(E+\kappa)+\kappa\rfloor\),
the estimate is
\begin{align}
 \mathbb E_{\mathbf n} b(\mathbf n,D)
 &\le\sum_{i=1}^t\left(\frac{2S_i^+}{H_i}
            +\frac{2\kappa+2S_i^+\Lambda/H_i}{E}\right)
          \mathbb E_{\mathbf n}
              \dim Y(\mathbf n_{i\leftarrow n_i^+}).\label{eq:ties}
\end{align}
The expectations give every point in the box equal weight.
In particular the dimension on the right is evaluated with just
\emph{one} degree replaced by its endpoint.

\end{proposition}

\begin{proof}
\proofstep{The test for a single pair of letters.}
Fix a pair \(i\), with \(b\) letters through its vector, and first
impose only the cutoffs through that vector. In the prefix
including the following covector \(\eta\), the highest space for
\(\mathrm{GL}_r\) of weight \(\lambda\) is the image of the \(\mathrm{GL}_{r-1}\) highest-weight space of weight
\(w=\lambda+n_i e_\rho\) under the map
\[
 f\longmapsto f_{n_i}=\eta_\rho^{\,n_i-S_i}T_\eta f ,
\]
where \(S_i\) is the total degree through the vector of pair \(i\) and
\(T_\eta\) is the substitution used in Lemma~\ref{lem:tied-gram-polynomial}.
Evaluation at \(\eta=e_\rho\) takes a highest-weight polynomial to the
weight-\(w\) space for \(\mathrm{GL}_{r-1}\).

Order the basis with the distinguished coordinate \(\rho\) last.
On the chart \(\eta_\rho\ne0\), upper triangular unipotent matrices
whose only off-diagonal entries are in column \(\rho\) remove the other covector
coordinates. Invariance and homogeneity uniquely reconstruct the
polynomial there from this restriction. Acting on the old
letters gives the substitution in the preceding section:
the last coordinate of each vector becomes
\((v\cdot\eta)/\eta_\rho\), and the other components of a covector become
\[
 u'_j=\frac{\eta_\rho u_j-u_\rho\eta_j}{\eta_\rho}=u_j-u_\rho\frac{\eta_j}{\eta_\rho}\qquad(j\ne\rho).
\]

Multiplication by \(\eta_\rho^{n_i}\) clears all denominators
since \(n_i\ge S_i\), the prefix degree through the vector.
This expression is invariant under the upper unipotents (the
subgroup on the other coordinates acts on the original highest
vector) and has weight \(\lambda\). The earlier cutoffs persist:
in each prefix the simultaneous substitution is an invertible
group change over the chart followed by a scalar on every fixed
multidegree. Evaluation and reconstruction preserve integral
coefficients. Evaluation is inverse to the polynomial lift:
\[
 \left.f_{n_i}\right|_{\eta=e_\rho}=f .
\]
Thus \(f\mapsto f_{n_i}\) takes an integral basis of the
\(\mathrm{GL}_{r-1}\) highest-weight space to an integral basis of the
\(\mathrm{GL}_r\) highest-weight space.

In the prefix multiplicity space for each \(w\), take left and
right Smith bases of \(p^{\lceil\kappa\rceil}\Gamma_{n_i}\).
Let \(h_j,k_j\) be the images of these basis vectors under
\(f\mapsto f_{n_i}\), and let \(d_j\) denote
the corresponding diagonal entries. A good index satisfies
\[
 v(d_j)\le E.
\]

With the same scalar \(c_i\) as in the tied Gram matrix
(including \(p^{\lceil\kappa\rceil}\)), we have
\[
 c_i \langle\Phi_{z_i}h_j,k_l\rangle=d_j\delta_{jl},
 \qquad G_{jl}=c_i\langle h_j,h_l\rangle\in\mathbb Z_p.
\]
For the latter assertion use the factorial count in the term
with zero contractions: the numerator and denominator factorial
arguments have equal total and at most \((b+1)r\) numerator
terms are needed. The binomial factor \(\binom Bo\) is integral.
For the first matrix the binomial-basis estimate above costs at
most \(\lceil\kappa\rceil\), as \(b\le N-1\) and \(3S_i\le6a\).

\proofstep{The retained spaces and their extension operators.}
Let \(\pi_i\) retain the good \(h_j\)'s and kill the rest.
On these highest spaces define a second operator by
\[
            T_i h_l=\sum_{j\ {\rm good}}
                   p^{\lfloor E\rfloor}d_j^{-1}G_{jl}\,k_j .
\]
All its coefficients are integral, since \(v(d_j)\) is an integer.
Both operators extend equivariantly to each prefix isotypic
summand, by acting on multiplicities, and act as the identity
on the tensor factor of later letters. Here this last convention
means tensoring the prefix operator with the identity.

The
operators map the space with cutoffs through the vector to itself.
They preserve subsequent cutoffs by equivariance under the
larger prefix groups. They commute with all later tie operators
on the later letters. The Gram identities on highest vectors
give on this earlier constrained space, and after adjoining tails,
\[
     B(\Phi_{z_i}u,T_i y)
                   =p^{\lfloor E\rfloor}B(\pi_i u,y),
\]
where \(B\) is the normalized untied Fock form. Indeed pairings
between distinct types vanish, and on equivalent irreducibles
the form with the matrix-unit transposition as adjoint is
determined by its value on the two highest vectors. The same
normalizing scalar occurs on both sides.

Here are bounds for extending the integral highest-vector
calculation to the coefficients of arbitrary vectors. Put
\[
 e_1=\binom r2\lfloor\log_p(2aN+r-1)\rfloor.
\]

On every prefix isotypic projection costs at most \(e_1\) by the
Hom construction, since the width of coordinate weights is at
most the total degree. In an isotypic integral lattice,
multiplication by \(p^{e_1}\) expresses any integral vector
in the divided-power span of its integral highest vectors.
Use the integral intertwiners evaluated on the lattice generated
by a highest vector, as in the proof of \eqref{eq:res-sat}.
Apply these facts in the constrained prefix lattice, which is
divided-power stable because the diagonal prefix action preserves
all earlier type bounds. For an integral word each coefficient
of a later monomial is integral and obeys those earlier bounds.
Consequently \(p^{2e_1}T_i\) is integral on the full constrained
word lattice. We may use the weaker allowance
\[
 (N+2)e_1\le\kappa
\]

for every \(i\), by the size condition on \(2aN+r\).

\proofstep{Simultaneous restrictions for all pairs.}
Intersect \(Y\) with the fixed spaces of all \(\pi_i\), obtaining
a rational space \(W\). For \(u\in W\), apply the identity
\[
 B(\Phi_{z_i}v,T_i y)=p^{\lfloor E\rfloor}B(\pi_i v,y)
\]
successively for \(i=1,\ldots,t\), keeping the \(T_i\) in the same
composition order. At step \(i\), the projector \(\pi_i\) commutes with
the remaining \(\Phi_{z_j}\) for \(j>i\) and satisfies \(\pi_i u=u\).
Each later \(T_j\) preserves \(Y\).
The successive identities reduce the twisted pairing with the
composition of the \(T_i\) to \(p^{t\lfloor E\rfloor}B(u,y)\).
For a primitive \(u\in W\cap Y_{\mathbb Z_p}\), choose a monomial
\(m\) with unit coefficient in \(u\).

Let \(P_Y\) be the product of
the commuting prefix type projectors. Self-adjointness and
\(P_Yu=u\) give
\[
 y=p^{Ne_1}P_Ym\in Y_{\mathbb Z_p},\qquad
 B(u,y)=p^{Ne_1}B(u,m),\qquad v(B(u,y))\le Ne_1.
\]
The integer exponent clears at most one projector per prefix.
The valuation bound uses that the normalized Fock norm of \(m\)
is the reciprocal of an integer.
Thus each such \(u\) has a partner in the full integral lattice
for the twisted form with valuation at most \(D\).
In Smith coordinates the directions with valuation greater
than \(D\), including the null directions, span a space of
dimension \(b(\mathbf n,D)\). Write \(B_{>D}\) for this Smith subspace. The dimension formula gives
\[
 \dim(B_{>D}\cap W)\ge b(\mathbf n,D)-\operatorname{codim}_Y W.
\]
If the right side were positive, their intersection with \(W\) would contain a primitive
integral vector with no such partner. Hence
\(b(\mathbf n,D)\le\operatorname{codim}_Y W\).

\proofstep{Counting the lost continuation dimensions.}
To compute this codimension bound, let \(m_{iw}\) be the
dimension of the old highest space of weight \(w\) and let
\(d_{iw}(\mathbf n)\) be the dimension of the allowed continuations
of one irreducible summand generated by the image
\(f_{n_i}\) of its highest-weight vector \(f\). This
includes the copy itself tensored with the tail and intersected
with all remaining cutoff spaces. All copies of the same type
have identical continuation dimensions, since the remaining
prefix actions intertwine identifications of the copies tensored
with the identity on the tail. Complete reducibility and the
highest-space bijection give
\[
     \sum_w m_{iw}d_{iw}(\mathbf n)=\dim Y(\mathbf n).
\]
A bad Smith direction for pair \(i\) costs at most
\(d_{iw}(\mathbf n)\) dimensions. Let \(Y\) be the ambient continuation space and \(K_i\subseteq Y\)
the subspace satisfying the \(i\)-th condition. The diagonal map gives
\[
 Y/\bigcap_i K_i\hookrightarrow\bigoplus_i Y/K_i,\qquad
 \operatorname{codim}_Y\bigcap_iK_i\le\sum_i\operatorname{codim}_Y K_i .
\]
This applies without a commutativity assumption on the \(\pi_i\).

These continuation dimensions are nondecreasing in \(n_i\).
For an increase of \(u\) multiply by \(\eta_\rho^u\) and project
on the prefix through \(\eta\) to the Cartan component of highest
weight \(\lambda-u e_\rho\). Its generating highest vector and new weight are
\[
 \eta_\rho^u f_{n_i},\qquad \operatorname{wt}(\eta_\rho^u f_{n_i})=\lambda-u e_\rho.
\]
Indeed in the product of the old
copy with \(\operatorname{Sym}^u(\mathrm{Std}^*)\) this is the
top tensor type of multiplicity one. Its image under polynomial
multiplication is nonzero.

The prefix map commutes with all upper
raising operators, shifts weights by \(-u e_\rho\), and is
nonzero on the old highest line. Its kernel on that copy is
zero: any nonzero stable kernel for the torus and upper
operators would contain a highest vector. Different tail
monomials give independent images. Multiplication by a
covector power preserves the earlier cutoffs and does not
increase the maximal first weight for longer prefixes.
The projection preserves the longer bounds by equivariance
and its range satisfies the earlier ones because its highest-weight generator \(\eta_\rho^u f_{n_i}\) satisfies them. Thus multiplication followed by Cartan projection injects the old continuation space into the new one. In particular
\[
 d_{iw}(\mathbf n)\le d_{iw}(\mathbf n+u e_i)\qquad(u\ge0),
\]
for the two allowed degrees on this fibre.

Finally fix every coordinate except \(i\). Both the old highest
space and its integral basis may be fixed on this fiber.
For each \(w\) the polynomial estimate gives the mean number
of bad Smith directions at most \(m_{iw}\) times the parenthesis
in \eqref{eq:ties}. Indeed take \(h=\lceil\log_p
\max(H_i,2S_i+1)\rceil\le\Lambda\) in
\eqref{eq:smith-average}. The graded determinant bound gives
\(v_*\le(N+2)r^2\Lambda\), and
\[
 \lceil\kappa\rceil+v_*+3\Lambda\le2\kappa
\]
by the assumed
slack. Weight those bad counts by the endpoint continuation
dimensions and sum over \(w\). The dimension identity gives
the endpoint on the right of \eqref{eq:ties}. Average over the remaining coordinates and sum in \(i\). The resulting bound, also valid for an empty set of ties, is the one applied to the enlarged box of sampled degrees.
\end{proof}

\subsection{The segment lattice and the last sampled degree}
\label{sec:segment-window}
\hypertarget{source-segment-window}{}

For the segment application, the variables in \eqref{eq:ties} range over an enlarged integer box. The final sampled degree contributes a separate rank loss.
\begin{paramdefinition}[The sampled degrees and size conditions for a segment]\label{hyp:segment-sampling}
Write \(s\) for the number of odd tops
between successive even tops. Put \(r=k-1\), \(P_i=p^{G-d_i}\)
where \(G\) is the length of this segment and \(d_i\) are the
increasing distances of its odd tops from its start. Its \(s+1\)
sampled degrees are \(a_i\) at those tops and \(a_+\) at the next
even top. They range over the same enlarged integer interval
\[
 W=[a,\lfloor a+a/(kY)+x_a\rfloor],
\]
where
\(0\le x_a\le2\mathcal L b'/k\) and
\(\mathcal L=C_0N_1^2\log a\). The \(b_i,b_+\) are held fixed
in the enlarged second interval.

In particular \(b_i,b_+\le5b'/2\).
Write \(\lambda=1+\log_p(8a+r)\), reserve \(N\) for the number of
tops in the cycle, and put \(Q=p^{G'}\), \(G\le G'\).
The local size hypotheses include
\[
 k\ge5,\; 2s+2\le r,\; s\le N_1,\;
 p\ge k^2Y^2,\quad pb'<a,\quad kQ\le b',\quad
 4\mathcal L<k,\quad kY^2\log a\le a/b'.
\]
The parameter \(a\) in the integer construction is the rounded
lower endpoint. We use the inequalities on that integer.
\end{paramdefinition}
The host realization uses the good cycle of Section~\ref{subsec:frobenius-selection}. The estimates preceding Proposition~\ref{prop:host-large-prime-test} verify these inequalities with the rounded endpoint and the enlarged box. In particular the two growing scales there satisfy
\[
        a/b'=X',\qquad \log Q=o(\log b').
\]


\begin{lemma}[One lattice prescription on a digit box]\label{lem:segment-common-digit-lattice}
Fix the second degrees and a nonempty digit class. On the
product of admissible first-degree progressions both the
retained label conditions and the allowed levels
\(\mathbf j\) are independent of the sampled first degrees.
The last shifted vector order \(o_+\) is less than \(a\).
On the lattice at each permitted level one has
\[
 B_0(u,v)=B_0(v,u)\qquad(u,v\in L_{\mathbf j}).
\]

\end{lemma}

\begin{proof}
\proofstep{The integral digit class.}
A digit class consists of vectors \(\rho_i,\tau_i\) of length \(k\)
with entries from \(0\) to \(P_i-1\). Write \(|\rho|=\sum_c\rho_c\).
Retain the class at these degrees when both coarse degrees are nonnegative integers, where
\[
 A_i=(b_i-|\rho_i|)/P_i,\qquad C_i=(a_i-|\tau_i|)/P_i .
\]
Put \(A_0^{\,*}=\lfloor(a-E_0)/p^G\rfloor\), using a star to
distinguish this starting capacity from the individual \(A_i\).
The odd tops split at an index \(\sigma\) into an
initial list and a final list of \(t=s-\sigma\) tied pairs
(\(0\le\sigma\le s\)).

For a block label use integers
\((w_i,h_i)\) at the initial positions, \(o_i\) at tied positions,
and \(o_+\) at the end. Set \(u_i=w_i+h_i\) or \(o_i\), respectively,
and \(K_l=A_0^{\,*}+\sum_{i<l}(A_i-u_i)\), numbering positions from zero.
The retained labels are the nonnegative choices satisfying
\begin{align*}
 w_i&\le A_i,\quad h_i\le C_i,\quad
 w_i+h_i\le A_i(1+4/p^2)+\mathcal L K_i/k,\quad
 h_i\le2\mathcal L A_0^{\,*}/k &&(i<\sigma),\\
 o_i&\le A_i,\quad o_i\le K_i+\mathcal L A_i/k &&(i\ge\sigma),\\
 o_+&\le K_s+(\mathcal L/k+4/p^2)\sum_i2a/P_i+\mathcal L b_+/k .
\end{align*}
For a nonempty class, the admissible \(a_i\)'s form a product of
arithmetic progressions of steps \(P_i\). On that entire product,
\[
 P_iC_i\ge a-kP_i\ge a-b',
 \qquad P_iA_i\le5b'/2,\qquad
 P_i\,2\mathcal L A_0^{\,*}/k\le a/2.
\]

Since \(pb'<a\) and \(p\ge k^2Y^2\), these imply both
\(C_i\ge A_i\) and \(C_i\ge 2\mathcal L A_0^{\,*}/k\).
Thus the restriction \(h_i\le C_i\) follows from its fixed
cutoff, and each minimum defining an internal tied level
is its fixed vector degree \(o_i\).
Also \(o_+\le a\). For this, \(A_0^{\,*}\le a/25\),
\(K_s\le a/25+5sb'/2\), \(sb'\le a/8\), and
\[
 \sum P_i^{-1}\le p/(p-1)\le25/24,
\]
by the strictly decreasing
powers. Inserting these and \(b_+\le5b'/2\), \(kb'\le a/4\),
in the displayed upper bound gives \(o_+<a\).

The chain now has \(2s+1\) alternating vector and covector letters,
starting and ending with a vector. Their degrees are
\((w_i,C_i-h_i)\) in the initial list, \((o_i,C_i)\) in the
tied list, and \(o_+\) at the end. Adjoin a covector \(\xi\) of
degree \(b_+\). The differences between the sampled and tied vector degrees give the decorations
\[
 z_i=A_i-o_i,\qquad z_+=a_+-o_+.
\]

The capacities at the prefixes ending in a tied vector and at
the end of the chain are \(K_i,K_s\). Subtract the prescribed
common margin \(m_0\) and take floors there to obtain the cutoffs
on first highest weights.

\proofstep{The level parameters and their order.}
A level consists of integers \(j_i\) for the ties, including the
last one \(j_+\), between zero and the minimum of the two
corresponding degrees. Subtract \(j_i\) from both degrees of that
pair. The index set of levels is fixed over the class since
\(C_i\ge A_i\ge o_i\) at the tied positions. Let \(V\) be the
constrained chain space at the reduced degrees and \(L_V\) its
coefficient lattice over \(\mathbb Z_p\). Before restoring the scalar
variables the lattice at a level is
\[
 L_{\mathbf j}=L_V\otimes_{\mathbb Z_p}
       \mathbb Z_p[\xi]_{b_+-j_+}.
\]
The factors use disjoint variables.

Write \(d_F\) for the product
of the factorials of all the reduced degrees, including the
degree of \(\xi\). Write \(T\) for the product of the \(t\) internal
operators \(\Phi_{z_i}\) on the reduced pairs.
The two forms required on \(L_{\mathbf j}\) are
\[
 B_0(f,g)=\langle Tf,g\rangle/d_F,\qquad
 B_z(f,g)=\langle\Phi_z Tf,g\rangle/d_F,\quad z=a_+-o_+ .
\]
The first form is symmetric since every \(U^lD^l\) is self-adjoint
and the operators on distinct ties commute.

\end{proof}

\begin{proposition}[Uniform mean on a segment]\label{prop:segment-discard}
Fix a segment \(I\) with the degree and size conditions of
Definition~\ref{hyp:segment-sampling}. Hold fixed the margin of
the lattice prescription in Lemma~\ref{lem:segment-common-digit-lattice}
and the \(s+1\) second weights. With \(s+1\) sampled first-weight
coordinates, the equal-weight average of the discarded
fractions over the enlarged interval box satisfies
\[
 \mathbb E_{\mathbf a,a_+}\kappa_I\le C_A(s+1)N_1/Y
\]
at the depth \eqref{eq:segment-mean-depth}. The same absolute constant
works for the depth allowance and the mean.
\end{proposition}

The proof uses the several-pair estimate separately at the
internal degrees and at the last degree. It then transports
the resulting ranks to the factorial lattice and changes back
from the auxiliary sampling law to equal weights on the box.

\begin{proof}
\proofstep{The internal degrees.}
Freeze the untied coordinates \(a_i\) and the digits and level.
At a tied position with radix \(P=P_i\) the new covector degree
\(n_i=(a_i-|\tau_i|)/P-j_i\) runs through an interval of
consecutive integers of length \(H_i\). Its law is uniform there.
One has \(H_iP\ge a/(2kY)\), since the interval in \(a_i\)
has length at least \(a/(kY)-1\), and
\(2(1+P)kY\le4b'Y\le a\). All points of this residue class above
\(a\) exceed \(|\tau_i|\).

Let \(S_i^+\) be the prefix degree
through its vector at the largest values of the other tied
coordinates. Then
\[
 P S_i^+\le 5b'/2+(2s+1)2a/p,\qquad
 P\min n_i\ge a-kP-5b'/2.
\]
The first bound uses \(P_i/P_l\le1/p\) for every earlier position
and the bound \(2a\) on each degree before division by its radix.
The vector at \(i\) has \(Po_i\le b_i\).

The second bound follows
from \(Pj_i\le b_i\) and \(|\tau_i|<kP\).
They imply \(S_i^+\le\min n_i\), since \(kP\le b'\),
\(b'<a/p\), \(2s+1\le k-2\) and \(p\ge k^2\).
They also give
\[
 S_i^+/H_i\le7/Y
\]
by multiplying the
first upper bound by \(2kY^2/a\) and using \(p\ge k^2Y^2\).
All the reduced degrees and internal decorations are at most \(2a\).
The last vector satisfies the stronger bound \(o_+\le a\) above.
The degree-size condition of \eqref{eq:ties} follows from
\(2s+2\le r\) and \(p\ge k^2\).

\proofstep{The allowance for the internal degrees.}
Use \(\kappa=3(s+2)r^2\lambda\), \(E=Y\kappa\), and write
\begin{equation}\label{eq:segment-mean-depth}
 d=t(Y+1)\kappa+\kappa,\qquad d'=\lfloor Yd\rfloor.
\end{equation}
In \eqref{eq:ties} the required lower bound on \(\kappa\) for
\(2s+1\) letters is \(\tfrac32(2s+3)r^2\lambda\).
The coefficient of each endpoint dimension is at most \(31/Y\)
because its three terms are bounded by \(14/Y,2/Y\) and
\(14\lambda/(Y^2\kappa)\).

Write \(b_V\) for the failed rank on
\(L_V\) at depth \(\lfloor d\rfloor\) and \(\widehat a_i\) for
replacement of \(a_i\) alone by its largest compatible value.
Averaging first over these progressions, then over the frozen
coordinates, gives
\begin{equation}\label{eq:seg-interior}
 \sum_{\mathbf a} b_V(\mathbf a)
   \le \frac{31}{Y}\sum_{i\ {\rm tied}}\sum_{\mathbf a}
                     \dim V(\widehat a_i).
\end{equation}
This includes \(t=0\), where the unmodified form already has
sufficient depth. Tensoring each pair of successful vectors with
each monomial \(\xi^\nu\) preserves the depth for \(B_0\).
On the partner put the integral multiplier
\((b_+-j_+)!/\prod\nu_c!\). Thus the failed rank for \(B_0\)
is at most \(b_V d_\xi\), with
\(d_\xi=\binom{b_+-j_++r-1}{r-1}\).

\proofstep{The last parameter.}
Choose an orthogonal basis of \(L_{\mathbf j}\) for \(B_0\) over
\(\mathbb Z_p\). After clearing a common denominator, a pairing
of least valuation can be realized on a single vector. If all
diagonals have greater valuation and \(B_0(u,v)\) is minimal,
use
\[
 B_0(u+v,u+v)=B_0(u,u)+2B_0(u,v)+B_0(v,v).
\]
Since \(2\) is a unit, the middle summand has strictly smallest valuation. The resulting
primitive vector splits off by subtracting integral multiples.
Induction, with a basis of the radical appended, gives the
assertion. Take the part \(f_1,\ldots,f_h\) of this basis with
nonzero diagonal of valuation at most \(\lfloor d\rfloor\).
Smith reduction identifies \(h\) with the successful rank.

Let \(o=o_+-j_+\), \(n=b_+-j_+\), \(J=\min(o,n)\),
\(q=o+n+r-1\). For the last pair write \(U=x\cdot\xi\) acting
by multiplication, \(D=\sum_c\partial_{x_c}\partial_{\xi_c}\).
Set
\[
 M_l(u,v)=
 \frac{\langle U^lD^l Tu,v\rangle}{(l!)^2\binom Jl\,d_F}\quad(0\le l\le J).
\]
The matrix with entries
\[
 p^{x_0}\sum_{l=0}^J(-1)^l
       \binom{q+z-l}{J-l} M_l(f_i,f_j)
   =p^{x_0}\binom{q+z}{J}B_z(f_i,f_j)
\]
is a polynomial in \(z\) with integral binomial-basis
coefficients for the following choice. Put \(a_*=(3s+2)a\),
\(L_*=\lfloor\log_p(8a_*+r)\rfloor+1\) and
\[
 x_0=\big((t+1)(r+1)+(2s+2)r+2(t+1)\big)L_* .
\]
For the coefficient bound, expand \(D^l\) by multi-indices
\(|\alpha|=l\). After the two divided derivatives in every
coordinate, division by \((l!)^2\) costs the inverse multinomial
\(\binom l\alpha^{-1}\). Its valuation loss is at most
\((r-1)L_*\) by the factorial formula.

The additional reciprocal
binomial costs at most \(L_*\). The same reasoning applies to
every operator in \(T\) with its own denominator
\[
 \binom{o_i'+n_i'+r+z_i-1}{l},
\]
where primes indicate
reduced degrees. Finally the monomial norm divided by \(d_F\)
costs at most one multinomial per letter, altogether at most
\(((2s+2)r+2(t+1))L_*\). Multiplication by \(U^l\) has integral
coefficients. Vandermonde's identity for each binomial in
\(q+z-l\) now proves integrality. The coefficient of \(\binom zJ\)
is exactly \(p^{x_0}B_0(f_i,f_j)\). Its diagonal entries are nonzero. Thus the coefficient of the highest power of \(z\) in the determinant is
\[
 (J!)^{-h}\prod_{i=1}^h p^{x_0} B_0(f_i,f_i).
\]
Consequently the determinant has degree \(hJ\), and the
leading coefficient has valuation at most
\(h(x_0+\lfloor d\rfloor-v_p(J!))\).

\proofstep{A Smith estimate over the last interval.}
Apply \eqref{eq:smith-average} with \(H=|W|\), shift \(a-o_+\),
\(e=\lceil\log_p(a+1)\rceil\) in place of \(h\) there, and depth
\(d'+1\). Both \(J+1,H\le a+1\).
The integral nonzero factor \(p^{x_0}\binom{q+z}{J}\) only
increases Smith valuations. Integer valuations mean that the
failed directions of the rescaled matrix at depth \(d'\) are
precisely those counted at \(d'+1\).
Legendre's formula gives
\[
 J/(p-1)-v_p(J!)\le L_*.
\]

Since \(3s+2\le p\), \(L_*\le2\lambda\), and for \(r\ge4,t\le s\)
direct multiplication gives \(x_0+L_*\le\kappa\).
Also \[
 \begin{gathered}
 J/H\le(5b'/2)kY/a\le1/(2Y), \\ e\le\kappa, \\ d\ge\kappa, \\ d'+1\ge Yd.
\end{gathered}
\]
The mean number lost on this orthogonal subsystem is therefore
at most
\[
 h\left(\frac{J}{H}+
   \frac{x_0+\lfloor d\rfloor+L_*+Je/H}{d'+1}\right)
 \le \frac{5h}{Y}.
\]
Combining with the directions already missing in \(B_0\)
bounds the mean failed rank of \(B_z\) on the full level by
\(d_\xi b_V+5d_\xi\dim V/Y\). This last average uses \(a_+\)
with \(\mathbf a,b_i,b_+\) still fixed.

\proofstep{From the levels to the factorial lattice.}
Introduce one extra pair of scalar variables \(Z_i,S_i\) at each
tie, including the last pair. In the following products \(i\)
runs over all \(t+1\) pairs. Write \(E_i=Z_iS_i+U_i\).
The rational block space is spanned by
\[
 \left(\prod_i E_i^{j_i}Z_i^{z_i}\right) f,\qquad
               f\in L_{\mathbf j},\quad \mathbf j\text{ admissible}.
\]
Extension of scalars to the coefficient field \(K=\operatorname{Frac}O\)
gives precisely the span in the definition of the segment lattice.
In fact a highest-weight cutoff is the largest prefix-stable
subspace of a monomial support condition. In fixed multidegree it is
the common kernel of maps obtained by applying words in the matrix
units and projecting to the forbidden weight spaces. There are only
finitely many independent equations, all with rational coefficients.
Kernels, their finite intersections and their dimensions commute
with field extension. The integral coefficient intersections span
these spaces over the fraction fields by clearing denominators.

\proofstep{The normalized form on one pair of letters.}
For one pair with reduced degrees \(o,n\) and decoration \(z\),
let \(H\) be harmonic projection in \(r+1\) variables at degrees
\(o+z,n\). Write \(F=D+\partial_Z\partial_S\).
For \(f\) independent of \(Z,S\),
\[
 H(Z^z f)=\sum_{l=0}^{\min(o,n)}
  \frac{(-1)^l E^l Z^z D^l f}{(l!)^2\binom{o+n+r+z-1}{l}}.
\]
Indeed \[
 F(E^l v)=E^lFv+l(r+1+\deg_x v+\deg_y v+l-1) E^{l-1}v\,;
\] the coefficients telescope and the map is the
identity modulo \(E\). The Fock adjoint of \(E\) is \(F\), so
\(H\) is self-adjoint. Replace \(Z_i^{z_i}\) in the level map
by \(f\mapsto H_i(Z_i^{z_i} f)\), taking products of these
commuting maps, followed by \(\prod E_i^{j_i}\).
This gives a map \(\Theta_{\mathbf j}\). The terms other than the
original generator have larger total level. They belong to the
indicated higher-level spaces. For a prefix including both letters
the contraction \(D_i\) is equivariant. Before the pair it acts on disjoint variables.

In either case, for the corresponding prefix action \(g\),
\[
 D_i(gf)=g(D_i f).
\]
At the intermediate prefix
it differentiates the vector. The span of the derivatives there
is an image of the old representation tensored with the dual
standard representation, so each highest weight is obtained by
subtracting one box by the dual Pieri rule. This preserves the
upper cutoff. At the last pair the covector derivative acts only
on \(\xi\). Thus the replacement is triangular, with identity
diagonal, among the rational generators. The original level spaces
are independent because in \(E_i^{j_i}\) the term of greatest \(S_i\)-degree
is \((Z_i S_i)^{j_i}\). Extract successive lexicographically
largest scalar multidegrees in a relation. Their coefficients are
exactly the separate polynomials \(f\). Summing the dimensions of these independent level spaces gives
\[
 \dim(\text{block})=\sum_{\mathbf j}d_\xi\dim V.
\]


\proofstep{Orthogonality of distinct levels.}
Images of distinct levels under \(\Theta\) are orthogonal. In a
pair with unequal levels move \(E\)'s across the pairing as
\(F\)'s until a positive power of \(F\) hits a harmonic vector.
For equal levels, moving all \(E\)'s across gives
\(j!(r+1+o+n+z)^{\overline{j}}\) at this pair. Self-adjointness
reduces the pairing of the two harmonic projections to that of
\(Z^zf\) with \(H(Z^zg)\). Only the term of \(S\)-degree zero
survives, yielding \(z!\langle f,\Phi_z g\rangle\).
Multiplying the identities gives the raw Fock pairing of the images
as \(\ell_{\mathbf j}d_F B_z\), where
\[
 \ell_{\mathbf j}=\prod_i (j_i!)^2
       \binom{o_i'+n_i'+r+z_i+j_i}{j_i} z_i!.
\]
Here \(o_i',n_i'\) are the reduced degrees of that pair.
The same operator bound used for \(M_l\) shows that
\(p^U\Theta_{\mathbf j}f\) has integral coefficients for every
\(f\in L_{\mathbf j}\), with
\[
 U=(t+1)\lceil (r+1)(1+\log_p(8a_*+r))\rceil.
\]

For this bound expand \(D_i^l\) before multiplying by the
extended \(E_i^l\), which is integral. All the relevant
factorial arguments are at most \(8a_*+r\).
This argument only needs membership in the rational block span
and the asserted coefficient bound.

\proofstep{Uniform denominator and depth bounds.}
Let \(M=\max_{\mathbf j}v_p(\ell_{\mathbf j}d_F)\).
For each of the good systems in \(B_z\) at depth \(d'\), apply
\(p^U\Theta_{\mathbf j}\) on both sides and multiply the partners
by \(p^M/(\ell_{\mathbf j}d_F)\). The scalar in the raw pairing is
\[
 p^{2U}\frac{p^M}{\ell_{\mathbf j}d_F}\,\ell_{\mathbf j}d_F\,p^{d'}=p^e,\qquad e=d'+2U+M.
\]
Orthogonality between distinct levels makes the combined matrix \(p^eI\). The factorial metric on the actual
block absorbs almost all of \(M\). More explicitly, with
\[
 v=v_p,\quad v(w)=1/2,\quad
 b_*=\prod_{i<\sigma}(A_i-w_i)!\,h_i!,
\]
that lattice consists of polynomials in the rational block span
whose monomials \(c_\nu X^\nu\) satisfy
\[
 2v(c_\nu)+v_p(\nu!)+v_p(b_*)\ge0;
\]

\(\nu!=\prod_j\nu_j!\). Its pairing is \(b_*\) times raw Fock.
For any supported \(\nu\) in this block and \emph{any} level,
the total of the factorial arguments in
\(\prod_i(j_i!)^2 z_i!\, d_F\) is exactly \(|\nu|\).
Each restored power of \(E_i\) contributes two degrees and each \(Z_i\) contributes one. Thus
\[
 |\nu|=\sum_i (2j_i+z_i)+\text{total reduced degree},
\]
the sum of the factorial arguments just listed. Legendre's formula shows at
each power of \(p\) that the difference of the sums of the
floors in these factorials and in \(\nu!\) is at most the number
of variables in \(\nu\). Each of the \(t+1\) binomials in
\(\ell_{\mathbf j}\) adds at most one at that power.

Each single
coordinate exponent is at most \(8a_*+r\), by the pairwise
bidegrees and disjoint support. Thus, writing
\(L_0=\lfloor\log_p(8a_*+r)\rfloor\),
\[
 M-v_p(\nu!)\le [(2s+2)r+3(t+1)]L_0.
\]
Since \(3s+2\le p\), \(L_0\le\lambda\) and
\(U\le(t+1)((r+1)(\lambda+1)+1)\).
It follows that
\[
 2U+M-v_p(\nu!)\le
 L_{\rm pass}=10(s+1)(r+1)\lambda+\kappa_0,\qquad
 \kappa_0=2Nk\log_p a .
\]
For \(F_0=\lfloor d'+L_{\rm pass}\rfloor\), multiply both systems
by \(w^{F_0-e-v_p(b_*)}\). The exponent can be negative.
The monomial inequality just proved, and the integrality before
this multiplication, put every resulting vector in the factorial
lattice.

Since \(v(w)=1/2\), the new cross matrix is scalar with valuation
\[
 v_p(b_*)+e+2v(w)\bigl(F_0-e-v_p(b_*)\bigr)=F_0.
\]
Smith reduction in the full lattice now gives at least
as many good directions as this scalar matrix has rows. For
example, apply Smith reduction to the full form after clearing
denominators. A restriction by integral columns and rows cannot
increase the number of factors dividing a fixed scalar. Modulo
one more uniformizer power than that scalar, the image of the
restriction factors through a module generated by the low Smith
columns. Nakayama's lemma gives the assertion. One can then take
the saturated span of the corresponding left Smith vectors in the
full lattice. Thus passage to saturation preserves the partner
criterion even if the initially constructed systems span a
nonsaturated module.

The depth needed in the final segment test is at most
\[
 Yd+L_{\rm pass}\le C_A Y^2 N_1^2 k^2\lambda+2\kappa_0=:D_I .
\]
For the inequality use \(t\le s\le N_1\), \(N_1,Y\ge1\),
\(s+2\le3N_1\), and \(d=(tY+t+1)\kappa\). A constant
\(C_A\ge56\) suffices for this depth estimate.

\proofstep{Returning to equal weights on the box.}

Write \(D_k(u)=\binom{u+k-1}{k-1}\) for \(u\ge0\), zero
otherwise. For a compatible digit class \(\delta\), put
\[
 m_\delta(\mathbf a)=D_k(A_0^{\,*})D_k(b_+)
                  \prod_i D_k(A_i)D_k(C_i).
\]
The ordinary retained-rank comparison needed here is
\begin{equation}\label{eq:ordinary-class}
 \tfrac12 m_\delta(\mathbf a)\ \le\
   \sum_{\beta\ {\rm retained\ in}\ \delta}
       \sum_{\mathbf j} d_\xi\dim V_{\beta,\mathbf j}(\mathbf a)
   \ \le\ m_\delta(\mathbf a).
\end{equation}
This comparison concerns dimensions in characteristic zero.
Its role and its required margin are recorded below.
Let \(g_\delta(\mathbf a,a_+)\) be the sum over these levels
of their failed ranks in \(B_z\) at depth \(d'\).
The preceding construction supplies a bound by \(g_\delta\) on
the failed rank of the actual block sum.

For each class the
level set is fixed over its box of compatible degrees.
Summing \eqref{eq:seg-interior} and the last-parameter estimate
therefore gives
\[
 \sum_{\mathbf a,a_+}g_\delta
 \le \frac{|W|}{Y}
      \sum_{\mathbf a}\left[31\sum_{i\ {\rm tied}}
         m_\delta(\widehat a_i)+5m_\delta(\mathbf a)\right].
\]
Only one \(D_k(C_i)\) changes at the endpoint. If it changes from
\(D_k(u)\) to \(D_k(v)\), then \(P_i u\ge a-b'\ge3a/4\),
\(P_i(v-u)\le a/(kY)+x_a\), and
\(r x_a\le 2\mathcal L b'\le a/8\).
The binomial product gives \(D_k(v)/D_k(u)\le
\exp(r(v-u)/u)\le e^3\).

\proofstep{Returning to ordinary rank weights.}
Sum over all classes and set \(m(\mathbf a)=\sum_\delta
m_\delta(\mathbf a)\). For each letter every degree-\(a_i\)
monomial has a unique expression for its exponent vector as
\(\tau_i+P_i\gamma\). Hence
\(\sum_{\tau_i}D_k(C_i)=D_k(a_i)\), and likewise for \(A_i,b_i\).
It follows that
\[
 m(\mathbf a)=D_k(A_0^{\,*})D_k(b_+)
                  \prod_i D_k(b_i)D_k(a_i).
\]
If \(C_A(t+1)N_1/Y\ge1\) the segment bound is automatic.
Otherwise \(s\le N_1\le Y\). Throughout the box
\(\max m/\min m\le\exp(sr(1/(kY)+x_a/a))\le e^3\)
when \(m>0\). Indeed \(s\le Y\),
\(rx_a\le2\mathcal L b'\), and
\(kY^2\log a\le a/b'\), \(4\mathcal L<k\) give
\(sr x_a/a\le1\) for \(a\ge e\).

The size assumptions imply
\(a\ge p b'\ge25\).
If \(D_k(A_0^{\,*})=0\), \eqref{eq:ordinary-class} forces all
ranks to vanish and the bad share is zero by definition.
Otherwise divide the total failed rank at each point by its
actual retained rank, which is at least \(m(\mathbf a)/2\).
The result for the \emph{unweighted} box average is
\begin{equation}\label{eq:operational-seg}
 \mathbb E_{\mathbf a,a_+}\kappa_I
    \le \frac{2 e^3(31 e^3 t+5)}{Y}
    \le \frac{C_A(s+1)N_1}{Y}
\end{equation}
for one absolute \(C_A\) in both the depth and the mean.
The margin and second weights were held fixed throughout. For
each segment its odd positions and its next even position are
distinct, hence each assignment of these \(s+1\) weights has the
same number of extensions to the full enlarged box. This proves
the required sampling law on the cycle. Section~\ref{sec:retained-rank} proves the ordinary-rank comparison
used here.

  The factorial-lattice calculation and the digit-class sum give the closed bound for the Gram estimate.
For the mean bound, use \(m_0=C_m kN_1Qc+12c\) and \(D_I\) as above. In its interval
\(a=\lceil a_{\rm real}\rceil\) and
\[
 x_a=\lfloor a_{\rm real}+a_{\rm real}/(kY)\rfloor+ \lfloor J_0\rfloor+c_{\rm cl}-(a+a/(kY)),
\]

where \(J_0=\mathcal L b'/k\) uses the integer lower endpoint.
Thus \(W\) is precisely the enlarged interval used in the
edge comparison. Its numerical hypotheses require
\(0\le x_a\le2J_0\), \(0\le m_0\le b'/(kQY)\).
\end{proof}

\subsection{The ordinary ranks of the retained digit classes}
\label{sec:retained-rank}
\hypertarget{source-retained-rank}{}

A successful specialization controls the Gram form on the retained space. Estimate its ordinary dimension before multiplying the segment contributions in \eqref{eq:ordinary-class}.
\begin{established}[Separation and degree reserves on the retained digit class]\label{hyp:retained-digit-degrees}
Work in one compatible digit class. Write \(B=A_0^{\,*}\),
\(M=\lceil m_0\rceil\), \(\delta=4/p^2\) and
\(\epsilon=N_1\mathcal L/k\). The separation of the odd powers
from the transition between the two lists is important here.
We give the numerical details used from the separated regime.
Let \(\Delta<1/6\) be its relative band width parameter, with
\(40\epsilon\le\Delta\).

At the first list
\(p^{d_i}<(1-2\Delta)a/b'\), at the second
\(p^{d_i}>(1+2\Delta)a/b'\). In the enlarged box
\(|b_i/b'-1|\le b_e\le\Delta/2-18\epsilon\),
\((E_0+Q)/a\le\epsilon\) and \(kQ\le\epsilon b'\).

\begin{equation}\label{eq:coarse-split}
 (1+\Delta)A_i<B\quad(i<\sigma),\qquad
 B<(1-\Delta)A_i\quad(i\ge\sigma).
\end{equation}
The prefix sums satisfy

\begin{equation}\label{eq:coarse-prefix}
 \sum_{l<i} A_l\le\delta A_i,\qquad
 \sum_{l<i} C_l\le\delta C_i .
\end{equation}
\end{established}
\begin{proof}
The inequalities for the box are verified on the selected host cycle in the size calculation of Proposition~\ref{prop:host-large-prime-test}. At each interior top the digit and floor identities give \[
 \begin{gathered}
 b_i-kP_i<P_i A_i\le b_i, \\ P_i p^{d_i}=p^G, \\ a-E_0-p^G<Bp^G\le a.
\end{gathered}
\] For the first inequality use
\((1-2\Delta)(1+b_e)(1+\Delta)\le1-\epsilon\).

For the second use
\((1+2\Delta)(1-b_e-\epsilon)(1-\Delta)\ge1\).
Both follow by expanding with the bounds on \(b_e,\Delta\).
For consecutive indices \(P_i/P_{i+1}\ge p^2\), as both
distances are odd. The bounds just used give
the two ratio bounds
\[
       p^2A_i\le2A_{i+1},\qquad p^2C_i\le2C_{i+1},
\]
using \(a_i\in[a,a+2a/(kY)]\) and the digit error \(kP_i\). Summing the resulting geometric progressions gives \eqref{eq:coarse-prefix}.

We may use \(B\ge a/(2Q)\), by the floor bound and
\((E_0+Q)/a\le\epsilon\). With \(4 k^3 QY\le b'\),
\(kY\mathcal L b'\le a\), \(\mathcal L\ge24\), and
\(m_0\le b'/(kQY)\), we have
\[
 kM/B\le4b'/a\le1/30,\qquad M\le\epsilon B.
\]
Also \(k^3\le A_i,C_i,B-M\) by the digit and floor bounds,
and \(k\ge10\) by \(40\epsilon\le\Delta\). These inequalities
apply with absolute size constants, simultaneously on the
enlarged box.
\end{proof}

\begin{proposition}[Ordinary rank of one retained digit class]\label{prop:retained-digit-rank}
Fix a compatible digit class satisfying
Lemma~\ref{hyp:retained-digit-degrees}, using the retained space
in \eqref{eq:ordinary-class}.
Let \(R_{\rm ret}\) be the ordinary rank of the retained space
and \(N_{\rm mon}\) the unrestricted monomial count of this class. Then
\[
 R_{\rm ret}\ge N_{\rm mon}
 (1-e^{-\mathcal L/8})(1-e^{-\mathcal L/3})(1-rM/B).
\]
For \(\mathcal L\ge24\) and \(kM/B\le1/30\), this fraction
is greater than one half. When \(s=0\) there is no loss
from the internal tests.
\end{proposition}

\begin{proof}
\proofstep{The two counts.}
For \(A\in\{B,B-M\}\), put \(V=\mathbb C^k\). The tensor and
its dimension are
\[
 \mathcal T_A=\operatorname{Sym}^{A}V\otimes
       \bigotimes_i(\operatorname{Sym}^{A_i}V\otimes
                         \operatorname{Sym}^{C_i}V^*),\qquad
 T(A)=D_k(A)\prod_i D_k(A_i)D_k(C_i).
\]
In characteristic zero, each Pieri step gives one copy of each
allowed highest weight. For decreasing integer rows write
\(\nu\succ\mu\) when \(\nu_j\ge\mu_j\ge\nu_{j+1}\) wherever
the entries exist. From the zero weight add the initial strip of
size \(A\). At each subsequent pair the allowable transitions are
\[
 \mu\ \longrightarrow\ \nu\ \longrightarrow\ \tau,\qquad
 \nu\succ\mu,\quad |\nu|-|\mu|=A_i,\qquad
 \nu\succ\tau,\quad |\nu|-|\tau|=C_i ,
\]
where \(|\mu|\) denotes the sum of the rows, including negative
entries. These rules follow from their polynomial versions by a
determinant twist and duality. Restriction to \(\mathrm{GL}_r\)
appends one interlacing weight \(\lambda\) below the last weight,
again with multiplicity one. \cite{fulton-harris}

Let \(\mathcal P_A\) denote these entire paths and
\(d_\lambda=\dim V_\lambda\). The terminal dimension and the
independent choice of the last-coordinate exponent \(j_+\)
give the two counts
\begin{equation}\label{eq:retained-path-dimensions}
 \sum_{\pi\in\mathcal P_A}d_{\lambda(\pi)}=T(A),\qquad
 \sum_{j_+=0}^{b_+}D_r(b_+-j_+)=D_k(b_+).
\end{equation}
The second identity partitions degree-\(b_+\) monomials
according to their exponent on the last variable.

\proofstep{The row recursion for a counted path.}
Number rows from zero. Let \(c_i\) be the first row before
the \(i\)-th vector addition, \(c_0=A\), and \(c_s\) the final
one. Let \(w_i\) be the increase in the sum of all other
rows at the addition and \(h_i\) the decrease in the first
row at the removal. Write
\[
 u_i=w_i+h_i,
\]
and let \(R_i\)
be row 1 just after the addition.

Write \(n_i=-\lambda_{r-s+i}\)
and let \(d\) be the sum of \(\lambda\) minus the sum of the
rows below the first row at the end. Deleting the first row
after the initial letter gives a path in dimension \(r\)
with alternating strip sizes \(w_i,C_i-h_i\) and one final
addition of size \(d\) ending at \(\lambda\). Indeed the last
interlacing is exactly such a horizontal addition. Record \((w_i,h_i)\) in the initial list. The tied positions and the endpoint have labels
\[
 (o_i,j_i)=(u_i,h_i),\qquad o_+=d+j_+.
\]
Conversely these data and
the reduced path uniquely recover the large path at a
specified \(A\). Its first rows at the three successive times
of a pair must be
\[
 c_i,\quad c_i+A_i-w_i,\quad c_{i+1}=c_i+A_i-u_i .
\]
Thus sizes determine the restored entries. Their interlacing
with the retained rows is the condition to check for the
converse.

\proofstep{Bounding the upper count by the unrestricted list.}
For the upper bound start with a counted block, level and
allowed path in dimension \(r\). The independent factors in \eqref{eq:retained-path-dimensions}
give the count on this side as the sum of \(d_\lambda D_r(b_+-j_+)\) because each prefix
cutoff takes the sum of the allowed prefix summands in the
successive Pieri decompositions. The triangular lift of levels
in Section~\ref{sec:segment-window} preserves this total rank.
Reconstruct at \(A=B\), where \(c_i=K_i\). The first row of the
reduced path just after its vector must be at most both
\(K_i,K_{i+1}\). At a tie its cutoff gives the former and
\(K_{i+1}\ge K_i\).

In the initial list this first row is at
most
\[
 \sum_{l\le i}w_l\le(1+\delta)A_i,
\]
since a covector
step never increases it. Up through this pair every
\(K_l\ge B-2s\mathcal L B/k\), by \(w_l\le A_l\) and the
bound on \(h_l\). By \eqref{eq:coarse-split} and
\[
 \delta+3\epsilon\le\Delta/2,
\]
both required inequalities
follow. The final cutoff gives \(\lambda_0\le K_s\). The strip
degrees give \(0\le h_i\le C_i\), \(0\le w_i\le A_i\)
also at the ties. These conditions suffice for every inserted
horizontal strip and the final restriction. We have a
weight-preserving injection into the unrestricted count,
proving the upper bound of \eqref{eq:ordinary-class}.

\proofstep{A lower bound by deletion.}
For the lower bound start with \(A=B-M\), so the capacities
reconstructed from the labels satisfy \(K_i=c_i+M\). Retain
only paths satisfying, at every \(i\),
\begin{equation}\label{eq:kept-tail}
 u_i-R_i+\delta(A_i+r)\le \frac{\mathcal L}{k}\max(c_i,A_i),
 \qquad n_i\ge (1-\mathcal L/k)C_i ,
\end{equation}
and \(j_+\le\mathcal L b_+/k\).
We verify that these paths contribute to the block count at
margin \(m_0\). Interlacing and \eqref{eq:coarse-prefix} give
\[
 0\le R_i\le c_i,\qquad
 R_i\le w_i+\delta A_i,\qquad
 0\le n_i\le(1+\delta)C_i .
\]

\proofstep{Keeping enough rows after deletion.}
For the second bound the previous row 1 was at most the sum of
previous vector additions below the first row. For the third,
follow the bottom row just after pair \(i\). Until then its
negative part was at most \(\sum_{l\le i} C_l\). A later vector
step increases each row, a covector step decreases a row at
most to the former next row, and \(\lambda_j\ge\mu_{j+1}\).
Following these comparisons back proves the upper bound. A horizontal strip creates at most one new row of the corresponding sign. After \(i\) pairs, starting with the initial vector, the row counts thus satisfy 
\[
 \#\{a:\lambda_a>0\}\le i+1,\qquad\#\{a:\lambda_a<0\}\le i,
\]
 for the intermediate highest weight \(\lambda\). The inequalities about signs follow
from the same interlacing, in particular row 1 is nonnegative
after a vector step and the last \(s\) rows of \(\lambda\) are
nonpositive (interlace downwards from the zero middle rows).

In the initial list
\[
 c_i\le B+\delta A_i\le2B.
\]

Thus \eqref{eq:kept-tail} gives
\(u_i\le A_i(1+\delta)+2\mathcal L B/k\) and
\[
 c_{i+1}\ge c_i-(\delta+2\mathcal L/k)B.
\]
Starting at
\((1-\epsilon)B\) or higher gives
\[
 c_i\ge(1-4\epsilon)B\ge A_i
\]
throughout that list.
Reinserting this comparison in \eqref{eq:kept-tail} gives
\[
 u_i\le A_i(1+\delta)+\mathcal L c_i/k,\qquad h_i=u_i-w_i\le\mathcal L c_i/k\le2\mathcal L B/k.
\]
In the tied list one has
\(c_i\le B+\delta A_i<(1-\Delta+\delta)A_i\).

Consequently \(u_i\le c_i+\mathcal L A_i/k\le A_i\).
The final degree satisfies
\[
 d\le c_s+\sum_i(C_i-n_i)\le c_s+(\mathcal L/k)\sum_i C_i.
\]
In fact, on the final upper rows the increments of the
restriction telescope by interlacing to at most \(c_s\)
through the last middle row. On the last \(s\) rows the
negative mass before restriction is at most \(\sum C_i\)
and afterward equals \(\sum n_i\). This proves the first
inequality even when some last rows vanish. Since
\(C_i\le 2a/P_i\), the cutoff for \(o_+\) holds.
The first rows of the reduced path at a tie and at the end
are at most \(c_i\) and \(c_s\) by interlacing, exactly the
cutoffs \(\lfloor K_i-m_0\rfloor\) and
\(\lfloor K_s-m_0\rfloor\). Finally
\[
 j_i=h_i\le\min(o_i,C_i)
\]

and \(j_+\le\min(o_+,b_+)\), giving legal levels. The map is
injective by the reconstruction above and preserves the
weight of each path.

\proofstep{Weight of discarded paths.}
Write \(\eta=1/k^2\) and \(\tau=\mathcal L/(4k)\). Use the following coordinates for one step,
\[
 \begin{aligned} x&=c_i+r, & y&=(A_i+r)(1+\delta),\\ v&=R_i+r, & z&=u_i-R_i+\delta(A_i+r).\end{aligned}
\]
Then
\(0\le v\le\min(x,y)\), \(z\ge0\), and
\[
 v(c_{i+1}+r)=v(x+y-v-z)
             =xy-(x-v)(y-v)-vz\le xy .
\]
If the first inequality of \eqref{eq:kept-tail} fails then
\(z>2\tau\max(x,y)\), since \(A_i\ge k^3\) and
\(\max(x,y)\le2\max(c_i,A_i)\).
For \(v\ge\min(x,y)/2\) the term \(vz\) exceeds \(\tau xy\).
For smaller \(v\) the other term is at least \(xy/4\ge\tau xy\).
Thus the step has the extra factor \(1-\tau\) on its upper
bound.

Also \(y\le e^{2\eta} A_i\) and
\(n_i+r\le e^{2\eta} C_i\) since \(\delta\le\eta\) and
\(r/C_i,r/A_i\le\eta\). A failure of the second inequality
gives
\[
 n_i+r\le (1-\tau)e^{2\eta} C_i,
\]
by
\(n_i<(1-4\tau)C_i\), \(r/C_i\le\eta\) and \(\eta\le\tau\).
Multiplying and cancelling the successive \(c_i+r\) gives
for any failed path, since \(\lambda_0\le c_s\) and
\(A+r\le e^\eta A\),
\[
 \frac{(\lambda_0+r)\prod_i(R_i+r)(n_i+r)}{A\prod_i A_iC_i}
 \le e^{(4s+1)\eta}(1-\tau).
\]
The \((r-1)\)-st power of the right side is at most
\(e^{2-\mathcal L/5}\). Indeed \(2s+1\le r\) and
\((r-1)/(4k)\ge1/5\).

\proofstep{Converting the path loss into a dimension loss.}
To turn this into a dimension bound, row \(s-i\) of
\(\lambda\) is nonnegative and at most \(R_i\).
Each remaining vector interlaces it with the next row,
each covector can only decrease a fixed row, and the
final restricted row is at most the corresponding full
row. Its lower bound follows by interlacing with the zero
middle rows, or with the nonnegative entries above them.
Rows with \(s<j<r-s\) vanish.

With \(V\) the standard representation, the desired occurrence is
\[
 V_\lambda\lhook\joinrel\longrightarrow \bigotimes_{\lambda_i>0}\operatorname{Sym}^{\lambda_i} V\,\otimes\!\bigotimes_{\lambda_i<0}\operatorname{Sym}^{-\lambda_i} V^*.
\]
Fill positive rows singly from the top, then negative
rows singly from the bottom, using Pieri at each step.
Consequently
\[
 d_\lambda\le D_r(\lambda_0)\prod_iD_r(R_i)D_r(n_i).
\]
Use \(z^{r-1}\le(r-1)!D_r(z)\le(z+r)^{r-1}\) and
\(D_r(z)\le D_k(z)\) on the unrestricted tensor. A single
failed path therefore has weight at most
\(e^{2-\mathcal L/5}T(A)\).

There are at most \((8a+1)^{(2s+3)(2s+1)}\) possible paths
for \(s\ge1\). At each of at most \(2s+3\) weights only the
top \(s+1\) and bottom \(s\) entries need be recorded.
Each has absolute value at most \(4a\), since at a vector or
covector step an individual absolute change is at most
the degree, and
\[
 A+\sum_i(A_i+C_i)\le4a
\]
by the
geometric progression of the \(P_i\), \(P_i\ge p\),
\(A\le a\) and \(A_i,C_i\le 2a/P_i\).

The same bound holds upon interlacing. Thus the log of the
number is at most \(30N_1^2\log a\), using \(s\le N_1\)
and \(a\ge25\). Increasing the absolute constant \(C_0\)
if necessary, this and the additive 2 in the exponent
are at most \(3\mathcal L/40\). The total failed weight
is at most \(e^{-\mathcal L/8}T(A)\). For \(s=0\)
there are no failures of \eqref{eq:kept-tail}.

\proofstep{The loss at the last label.}
For the last label put
\[
 J=\lfloor\mathcal L b_+/k\rfloor.
\]

The discarded monomials have last-coordinate exponent
greater than \(J\). If \(J<b_+\), subtract \(J+1\) from the last coordinate
of each discarded monomial. This bijects them with monomials
of degree \(b_+-J-1\), so their proportion is
\[
 \frac{D_k(b_+-J-1)}{D_k(b_+)}
 =\prod_{i=1}^r\left(1-\frac{J+1}{b_++i}\right)
 \le \exp[-r(J+1)/(b_++r)]\le e^{-\mathcal L/3},
\]
using \(b_+\ge k\) and \(k\ge 3\). For \(J\ge b_+\)
the discard is zero. Lastly
\(D_k(B-M)/D_k(B)\ge1-rM/B\) by the binomial product.
The two counts now give the lower bound in
\eqref{eq:ordinary-class} with coefficient
\begin{equation}\label{eq:rank-fraction}
 (1-e^{-\mathcal L/8})(1-e^{-\mathcal L/3})(1-rM/B).
\end{equation}
It exceeds \(1/2\) because \(\mathcal L\ge24\) and
\(kM/B\le1/30\). This completes the proof of
\eqref{eq:operational-seg} at the stated depths.

\end{proof}

\paragraph{Multiplication around the cycle.}

The ranks just estimated belong to individual segments. On closing the cycle the retained space has the product rank, so each discarded fraction must be charged before multiplying.

On the cycle let \(R_I\) be the rank of the retained digit
classes on segment \(I\), \(U\) the two-letter tensor space,
and \(g\) the number of neighboring edges. Denote by
\(\kappa_P\) the margin loss
\[
 \kappa_P=C_\kappa N\left(e^{-\mathcal L/8}
   +\frac{kQ(m_0+N_1 k)}{b'}+\frac{1}{Y}+\frac{kE_0}{a}\right),
\]
where \(C_\kappa\) is absolute.

\begin{lemma}[Product of retained ranks]\label{lem:cycle-retained-rank}
For each segment of the cycle impose the degree separation
and reserves of Lemma~\ref{hyp:retained-digit-degrees}.
The retained ranks satisfy
\[
 \prod_I R_I\ge(1-\kappa_P)\frac{\dim U}{g}.
\]

This is the digit-class loss to be added to the exceptional
fractions of Proposition~\ref{prop:segment-discard}.
\end{lemma}

\begin{proof}
\proofstep{The losses on neighboring segments.}

Let \(I\) index a segment, \(Q_I=p^{G_I}\), and let \(a_0\)
denote the actual degree at its starting even top. Summing
over the digits gives the full monomial dimensions at each
odd top, as in Section~\ref{sec:segment-window}. At the start,
\(B=\lfloor(a-E_0)/Q_I\rfloor\) and \(a\le a_0\le a+2a/(kY)\).
For every \(1\le j\le r\),
\[
 \frac{Q_I(B+j)}{a_0+j}\ge 1-z,\qquad
 z=\frac{2a/(kY)+E_0+Q}{a}.
\]
Indeed the difference between denominator and numerator
is at most \(2a/(kY)+E_0+Q_I\), and \(Q_I\le Q\).
In a nonvacuous lower bound \(z<1\), hence
\(Q_I^r D_k(B)/D_k(a_0)\ge1-rz\).
Every even start and every even end occur once.

Also
\(\prod_I Q_I=p^{n'}\) and the number of neighboring edges
is \(g=1+p^{n'}+\cdots+p^{rn'}\ge\prod_I Q_I^r\).
Multiplying the start factors and \eqref{eq:rank-fraction}
therefore gives
\begin{equation}\label{eq:retained-product}
 \prod_I R_I\ge(1-\kappa_P)\dim U/g,\qquad
 \kappa_P=C_\kappa N\left(e^{-\mathcal L/8}
   +\frac{kQ(m_0+N_1 k)}{b'}+\frac{1}{Y}+\frac{kE_0}{a}\right).
\end{equation}
Here \(C_\kappa\) is absolute. In detail, per segment use
\(rM/B\le2kQ(m_0+1)/a\), \(e^{-\mathcal L/3}\le
e^{-\mathcal L/8}\), and \(rz\le2/Y+k(E_0+Q)/a\).
For nonnegative losses whose sum is below one,
the product of their complements is at least one minus
their sum. Otherwise the assertion with enlarged
\(C_\kappa\) is automatic. This argument explains why a
fixed positive retained fraction per segment alone would
not give the required rank.

\proofstep{The specialization of the cycle parameters.}
For the cycle with parameters from \eqref{eq:retained-product}, take
\[
 E_0=C_{E1}\left(\frac{a}{p^2}+\frac{b'}{p}
       +\frac{\mathcal L N_1a}{pk}+\frac{\mathcal L b'}{k}\right)
       +C_{E2}k^2N_1Q^2c,\qquad
 m_0=C_m kN_1Qc+12c.
\]
In \eqref{eq:retained-product} take \(k^3Q^2cY\le b'\),
\(k^3QY\le b'\), \(kb'\le a\), \(\mathcal L Y b'\le a\),
\(8\log(NY)\le\mathcal L\le k\), and \(p\ge k^2Y^2\).
The exponential is at most \(1/Y\). The three parts of the
margin term, after division by \(N_1/Y\), are bounded by
\(C_m,12,1\) respectively.

After the same division each of
the four \(C_{E1}\) terms in \(kE_0/a\) is at most \(C_{E1}\).
The last term is at most \(C_{E2}\) since \(k^3Q^2cY\le b'\le a\).
It follows that
\(\kappa_P\le C_\kappa(C_m+4C_{E1}+C_{E2}+15)NN_1/Y\).
  All these estimates use the
same \(c,m_0\) and digit classes as the pairing tests.
\end{proof}

\subsection{The normalized pairing of two digit symbols}
\label{sec:block-pairing}
\hypertarget{source-block-pairing}{}

Ranks count the retained symbols, but combining the segment tests at an edge also requires their pairings. The two models use factorials at different degrees. Here we
compare their pairings before lifting off the symbol. All
valuations in this calculation have \(v(p)=1\) and \(w^2=p\).

\paragraph{Factorial units.}

To compare the coarse and true factorial pairings, separate the unit part of the factorial of an exponent with prescribed low digit.

For \(P=p^d\), \(d\ge1\), and \(0\le h<P\), write
\[
 (h+Pn)!=h!\,n!\,(P!)^n C_P(n)D_{P,h}(n),\qquad
 D_{P,h}(n)=\prod_{1\le j\le h}(1+Pn/j).
\]
Here \(C_P(0)=1\) and
\(C_P(n+1)/C_P(n)=\prod_{1\le j<P}(1+Pn/j)\),
as is seen by cancelling the last multiple of \(P\) in the
factorial quotient. Every \(P/j\) in these products belongs
to \(p\mathbb Z_p\). Both factors are units.

\begin{lemma}[Approximating inverse factorial units]
\label{lem:factorial-unit-polynomial}
For any finite product of the factors \(C_P(\mu_i)^{\pm1}\) and
\(D_{P,h}(\mu_i)^{\pm1}\), with the indicated restrictions on \(P,h\),
and any positive integer \(c\), write \(A(\mu)\) for that product. One polynomial \(P_c\) works for every nonnegative integral vector \(\mu\),
\[
 \deg P_c\le4c,\qquad P_c(\mu)\equiv A(\mu)\pmod{p^c}.
\]
 Its coefficients in the product
binomial basis belong to \(\mathbb Q\cap\mathbb Z_p\).
\end{lemma}

\begin{proof}
\proofstep{Polynomial interpolation of the factorial units.}
We need a uniform polynomial for the inverse units modulo
\(p^c\), with degree controlled independently of the radices,
digits and number of variables. We include a finite difference
proof. In each product \(B(x)=\prod_{j\le h}(1+Px/j)^{\pm1}\),
the coefficient of \(x^e\) has valuation at least \(e\). This
also holds for negative exponents by the convergent geometric
series, uniformly on \(\mathbb Z_p\). For the forward difference
\(\nabla f(x)=f(x+1)-f(x)\),
\[
 v(\nabla^l(B-1)(x))\ge\max(1,l)
       \quad (l\ge0,\ x\in\mathbb Z_{\ge0}).
\]
Indeed the constant term of \(B-1\) is zero and the \(l\)-th
difference kills powers of degree less than \(l\).

For \(F=C_P^{\pm1}\) we have \(\nabla F=(B-1)F\) with \(h=P-1\).
The shifted product rule gives
\[
 \nabla^j F(x)
   =\sum_{l=0}^{j-1}\binom{j-1}{l}
       \nabla^l(B-1)(x)\,\nabla^{j-1-l}F(x+l).
\]
Using the bound on \(B-1\) and induction on \(j\), each term
has valuation at least \(\max(1,l)+(j-1-l)/2\ge j/2\). The same bound holds for
\(D_{P,h}^{\pm1}\) by the coefficient estimate. Repeated use of
the product rule preserves this bound for products in one
variable.

For a product of such functions in several variables
the mixed difference of order \(\nu\) at zero has valuation
at least \(|\nu|/2\). Newton's finite difference formula at a
nonnegative integer vector \(\mu\) is exact with indices
\(\nu\le\mu\), by binomial inversion in each coordinate.
Truncating to \(|\nu|\le4c\) thus gives a polynomial
\[
 H(\mu)=\sum_{|\nu|\le4c} q_\nu\prod_v\binom{\mu_v}{\nu_v},
 \qquad q_\nu\in\mathbb Q\cap\mathbb Z_p,
\]
congruent to the desired product modulo \(p^c\) for every
\(\mu\). A variable on which the unit is constant occurs with
degree zero. The decomposition applies to any finite product of digit factors.

\end{proof}

\begin{established}[Injectivity of the block monomial embedding]\label{hyp:digit-block-embedding}
Embed the block variables in the \(k\) coordinates of the
two letters at each interior top and at the next even top.
The \(r\) ordinary coordinates occupy columns 1 through \(r\).
At a tie the extra coordinates \(Z,S\) occupy column zero in
the vector and covector, respectively, including at the final
pair. A block monomial \(X^\mu\) maps to \(X^{I(\mu)}\) by
multiplying its exponents at interior top \(i\) by \(P_i\),
then adding the digit vectors. At an untied interior top also
add \(P_i(A_i-w_i),P_i h_i\) in the two columns zero.
There are no free variables in those two columns.
Writing \(\mathcal M_B\) for the monomials in block \(B\), the map satisfies
\[
 I_B(\mu)=I_B(\nu)\Longrightarrow\mu=\nu,\qquad
 I_B(\mathcal M_B)\cap I_{B'}(\mathcal M_{B'})=\varnothing\quad(B\ne B').
\]
\end{established}
\begin{proof}
At the
end the radix is one and the digits are zero.  On each ordinary coordinate the recovery is Euclidean division. If \(m\) is its true exponent at top \(i\), then
\[
 m=P_i q+\delta,\qquad q=\lfloor m/P_i\rfloor,\quad \delta=m\bmod P_i.
\]
At an untied top the
quotients in column zero recover \(w_i,h_i\). All remaining
coarse exponents are then recovered by division by the radices.
At any tie the difference of the exponents of \(Z\) and \(S\)
is its decoration \(z\), because each power of
\(E=ZS+x\cdot y\) has weight zero for that difference.
This recovers \(o_i=A_i-z_i\) or \(o_+=a_+-z_+\).
Thus different blocks have disjoint monomial supports and
\(I\) is injective within a block.
\end{proof}

\begin{lemma}[The scale of the digit block]\label{lem:digit-block-scale}
For the embedding of Lemma~\ref{hyp:digit-block-embedding}, the factor
\(\Phi\) in \(I(\mu)!=\Phi\,\mu!\,U(\mu)\) is independent
of \(\mu\). For the scale \(\gamma_I\) of \eqref{eq:digit-block-scale} one has
\[
 0\le v(\gamma_I)\le(k+1)L_I.
\]
 These bounds use the
same valuation \(v(p)=1\) as the digit pairing.
\end{lemma}

\begin{proof}
\proofstep{The degrees and the common scalar factor.}
Write \(g_x\) for the sum of the exponents in the free
variables belonging to a particular true letter \(x\), before
inserting digits. It is constant on all the levels of the
block. At an untied pair these sums are \(w_i,C_i-h_i\), at a
tie they are \(A_i,C_i\), and at the end they are \(a_+,b_+\).
Indeed each term of \(E\) has bidegree \((1,1)\), and each
\(Z^z\) restores its prescribed decoration. Write \(D_x\) for
the true degree after adding digits and the two fixed columns.
Put
\[
 b_*=\prod_{i<\sigma}(A_i-w_i)!\,h_i!,\qquad
 T=b_*\prod_x g_x!,\qquad V=v_p(T),\qquad
 L_I=\sum_x\lfloor\log_p D_x\rfloor .
\]
Use a zero log for a zero degree. We have \(b_*\mu!\mid T\)
on this block, by multiplying the letter multinomials.

Let \(U\) be the product of the \(C_PD_{P,h}\) over the free
coordinates at the odd top indices.
The factorial-unit identity gives, for a fixed integer
\(\Phi>0\) depending on this block and all its ambient degrees,
\[
 I(\mu)!=\Phi\,\mu!\,U(\mu).
\]
The factors independent of \(\mu\)
in \(\Phi\) are the factorials of the two fixed column
exponents at each untied top and the digit factorials in the
free coordinates. The remaining factor is
\(\prod_x(P_x!)^{g_x}\), independent of \(\mu\) by the preceding
homogeneity.

Write \(\partial^{[\nu]}\) for the divided derivative.
Apply the polynomial just proved to \(U^{-1}\),
with variables restricted to the odd tops. It gives the
diagonal operator
\[
 \mathcal H=\sum_\nu q_\nu X^\nu\partial^{[\nu]},
 \qquad \mathcal H X^\mu=H(\mu)X^\mu .
\]
Define
\begin{equation}\label{eq:digit-block-scale}
 Z_I=\frac{p^{kL_I}}{\prod_x D_x!},\qquad
 \gamma_I=Z_I p^V\Phi/b_* .
\end{equation}
Then
\[
 0\le v(\gamma_I)\le(k+1)L_I.
\]
To see the
cancellation in this assertion, form an exponent \(\mu^0\)
by putting each undecorated letter degree in its first ordinary
coordinate, all decorations in \(Z\) and zero in \(S\).
This uses level zero and has the right letter degrees,
irrespective of whether it obeys the prefix cuts.

In each
letter \(\mu^0\) has at most two nonzero entries. The
integer \(\prod_xg_x!/\mu^0!\) is a product of binomials,
of valuation between zero and \(L_I\) by the carry bound.
The true multinomials in \(\prod_xD_x!/I(\mu^0)!\) have
valuations summing to at most \((k-1)L_I\).
Legendre's formula proves both bounds, since at each positive
power of \(p\) the difference of the floor sums is at most one
for two entries and \(k-1\) for \(k\) entries. As
\[
 v(I(\mu^0)!)=v(\Phi)+v(\mu^0!),
\]
subtraction gives
\[
 v(\gamma_I)=kL_I+
 v\!\left(\frac{\prod g_x!}{\mu^0!}\right)
 -v\!\left(\frac{\prod D_x!}{I(\mu^0)!}\right).
\]
Both factorial quotients in this formula have nonnegative
valuation. Their upper bounds from the selected monomial
are \(L_I\) and \(kL_I\), respectively. Substituting them
gives the two endpoints for \(v(\gamma_I)\); in particular
\(\gamma_I\) is nonzero and integral.

\end{proof}

\paragraph{The two symbols.}
For two inputs \(u,u'\) in the factorial lattice of the block
take \(s(u)=w^V\mathcal H u\), \(s'(u')=w^V u'\).
As \(b_*\mu!\mid T\), the defining monomial inequality of
that lattice makes \(w^V u\) coefficient-integral. The same
holds after applying \(\mathcal H\).

We verify membership in the symbol lattice at the required
margin. First the integral span of the generators at a fixed
margin is saturated inside its rational span and the
coefficient lattice. Write a rational combination in levels
as a sum of \(\prod E_i^{j_i}Z_i^{z_i}\) times polynomials in
the reduced ordinary letters and the residual \(\xi\)'s.
Set every \(Z_i=1\) and then substitute
\[
 S_i\mapsto S_i-x_i\cdot y_i,
\]
including the last pair. These integral substitutions
take \(E_i\) to \(S_i\). Extracting the coefficients of
\(\prod S_i^{j_i}\xi^\alpha\) recovers separately the
constrained ordinary-letter polynomials. If the original
combination was integral these polynomials belong to the
coefficient lattice of their constrained rational spaces,
and reassembling them gives integral generators.

\begin{proposition}[The two lifts reproduce the block pairing]
\label{prop:digit-symbol-pairing}
Write \(B(u,u')\) for the normalized diagonal Fock pairing of the two block symbols. Then
\[
 B(u,u')-\gamma_I b_*\langle u,u'\rangle\in\gamma_I w^c O.
\]
 This holds for inputs in the two coefficient
lattices, including the allowed denominators. If the generators satisfy
the raw-tail assertion \eqref{eq:raw-tail-required}, the simultaneous
lifts preserve this congruence and give the two integral tests in
\eqref{eq:two-letter-tensor-rank}.
\end{proposition}

\begin{proof}
\proofstep{Compatibility with a variation of one coordinate.}
Next, replacing one coordinate \(v\) by \((1+t)v\) and taking
\([t^b]\) realizes \(v^b\partial_v^{[b]}\).
On the constrained letter part it increases each allowed first
row by at most \(b\). In fact in a prefix frame change a
rescaling of one coordinate becomes a change by \(t\) times a
rank-one endomorphism on that letter. Its degree-\(a\) term
is a sum of polynomials with \(a\) variables replaced and
\(a\) derivatives.

In any conjugate frame, replacing one
variable of a given vector (or covector) letter by another
changes the first torus-coordinate weight by at most one.
This proves the prefix-weight bound for the whole orbit, hence
the highest-weight cutoff. The integral divided-derivative
formula preserves integrality. Rescaling an odd coordinate also
changes at most one \(E_i\) and its adjacent \(Z_i^{z_i}\).
If a term of \(E_i\) is hit, write
\[
 E_i\mapsto E_i+tM.
\]
For an ordinary variable \(M=x_h y_h\).
For \(Z_i\) or \(S_i\),
\[
 M=Z_iS_i=E_i-x_i\cdot y_i.
\]

After expanding, removing \(l\) powers of \(E_i\) inserts
\(l\) ordinary bilinears. Each has vector and covector degrees
one and first-coordinate weight at most one in every prefix
frame (also at the prefix between its letters). Moving from
level \(j_i\) to \(j_i-l\) supplies precisely those degrees.
The power of \(Z_i\) contributes just an integer binomial
coefficient. Thus \([t^b]\) relaxes every cutoff by at most
\(b\) and preserves the span of the integral generators at the
relaxed cutoffs. Iteration costs at most \(|\nu|\le4c\) in
every cutoff. Since
\[
 m_0(C_m)-m_0(13)=(C_m-13)kN_1Qc\ge4c
\]
for \(C_m\ge17\),
both symbols belong to the required integral span at
\(m_0(13)\).

\proofstep{The diagonal Fock pairing.}
The Fock form on monomials has diagonal \(I(\mu)!\).
For all inputs in one block, including those with denominators,
we now have
\[
 Z_I\langle I(s(u)), I(s'(u'))\rangle
   \equiv \gamma_I b_*\langle u,u'\rangle
                  \pmod{\gamma_I w^c O}.
\]
Indeed \(H(\mu)U(\mu)\equiv1\pmod{p^c}\), the two factors of
\(w^V\) give \(p^V\), and each
\(b_*\mu!\,u_\mu u'_\mu\) is integral by the two defining
coefficient inequalities. Distinct blocks pair to zero by
disjoint support. Products of the symbol maps over the segments
multiply the pairings since their true letters are disjoint,
and \(Z_{\mathrm{ten}}=\prod Z_I\). The diagonal multiplier
has depth at most
\[
 (k+1)\sum_I L_I\le4\kappa_0,
\]
since the \(2N\) degrees
are at most \(2a\), \(\log_p(2a)\le2\log_p a\), and
\(\kappa_0=2Nk\log_p a\). This gives the segment-to-symbol
identity, including its uniform relative congruence.

To use the identity on two lifted polynomials it suffices that
their normalized edge substitutions be coefficient-integral
and congruent to the respective product symbols modulo \(w^c\).
On the true letter degrees the monomial diagonal of
\(Z_{\mathrm{ten}}\langle\ ,\ \rangle\) is integral, by the
same multinomial carry bound. Bilinearity then preserves the
congruence modulo \(w^c\). \paragraph{Orders, corrections and duality.}
Sections~\ref{sec:column-lifts} and~\ref{sec:symbol-correction}
construct the lifts used here. The order of a
block at a segment and the order of a tuple are
\[
 o_I=o_++\sum_i\left(P_i u_i+\sum_{h=1}^r\rho_{ih}+\tau_{i0}\right),
 \qquad o=\sum_I o_I .
\]
Let \(S_I=\sum_i\min(b_i,a_{\rm start}/p^{d_i})\). Write
\(R=a/p^2+b'/p+\mathcal L a/(pk)+\mathcal L b'/k+kQ\).
Then \(o_I\le S_I+30(s+1)R\). Indeed the digits in each summand
cost at most \(kQ\), and
\(P_i K_i\le a/p^{d_i}+2b'/p\) since \(P_i B\le a/p^{d_i}\)
and \(P_i\sum_{l<i}A_l\le (4/p^2)P_i A_i\le8b'/p^2\).

In the initial list the minimum in \(S_I\) is \(b_i\) by
the excluded band. The retained upper bound for \(u_i\) gives
\[
 P_iu_i\le b_i(1+4/p^2)+(\mathcal L/k)P_i K_i.
\]
In the tied list the minimum is at least \(a/p^{d_i}\), and
\(P_i u_i\le P_i K_i+2\mathcal L b'/k\).
For the end, put
\[
 C_{\rm end}=a/p^2+3b'/p+
       (\mathcal L/k+4/p^2)4a/p+2\mathcal L b'/k .
\]
The increasing distances are odd, \(P_i\ge p\), and \(G\ge2\).
Hence \(B\le a/p^2\), \(\sum A_i\le3b'/p\) and
\[
 \sum 2a/P_i\le4a/p
\]
by summing the geometric series. The
last retained cutoff implies \(o_+\le C_{\rm end}\). These bounds
give the claimed constant for \(p\ge5\). Summing them and using
\(\sum(s+1)\le N\) proves \(o\le\lfloor t_0\rfloor\)
for \(C_t\ge30\), as the additional \(Qc\) term is nonnegative.

\proofstep{The two lift maps in a basis adapted to the edge.}
Fix an edge line and use its adapted frame as the identity.
Denote by \(S_g\) the normalized substitution on polynomials
in the frame \(g\). The residue lines parametrize neighbors of
the source. Write \(A_1=13kN_1Q\) and \(M_0=12c\).
At precision \(0\le n\le c\) take the integral span of
product symbols with order at most \(o\), prefix margin
\(M_0+nA_1\), and excess
\[
 e=(c-n)A_1
\]
in their block bounds.
The relaxed bounds keep digit compatibility, the unused labels
equal to zero and the usual integer degree ranges. They require
\[
 h_i\le2\mathcal L B/k+e\ (i<\sigma),\qquad
 o_i\le (1-\Delta/4)A_i+e\ (i\ge\sigma),\qquad
 o_+\le\min(a_+,C_{\rm end}+e).
\]
Every retained block satisfies these bounds at \(n=c\).

For
the middle inequality use \(K_i\le B+\delta A_i\), the second
part of \eqref{eq:coarse-split}, and
\(\delta+\mathcal L/k\le3\Delta/4\).
Call this span \(\mathcal S(n,o)\). A polynomial is clean at
order \(o\) if it has the true multidegrees, its product by
\(w^o\) has integral coefficients in the source frame, and
its substitutions at every other residue line belong to
\(w^c O[X]\).

\proofstep{The required raw polynomial tail.}
The raw polynomial assertion is the following. For each
product generator \(b\) of \(\mathcal S(n,o)\), \(1\le n\le c\),
of actual order \(o_b\le o\), there is a clean \(F_0\)
at that order with
\begin{equation}\label{eq:raw-tail-required}
 S_1F_0-b\equiv\sum_{\epsilon=1}^{n-1}w^\epsilon b_\epsilon
 \pmod {w^nO[X]},\qquad
 b_\epsilon\in\mathcal S(n-\epsilon,o_b+\epsilon).
\end{equation}
Given this, induction lifts the entire span exactly to the
required precision. At zero use the zero polynomial. At \(n\)
lift \(b_\epsilon\) by induction and subtract \(w^\epsilon\)
times its lift. The multiplier recovers the lost source order
\(o_b\) and preserves cleanliness at every other line. It
raises the precision of the correction to \(n\). Integral
linear combinations finish the induction. At \(n=c\) the
margin is \(m_0(13)\), precisely the one obtained above, and
we can do this independently on an edge and on its dual.

\proofstep{The covector in the integral dual.}
Here is the identification of the second lift with a covector.
In polynomial coordinates let
\[
 \Xi(y)(x)=Z_{\rm ten}\langle x,y\rangle.
\]
Multihomogeneous polynomial coordinates realize
the tensor of one-row vector and covector representations,
with the monomial coefficient lattices in the corresponding
frames. The Fock form transports between a frame and its
inverse transpose. Indeed it is obtained by contracting
all indices of two symmetric tensors, summed over the
permutations in each letter, or equivalently by evaluating
the constant term of \(P(\partial)Q\). This works for both
letter orientations. On integral coefficients its normalized
form is integral in any such paired frames.

At a vertex
with frame \(b\), the dual frame is \(b^{-t}\). Along an edge
write \(D(u)=\operatorname{diag}(u,1,\ldots,1)\).
The intermediate frame is \(bD(w^{-1})\).
On the dual edge starting at the dual of the original target
the intermediate frame agrees with the inverse transpose, since
\(w^2=p\) and \(D\) is diagonal:
\[
 b^{-t}D(p)D(w^{-1})=b^{-t}D(w)=(bD(w^{-1}))^{-t}.
\] This remains true on all
Frobenius-twisted coefficient frames, since the twist on
\(b\) fixes \(p\). Thus \(\Xi\) takes the vertex and intermediate
integral lattices on the dual edge into the duals of those
on the original edge. In particular for a lattice \(L\) and an integer \(a\), linearity gives
\[
 \Xi(w^a L)=w^a\Xi(L).
\]

In particular the two normalized intermediate lattices both
carry \(w^{-A_H}\), where \(A_H\) is the sum of vector
degrees. The order multiplier \(w^{-t}\) at the vertex and the
conductor multiplier \(w^c\) at every other intermediate
lattice transport with the same exponents. Duality bijects
outgoing neighbors of the dual source with incoming neighbors
of the original target, preserving the distinguished pair.
This proves the target membership required of the second lift.
In the distinguished frames the normalized substitution of
each polynomial carries \(w^{A_H}\) times its expression in
its intermediate frame. Hence pairing its two substitutions
with \(Z_{\rm ten}\) agrees with evaluation by \(\Xi\) after applying the edge normalization
\[
 w^{A_H}\cdot w^{A_H}=w^{2A_H}=p^{A_H}.
\]
This gives the congruence used in
\eqref{eq:two-letter-tensor-rank} using
\eqref{eq:raw-tail-required}, proved in Section~\ref{sec:symbol-correction}.
\end{proof}

\subsection{Lifting the columns of a path}
\label{sec:column-lifts}
\hypertarget{source-column-lifts}{}

To obtain \eqref{eq:raw-tail-required}, one needs integral polynomials realizing the retained columns with control of the terms beyond their initial symbols. Fix a level of
a block as in Section~\ref{sec:segment-window}. The chain at this
level consists of the reduced vector and covector degrees in
dimension \(r=k-1\). Its prefixes have the first-row bounds
\(\lfloor K_l-\mu\rfloor\), where
\(\mu=M_0+nA_1\). We allow \(M_0=0\) or \(12c\).
Powers of the pairings \(E_i,E_+\), scalar decorations, and the
remaining endpoint covector monomials multiply this chain space.

\begin{paramdefinition}[Standard words of a Pieri path]\label{hyp:standard-path-words}
For the chain alone take a Pieri path as in
Section~\ref{sec:retained-rank}, with successive dominant weights
\(\lambda(t)\) starting at zero. For each integer column \(b\)
put
\[
 h_b(t)=\#\{j:\lambda_j(t)\ge b\}.
\]

In a positive column keep an exterior tensor of degree \(h_b\)
in the vector representation. For \(b\le0\) use the dual
representation and degree \(r-h_b\). Begin with 1. When the degree
rises, wedge on the right with the letter, and when it falls,
contract on the right with the letter. Monotonicity and interlacing
give only changes by one. In the dual exterior tensor a covector is
wedged and a vector contracted. When the height stays fixed, do
nothing. In the ordered subset basis, inserting or deleting index
\(j\) contributes the sign
\[
 (-1)^{\#\{u\in J:u>j\}},
\]
with \(J\)
the subset after omitting the inserted or deleted index.

Choose an output subset in each column. Complement these subsets
in the nonpositive columns, then require the resulting columns,
sorted increasingly, to form a weakly increasing tableau along
each row. Only finitely many columns act, since far to the left
the complemented column is full and far to the right it is empty.
The product of the selected exterior coordinates is the standard
word of the path and tableau. It is integral and of exactly the
specified multidegree, since a strip uses as many distinct columns
as its size.
\end{paramdefinition}

\begin{lemma}[Integral Pieri path words]
\label{lem:integral-path-basis}
The words of Definition~\ref{hyp:standard-path-words}
form an integral basis in the specified multidegree. If \(\mathcal W_{\rm pre}\) is the subset of words satisfying the prefix conditions and \(U_{\rm pre}\) the corresponding rational subspace of the lattice \(L\) tensored with \(\mathbb Q\), then
\[
 L\cap U_{\rm pre}=\bigoplus_{w\in\mathcal W_{\rm pre}}\mathbb Z w.
\]

\end{lemma}

\begin{proof}
\proofstep{The triangularity of the path basis.}
Here is a basis proof which also explains the importance of using
integral words. In a putative linear relation over a field insert
after each letter \(t\) an independent formal diagonal matrix
\[
 D_t=\operatorname{diag}(T_{t1},\ldots,T_{tr})
\]
acting on every
live exterior tensor, and after all letters insert one generic
invertible matrix \(g\). The relation still holds over the field
of rational functions. Indeed wedge and contraction commute with
a simultaneous change of basis, using its contragredient on dual
tensors. Moving the matrices past the subsequent operations to
the inputs changes each letter by the product of matrices following
it. This is one substitution, common to all paths.

To find the largest exponent vector of \(D_t\) on a fixed
path, take the initial \(h_b\) indices in a positive column \(b\)
and the last \(r-h_b\) indices in a dual column. If \(I_b(t)\)
is the chosen set and \(\varepsilon_b=1,-1\) for the two column
signs, the exponent of \(T_{ti}\) in the product is
\[
 \lambda_i(t)=\sum_b \varepsilon_b\,\mathbf1_{i\in I_b(t)} .
\]
Each chosen exponent vector is lexicographically maximal in its
column, and a strict decrease in any column strictly decreases
the sum in that order. Thus the maximal term forces every
choice \(I_b(t)\). These choices are compatible at consecutive
steps.

In the one changing column the old and new sets differ
by adding or deleting a single index \(i\); the corresponding
matrix entry of wedge or contraction by the input letter is
its \(i\)-th coordinate times a sign. Hence the coefficient
at the tuple of all these maxima is a signed monomial in the
inputs times the terminal matrix coefficient in \(g\). Index the insertions and deletions in the selected columns by \(a\), and let \(v_{a,i_a}\) be the input coordinate at occurrence \(a\). Their path factor is
\[
 \varepsilon\prod_a v_{a,i_a},\qquad \varepsilon\in\{1,-1\}.
\]
All dependence on the terminal tableau lies in the matrix coefficient. The matrix coefficient is a
product of minors with initial rows.

For a dual column, the
complementary minor identity introduces a sign and one inverse
determinant. After a common determinant power clears these for a
fixed path, the result consists of ordinary flag minors with
the tableau columns. In lexicographic order on the matrix entries,
row by row with smaller column index preferred, a flag minor with sorted column indices \(i_1<\cdots<i_h\) has leading monomial
\[
 \operatorname{in}\det(g_{a,i_b})_{a,b=1}^h=\prod_{a=1}^h g_{a,i_a}.
\]
The determinant signs from the dual columns change only the overall sign. The diagonal
monomial recovers each sorted tableau row. Consequently these
matrix coefficients are independent for each path, in every
characteristic. Choose the lexicographically largest sequence of
weights in the relation, comparing the diagonal matrices in
temporal order. The coefficient just computed forces all the
coefficients for that sequence to vanish. This proves independence.

\proofstep{The rank and prefix operators.}
Let \(p_\lambda\) count paths ending at \(\lambda\), and let
\(d_\lambda\) count its semistandard tableaux. For integer weights
the tableau can be shifted by a whole determinant power, which
does not change \(d_\lambda\). Pieri gives
\[
 \sum_\lambda p_\lambda d_\lambda
       =\prod_j\binom{r+o_j-1}{r-1},
\]
the monomial rank for the prescribed letter degrees \(o_j\).
The integer coefficient matrix is therefore square. Write \(A\) for this matrix. Independence in each residue characteristic gives
\[
 \det A\not\equiv0\pmod p\quad(\text{every prime }p),\qquad \det A\in\mathbb Z^\times=\{1,-1\}.
\]


At a prefix \(t\), varying the first \(t\) letters by a common
matrix factors their contribution through the tensor of the exterior
powers at that cut. In each torus coordinate its largest possible
weight is at most \(\lambda_1(t)\), by the same signed subset
calculation. The word therefore belongs to the first-row filtered
space at that bound. Over characteristic zero the multiplicity-one Pieri decompositions give
\[
 \dim\bigcap_t\{\text{prefix cutoff at }t\}=\#\{\text{basis words whose paths satisfy every cutoff}\}.
\]
Indeed one
may first split a prefix tensor into irreducibles and then tensor
each summand with the next letter. The later isotypic projections
preserve these summands. The independent allowed words consequently
span that intersection. Since they form a subset of an integral
basis, they also span its intersection with the coefficient lattice.
Expanding at each level and using the monomial basis for the
endpoint covector gives the desired generators of the integral
symbol space.

\end{proof}

\paragraph{Putting back the first row.}
Write \(q_j=p^{G_j}\) for the radix at the start of segment \(j\),
\(B=\lfloor(a-E_0)/q_j\rfloor\), and \(z_i=A_i-w_i\) or
\(A_i-o_i\) according as \(i<\sigma\) or \(i\ge\sigma\).
Prepend an even vector letter. Enlarge the path by a new first
row, initially zero, which at that first letter rises to
\[
             b_0=\lfloor B-\mu\rfloor-\sum_{i<\sigma} h_i .
\]
Increase this row by \(z_i\) at the odd vector \(i\). Keep it
fixed at the covectors and at the final vector. The reduced
path is unchanged in the other rows, delayed by the prepended
step.

For every vector the first reduced row after the addition
must be at most the enlarged first row \emph{before} the addition.
At a tied vector and the final vector this first row before
addition is exactly \(\lfloor K_i-\mu\rfloor\) and
\(\lfloor K_s-\mu\rfloor\), respectively. The initial \(h_i\)
are all accounted for in \(b_0\).

For the initial vectors the reduced first row is at most
\(\sum_{l\le i}w_l\). Put \(\epsilon_0=N_1\mathcal L/k\) and
\[
 e_0=(c-n)A_1.
\]
In the local size range
\(b'\ge k^3NQ^4cY,\; kb'\le a,\;N\ge2,\;Y\ge1\).
The floor estimate
\[
 B\ge a/(2Q)
\]
gives
\[
 \epsilon_0B\ge k^3Q^3 c,
\]
using \(\epsilon_0\ge1/k\).
Since \(2N_1+3\le k\) and \(k\ge5\),
\[
 \mu+N_1e_0\le12c+13kN_1^2Qc\le4\epsilon_0 B,\qquad
 \sum_{i<\sigma}h_i\le 2\epsilon_0B+N_1e_0 .
\]
In particular \(b_0\ge0\).

By \eqref{eq:coarse-prefix},
\(\sum_{l\le i} w_l\le(1+4/p^2)A_i\).
Using \((1+\Delta)A_i<B\), \(4/p^2\le\epsilon_0\),
\(40\epsilon_0\le\Delta<1/6\), we get
\[
 (1+4/p^2)A_i+6\epsilon_0 B\le B.
\]
Thus the required interlacing holds at every initial vector too.
The first row stays fixed at a removal, for which old dominance
and reduced interlacing suffice. This constructs a valid path in
dimension \(1+r\), including for an empty odd list.

\paragraph{Frobenius powers in exterior coordinates.}
For a column of the enlarged path perform the same exterior
recursion on the true letters, with the following insertions.
For successive used Frobenius indices \(s,s'\), the intervening operation on an exterior coordinate \(c_I\) is
\[
 c_I\longmapsto c_I^{\,p^{s'-s}}.
\]
Use the terminal index of the segment for the last such operation. Use increasing integer lifts of
the positions through wrap-around. The vector and covector at
one odd position have the same index. The result is integral.
At a used position \(s\) its degree in that letter is
\(p^{s_{\rm end}-s}\), and zero at unused positions. Positions
of the same orientation in a segment are distinct, including
the start and end by the two-segment hypothesis. Each true
letter degree in this column is therefore at most \(Q\).

Use coordinate 0 for the added line and the remaining coordinates
for the reduced tensor. We claim that the new column has order
at least
\begin{equation}\label{eq:lifted-path-column-order}
 d_b=\sum_{s:\,\text{reduced column height rises at }s}
                         p^{s_{\rm end}-s}.
\end{equation}
at every residue line. Its symbol after division by \(w^{d_b}\)
at the identity is the reduced column with every letter raised
coordinatewise to its radix, times the pure first-coordinate
power at the step at which the new first row crosses \(b\),
if it crosses at all. For a positive column include coordinate
0 in the output subset exactly when the first row reaches \(b\).
For a nonpositive column keep the dual output subset in the
remaining coordinates. Use the ordering already specified, so
right wedge and right contraction on ordinary coordinates have
the reduced signs.

\begin{lemma}[Order of a lifted path column]
\label{lem:column-graded-lift}
For the enlarged path at this block level, let \(C_\ell\) be its column at a residue line \(\ell\), with integral coefficient lattice \(L_\ell\). The order \(d_b\) in \eqref{eq:lifted-path-column-order} satisfies
\[
 C_\ell\in w^{d_b}L_\ell\qquad(\text{each residue line }\ell).
\]

At the identity line its symbol modulo \(w\) after this division is
the reduced column evaluated on the coordinatewise radix powers,
multiplied by the crossing scalar when the first row is crossed. The
symbol is homogeneous of normal degree \(d_b\).
\end{lemma}

\begin{proof}
\proofstep{The orders after the Frobenius powers are lifted.}
We give the order proof since taking Frobenius powers in a
moving frame requires two orders of precision. Call a subset
allowed in a positive column of positive height when it contains
0 and has that height, and at height zero when it is empty.
In a dual column the allowed subsets omit 0 and have the
complementary height. Measure the entries in the adapted frame
of the tested line after each step.

The induction asserts that
all are divisible by \(w^d\) and the entries outside the allowed
subsets by \(w^{d+1}\). Initially this holds for the scalar 1.
Between indices let \(P=p^{s'-s}\) and let \(L_s\) convert actual
exterior coordinates into adapted coordinates, using the
appropriate twisted exterior matrix. Write the adapted row as
\(w^d U\). After entrywise powering and conversion at \(s'\) it is
\[
         w^{Pd}\big(U^{[P]}+pV\big)
\]
for an integral polynomial row \(V\). Indeed expand
\((U L_s^{-1})^{[P]} L_{s'}\) modulo \(p\), where powering is
additive and
\[
 (L_s^{-1})^{[P]}\equiv L_{s'}^{-1}.
\]

The last congruence holds in the unramified Witt ring before
extension to \(O\), for inverse minors as well as minors.
At \(s'=s\) the identity is exact. If the index strictly
increases, the extra \(w\) in a disallowed entry becomes at least
\(w^2\) after the common factor is removed, since
\[
 p=w^2
\]

and \(P\ge2\).

In a positive column a first-row crossing wedges the free
vector coordinate 0 into an empty exterior tensor. For a reduced height increase the tensor already contains \(e_0\). Since
\[
 e_0\wedge(e_0\wedge v)=0,
\]
a surviving insertion must use an ordinary coordinate carrying \(w\).
A height decrease contracts an ordinary covector at no cost.
Contracting coordinate zero costs an extra \(w\) and is disallowed.
For a nonpositive column the complemented tableau always contains
the new first row. A decrease wedges an ordinary covector. An
increase contracts an ordinary vector at cost \(w\). These exhaust
the possible degree changes.

At either costly step a free coordinate acting on a disallowed entry could consume the spare power of \(w\). Every costly step is a vector
at a strictly later index than all preceding used steps, as the
vector precedes the covector at an equal index. After the common factor is removed, the disallowed term has order at least two. Paying the one power for a costly step therefore still leaves the required extra order:
\[
 \operatorname{ord}_w(\text{disallowed term})-1\ge2-1=1 .
\] If it is the first used step,
all the disallowed entries are zero. At a free step the extra
order is preserved directly. If \(d_{\rm old}\) is the common order before a step and
\(P\) is its intervening Frobenius radix, we have proved
\[
 d_{\rm new}=P d_{\rm old}+
       \mathbf1_{\{\text{reduced height rises at this step}\}} .
\]
Iterating from order zero weights each costly step by its
remaining radix \(p^{s_{\rm end}-s}\), exactly
\eqref{eq:lifted-path-column-order}. Inverting the final
integral frame change preserves this common divisibility.

\proofstep{Returning to the identity frame.}
At the identity the frame matrices are identities. In the
normalized leading part the steps just described act on the
ordinary exterior coordinates exactly as in the reduced word.
Crossing the first row occurs while the ordinary positive
column is empty, and contributes its scalar vector entry.
If \(P\) is the intervening Frobenius power, then in the residue field
\[
 \Big(\sum_I\varepsilon_I a_I b_I\Big)^P=\sum_I\varepsilon_I a_I^P b_I^P\qquad(\varepsilon_I\in\{1,-1\}).
\]
Thus applying Frobenius before the exterior operation gives the same residue. The resulting ordinary column is thus evaluated
on the stated radix powers.

The leading expression in
characteristic zero lifted this way is homogeneous of normal
degree \(d_b\), counting ordinary vector coordinates and scalar
covector coordinates with weight one. The argument proves a
congruence modulo \(w\) after normalization, as required.
Taking the product of the columns reproduces all the reduced
words and spends \(q_j b_0\) at the start and \(P_i z_i\)
in vector coordinate zero at each odd top.
\end{proof}

\subsection{Completing and correcting the polynomial}
\label{sec:symbol-correction}
\hypertarget{source-symbol-correction}{}

The lifted columns prescribe the initial part of a polynomial. Completing it at the shared endpoints introduces error terms whose orders must be estimated before the edge tests are combined.

Use \(x\) at an even start, \(v,y\) at an odd top, and let
\(H_i=p^{d_i}\), so \(H_iP_i=q_j\). The other factors are
formed from the integral polynomials
\[
 K_i=\sum_{a=0}^{k-1}x_a^{H_i}y_a,\qquad
 W_{ih}=v_hx_0^{H_i}-v_0x_h^{H_i}\quad(1\le h<k),\qquad
 D_i=\sum_a v_a y_a .
\]
Here \(K_i\) denotes a polynomial only in this paragraph and the
factor table below. The capacities elsewhere are those of
Section~\ref{sec:segment-window}.

\begin{lemma}[Initial symbols on an adapted line]
At the identity line \(w^{-1}S_1K_i\) and \(w^{-1}S_1W_{ih}\)
have leading terms \(x_0^{H_i}y_0\) and \(v_h x_0^{H_i}\),
respectively, with remainders divisible by \(w^{H_i-1}\).
Both unnormalized polynomials vanish modulo \(w\) on every
adapted line. For the same-position pairing one has
\[
 S_gD_i=wD_i.
\]

\end{lemma}
\begin{proof}
At the identity, ordinary vector and scalar covector coordinates
acquire a factor \(w\). Term by term,
\[
  S_1K_i\equiv w x_0^{H_i}y_0,\qquad
  S_1W_{ih}\equiv w v_hx_0^{H_i}\pmod {w^{H_i}} .
\]
At another adapted line the vector modulo \(w\) is proportional
to the Frobenius twist of the first inverse row. Raising the
start row to \(H_i\) gives the twist at the odd top, since the
residue frame has the cycle period. This kills the minor \(W_{ih}\).
The ordinary covector columns of the forward frame, with indices
other than zero, annihilate that row, killing \(K_i\). In the
same-position pairing the frame matrices cancel exactly, giving
the asserted formula for \(D_i\).
\end{proof}

At a tie of level \(j_i\) use \(D_i^{P_i}\) a total of \(j_i\)
times. Its normalized symbol is \(\sum_a(v_ay_a)^{P_i}\) by
Frobenius modulo \(w\). The last even pairing \(D_+^{j_+}\) has
exact normalized symbol \(D_+^{j_+}\).

\begin{lemma}[Raw polynomial tail at a distinguished line]
\label{lem:raw-tail-correction}
For \(1\le n\le c\) and a product generator \(b\) of
\(\mathcal S(n,o)\) of actual order \(o_b\le o\), there is a clean
polynomial \(F_0\) of that order for which
\[
 S_1 F_0-b\equiv\sum_{e=1}^{n-1}w^e b_e\pmod{w^nO[X]},
 \qquad b_e\in\mathcal S(n-e,o_b+e).
\]
The assertion uses the bounds and margins in
Section~\ref{sec:block-pairing}.
\end{lemma}

\begin{proof}
\proofstep{The missing exponents in the true weights.}
Writing \(\rho_{ih},\tau_{ih}\) for the vector and covector
digits, the remaining order-bearing factors are
\[
\begin{array}{c|c|c}
 \text{factor}&\text{multiplicity}&\text{order per factor}\\ \hline
 K_i^{P_i}&h_i\ (i<\sigma)&P_i\\
 W_{ih}&\rho_{ih}\ (h>0)&1\\
 K_i&\tau_{i0}&1 .
\end{array}
\]
Set \(h_i=0\) at ties when using this table.
The free digits \(v_0^{\rho_{i0}}\prod_{h>0}y_h^{\tau_{ih}}\)
and the endpoint covector monomial of degree \(b_+-j_+\) are
unchanged in the identity substitution and integral at all lines.

Choose one even start \(j_0\). Using the vector letters \(x_j\)
at all starts, put
\[
 C_h=\prod_j x_{j0}^{q_j-1}-\prod_j x_{jh}^{q_j-1},
 \qquad R=x_{j_0,0}\prod_{h=1}^{k-1} C_h ,
\]
where the subscripts denote the start and coordinate.
At the identity \(C_h\) has order zero and its symbol is the
first monomial. At another line write \((a_h)_h\) for the first
inverse row over the residue field and \(s_j\) for the start indices.
Each substituted \(x_{jh}\) reduces to \(x_{j0}a_h^{p^{s_j}}\).

For \(a_h\ne0\) we have
\[
 \prod_j(a_h^{p^{s_j}})^{q_j-1}=1,
\]
since the factors are
successive conjugate ratios around the cycle. If \(a_0=0\)
the initial factor in \(R\) vanishes. If \(a_0\ne0\) there is
some \(h>0\) with \(a_h\ne0\), since this is a different line.
Then \(C_h\) vanishes. Thus \(R^c\) supplies the extra \(c\)
orders at all competitors. In counting perturbations regard
each \(C_h\) as a separate factor repeated \(c\) times and
put the pure power at \(j_0\) with the exact monomial factors.

\paragraph{An endpoint shared by adjacent segments.}

An endpoint shared by two segments has only one degree to spend. The contributions from both adjacent paths must be included when that degree is restored.

The starting scalar degree spent at \(j\) is
\[
 U_j=q_j\Big(b_0+\sum_i h_i\Big)
 +\sum_i H_i\Big(\sum_{h>0}\rho_{ih}+\tau_{i0}\Big)
 + c\big((k-1)(q_j-1)+{\bf1}_{j=j_0}\big).
\]
Here \(b_0\) belongs to segment \(j\). The available degree
is \(a_{+,j-1}-o_{+,j-1}\), the decoration at the end of
the preceding segment, with cyclic indices. The first term is
at most \(a-E_0\) by the definition of \(b_0\). Each summand
in the second term is at most \(kq_j\), since each digit is
less than \(P_i\). The last term is at most \(ckQ\).

For the relaxed bounds of Section~\ref{sec:block-pairing} we have
\(o_{+,j-1}\le C_{\rm end}+c A_1\). The reserve
\[
 E_0=C_{E1}\left(a/p^2+b'/p+
     \mathcal L N_1a/(pk)+\mathcal L b'/k\right)
      + C_{E2} k^2 N_1 Q^2 c
\]
with both constants at least 17 covers
\(C_{\rm end}+cA_1+N_1 kQ+ckQ\).
Indeed the terms of \(C_{\rm end}\) have respective coefficients
at most \(5,3,4,2\) in the first parentheses for \(p\ge5\).
The other three terms satisfy
\[
 cA_1+N_1kQ+ckQ\le15kN_1Qc.
\]

Since \(a_{+,j-1}\ge a\), the unused scalar degree is nonnegative.
Fill it by the pure power of \(x_{j0}\).

\proofstep{The common degrees at the endpoint.}
The product \(F\) now has exactly the true multidegrees. This can
also be checked directly on its leading symbol. The columns give
the chain word with its radix substitutions and the odd powers
\(v_0^{P_i z_i}\). The factors in the table give the remaining
digits and the initial covector scalar powers. The tied and final
pairings give the powers of \(E_i,E_+\), and filling at all starts
gives exactly the final decorations. Every factor and its symbol
have identical letter degrees. The sum of the orders in the
columns counts the reduced vector degrees with their radix
multipliers. Restoring the levels and adding the normal digits and
the initial scalar covector powers gives
\[
 o_b=\sum_{\text{segments}}\left\{
 \sum_i\left[P_i u_i+\sum_{h>0}\rho_{ih}+\tau_{i0}\right]+o_+
 \right\},
\]
the order of the generator.

We have proved that \(F_0=w^{-o_b}F\)
is clean, has integral substitution at the chosen line, and that
its normalized symbol there is the generator modulo \(w\).
The extra \(c\) at another line is added to the full \(o_b\)
because the universal orders of all the factors are still available there.

\paragraph{The error of a product.}
Consider a variable factor of the construction, of order \(d\)
and leading symbol \(s\). Its normalized error is
\(w^{-d}S_1 f-s\). For a monomial of normal degree \(a\) its
coefficient is divisible by
\[
 w^{\max(1,a-d)}.
\]
Indeed
the symbol is homogeneous of normal degree \(d\), and the
diagonal substitution on the integral polynomial supplies
\(w^a\) on this monomial. The exact factors, including the final
pairing and its residual covector, have zero error. Every variable
factor has degree at most \(Q\) in each true letter, and degree
zero in the endpoint covector. The three degree bounds used here are
\[
 H_i P_i=q_j\le Q\quad(K_i^{P_i}),\qquad P_i\le Q\quad(D_i^{P_i}),\qquad q_j-1\le Q\quad(C_h).
\]


\proofstep{Terms indexed by the error factors.}
Expand the product of normalized factors, choosing errors in a
nonempty multiset of factors. Expand those errors in monomials.
A term modulo \(w^n\) may be written \(w^e\) times an integral
scalar times the product of the remaining symbols and one monomial,
where \(1\le e<n\). Take \(e\) to be the sum of the displayed
lower bounds on its individual coefficient orders, ignoring terms
with this sum at least \(n\). At most \(e\) factors have been
replaced. The monomial replacing them has their total
multidegrees and normal degree at most their total order plus
\(e\). The entire term consequently has normal degree at most
\(o_b+e\).

In each segment separate the remaining reduced column product
\(G'\), of degrees \(d_{vi},d_{yi},d_{v+}\), and the surviving
pairing powers \(E_i^{j'_i},E_+^{j_+}\). Here
\[
 0\le j'_i\le j_i
\]

at the ties, and both are zero in the initial list. In the
true letters all remaining pieces comprise a single monomial
\(X^A\). Before the replacements its exponents at an odd top were
\[
 (\rho_{i0}+P_i z_i,\ (\rho_{ih})_{h>0}),\qquad
 (\tau_{i0}+P_i h_i,\ (\tau_{ih})_{h>0}).
\]
At the endpoint these were the scalar vector decoration
\(a_+-o_+\) and the residual ordinary covector exponents
\(\gamma_h\). Write \(\Delta_h\) for the signed change in the exponent of coordinate \(h\) of a given true letter. The two directional bounds are
\[
 \sum_h(\Delta_h)_+\le Qe,\qquad \sum_h(-\Delta_h)_+\le Qe,\qquad q=Qe.
\]
For each replacement one removes its monomial
part, possibly only the pure scalar part of a column, and inserts
a monomial of degree at most \(Q\). If part of the factor was kept
separate in \(G'\) or the \(E_i\)'s, it is only the former loss
which becomes smaller. All endpoint covector exponents remain
exactly as before.

\proofstep{The additional order from deleted columns.}
There is a second effect of deleting columns. The largest
first-coordinate weight in a generic prefix frame is additive on
products of nonzero polynomials. Apply the generic change of
basis to the prefix and then its diagonal torus at the first
coordinate. Over the resulting coefficient domain, Laurent degree satisfies
\[
 \deg_T(FG)=\deg_T F+\deg_T G\qquad(FG\ne0).
\]
Indeed the product of the two leading coefficients is nonzero. Each deleted reduced column
has degree at most one in each of the at most \(N_1\) covectors.
Its largest weight is at least \(-N_1\), a lower bound even
for each of its monomials. Thus \(G'\) has every prefix cutoff
of the full word relaxed by at most \(N_1 e\).

\paragraph{Recovering a block from the error.}
Divide each odd-top exponent of \(A\) by \(P_i\), using its
remainders as the new digits. Write \(Z,Y\) for the quotients
in coordinate zero of the vector and covector, and
\(V_h,T_h\) for the ordinary quotients, at the top under
consideration. Write \(V=\sum_{h>0}V_h\) and \(V_+\)
for the total ordinary vector exponent at the endpoint of
radix one. The unchanged true degrees give new coarse totals
\[
 A'_i=j'_i+d_{vi}+Z+V,\qquad
 C'_i=j'_i+d_{yi}+Y+\sum_{h>0}T_h .
\]

In particular the new digit class is compatible and
\(A_i-A'_i\le k\), since two sums of \(k\) remainders
differ by less than \(kP_i\). We have \(V,V_+\le q\),
\(Y\le h_i+q\), and \(Z\ge z_i-q\). For example, each
positive unit of an ordinary vector quotient costs at least
one unit of positive variation from a digit less than \(P_i\).
The other assertions follow by the same division with remainder.

At an initial top set \(w'_i=d_{vi}+V\), \(h'_i=Y\).
Then \(w'_i\le w_i+q\), \(h'_i\le h_i+q\), and the two
new labels are bounded by \(A'_i,C'_i\), respectively. At a
tie set
\[
              o'_i=j'_i+d_{vi}+Y+V .
\]
Here \(j'_i+d_{vi}\le o_i\) and \(Y\le q\). In order to
have \(o'_i\le A'_i\) we need \(Y\le Z\). This uses the
excluded band quantitatively. For any compatible digits,
\(A_i\ge(11/12)b'/Q-k\), since \(b_i\ge(11/12)b'\).
Using \(b'\ge2k^3Q^4c\), \(\Delta\ge40/k\),
\(2N_1+3\le k\) and \(k\ge5\), we obtain
\[
 \Delta A_i/4\ge (55/3) k^2 Q^3 c-10
              \ge c(A_1+2Q+3k)+4 .
\]

\proofstep{The degree reserves for the reconstructed block.}
In the second inequality use \(A_1\le13k^2Q\),
\(Q,c\ge1\). The tied cutoff
\(o_i\le(1-\Delta/4)A_i+(c-n)A_1\) therefore gives
\(z_i=A_i-o_i\ge 1+2q\). Thus \[
 \begin{gathered}
 Y\le Z, \\ o'_i\le A'_i, \\ o'_i\le o_i+2q.
\end{gathered}
\]
Set the unused labels to zero. Finally put
\(o'_+=j_++d_{v+}+V_+\). It is bounded by \(a_+\)
by homogeneity and by \(o_++q\) by the reduced final degree.
The relaxed block inequalities now all hold at precision
\(n-e\). In particular the old tied upper bound loses at
most \(k+2q\) on changing \(A_i,o_i\), and the extra
allowance is \(eA_1\).

To realize the term in this block, multiply \(G'\) by all the
ordinary quotient monomials in the corresponding coarse letters,
including those at the endpoint. At an initial top the scalar
quotients already agree with the required decorations. At a
tie factor the scalar powers as
\[
 Z_i^{Z}S_i^{Y}
       = Z_i^{Z-Y}(E_i-v_i\cdot y_i)^Y ,
\]
now writing \(v_i,y_i\) for the ordinary coarse letters and
\(Z_i,S_i\) for their extra scalar coordinates.
Here
\[
 Z-Y=A'_i-o'_i.
\]

Binomial expansion gives new levels
\(j'_i+m_i\), \(0\le m_i\le Y\), and
\(Y-m_i\) new ordinary pairings. The identities for
\(A'_i,C'_i\) prove that both letter degrees at each such level
are exact and that its level is allowed. They also check the
degrees in the initial list. At the endpoint the scalar vector
power is exactly \(a_+-o'_+\), and the residual covector still
has degree \(b_+-j_+\).

\proofstep{The order of the reconstructed top term.}
An ordinary covector monomial cannot increase the top weight
of a prefix, and the inserted ordinary vector degrees cost at
most \(s q\) up to the odd tops. A new ordinary pairing costs
nothing unless the prefix stops between its two letters, where
it costs at most its power
\[
 Y-m_i\le q.
\]
Only one such pair
can straddle a prefix. If the endpoint vector is included, no
interior pair straddles the cut and the endpoint costs at most
\(q\). The total added prefix weight is therefore at most
\((s+1)q\), in addition to the deletion cost \(N_1e\).
Meanwhile each increment \(A_i-u_i\) in the capacity decreases
by at most \(k+2q\). The capacity can fall by at most
\(s(k+2q)\) at any prefix.

All three losses fit inside the
margin \(eA_1\), since
\[
    N_1e+(s+1)q+s(k+2q)\le13kN_1Qe=eA_1 .
\]
All costs here are integers, so taking floors after subtracting
the margin preserves the inequality. This proves membership in
the required integral chain lattice at the reduced precision.
The order of the new block equals the normal degree of the
term. Each occurrence of a coarse ordinary vector contributes its radix
and every pairing contributes that radix once. The digits supply
the remaining ordinary vector and scalar covector exponents.
Alternatively this follows by inserting the reconstructed labels
in the displayed formula for \(o_b\). The global monomial
decomposes by its odd tops and the endpoint of each segment,
each true letter being included exactly once. Thus an error term \(T_e\) of order \(e\) satisfies
\[
 w^{-e}T_e\in\mathcal S(n-e,o_b+e).
\]
This proves
\eqref{eq:raw-tail-required} and completes the simultaneous lift
by the induction in Section~\ref{sec:block-pairing}.

The constants in these inequalities are absolute under the stated local size bounds.
Neither the number of columns nor the number of competing
lines appears in the loss at one precision step.
\end{proof}


\subsection{What the neighboring lattices measure}

The polynomial construction supplies vectors and covectors on individual edges. To convert them into a rank at a lattice, distinguish edges entering that lattice from edges leaving it.

Let \(W\) be the ring of integers in the unramified degree
\(f\) extension of \(\mathbb Q_p\), with residue cardinality
\(q_f=p^f\), and let \(O\) be the coefficient DVR in the local
setting with uniformizer \(w\) satisfying \(w^2=p\).
The selected cycles under consideration have even \(f\).
For a rank-\(k\) lattice \(L\) over \(W\), an edge
\(e=(L_-,L_+)\) satisfies \(pL_+\subset L_-\subset L_+\)
with \(L_+/L_-\) one-dimensional over the residue field.
For fixed \(L_+=L\) these edges are the hyperplanes of \(L/pL\).
For fixed \(L_-=L\) they are the lines of \(p^{-1}L/L\).
Both numbers are
\[
                       g=(q_f^k-1)/(q_f-1).
\]

At selected positions \(s\) the dominant weights have endpoints
\(a_s,-b_s\), with bounded intermediate entries. All entries
are bounded at the remaining positions. At odd positions replace
the weight by its negative reversal. Denote the resulting weights
by \(\mu_s\) and put
\[
 A_1=\sum_s\mu_{s,1},
\]
using first
coordinate index one. Let \(V\) be the tensor product of the
characteristic-zero dual Weyl spaces for the partitions
\(\mu_s-\mu_{s,k}{\bf1}\). For \(L\) write \(H(L)\) for the
tensor of their integral lattices in the Frobenius-twisted
component frames, including the determinant twist \(\mu_{s,k}\).
Equivalently in each component take the integral polynomial
lattice in every basis of the given module.

In the component at
\(s\), write \(L^{(s)}\) for the \(O\)-span of its image under
the \(s\)-th Frobenius embedding. Define \(H(e)\) by applying
the same construction to the intermediate component lattices
\(L_-^{(s)}+w L_+^{(s)}\). The two normalized modules and their
pairing are
\[
 R_e=w^{-A_1}H(e),\qquad R'_e=w^{-A_1}H(e)^\vee,\qquad
 B(x,y)=p^{A_1}y(x).
\]
They pair integrally, since \(p^{A_1}=w^{2A_1}\).
Write \(e^-,e^+\) for the source and target of a directed edge
\(e\). At a chosen order \(t\), use the source and target modules
\[
 N_-(L)=w^{-t}H(L)\cap\!\bigcap_{e^-=L}R_e,\qquad
 N_+(L)=w^{-t}H(L)^\vee\cap\!\bigcap_{e^+=L}R'_e .
\]
For an edge \(e\) put \(C_e\subset N_-(e^-)\) equal to the
submodule satisfying in addition \(x\in w^cR_{e'}\) at every
other outgoing edge. Define \(C'_e\subset N_+(e^+)\) by the
analogous condition at every other incoming edge.

\begin{construction}
\[
\begin{array}{ccccc}
L_-&\lhook\joinrel\longrightarrow& L &\lhook\joinrel\longrightarrow &L_+\\
\multicolumn{2}{c}{pL\subset L_-\subset L}&&\multicolumn{2}{c}{L\subset L_+\subset p^{-1}L}\\
\multicolumn{2}{c}{\text{hyperplanes in }L/pL}&&\multicolumn{2}{c}{\text{lines in }p^{-1}L/L}\\
\multicolumn{5}{c}{g=1+q_f+\cdots+q_f^{k-1}\quad\text{choices on each side}}
\end{array}
\]
\caption{The incoming and outgoing families of a vertex have the same cardinality. Incoming tests are combined at a common target, whereas the outgoing rank bound counts directions at a common source.}\label{fig:incident-lattice-edges}
\end{construction}

\begin{proposition}[Combining tests on incident edges]\label{prop:incident-edge-tests}
Suppose that for each edge there are \(m_e\) vectors in
\(C_e\) and \(m_e\) covectors in \(C'_e\) whose pairing matrix
has an integral right multiple equal to \(z_e I_{m_e}\), with
\(z_e\ne0\) and \(v(z_e)\le D\). Here \(D\ge0\), \(c>2D\), and zero values of \(m_e\) are permitted. If
\[
                m_e\ge (1-\kappa)\dim(V)/g
\]
at every incoming edge, then the tests at target \(L\) recover
at least \((1-\kappa)\dim(V)\) directions with an integral
inverse after multiplication by \(w^{\lfloor 2D\rfloor}\). A uniform number \(m\) of directions
recovered on every outgoing edge at a fixed source also satisfies
\(gm\le\dim V\).
\end{proposition}

\begin{proof}
\proofstep{Block normalization at a target.}
Apply the indicated integral column operations within
each block and multiply the covectors of that block by
\(w^{\lfloor 2D\rfloor}/z_e\in O\). The diagonal blocks of the
combined pairing matrix become \(w^{\lfloor 2D\rfloor} I\).
At distinct edges with the same target
the covector from the second block belongs to \(w^cR'_e\) at
the first. Every off-block entry therefore belongs to \(w^c O\).
For the combined matrix \(M\), the remaining divisibility is
\[
 w^{-\lfloor2D\rfloor}M=I+w^{c-\lfloor2D\rfloor}E\equiv I\pmod w,\qquad E\in M_{\sum m_e}(O).
\]
Its inverse is integral by the adjugate formula.
Combining the covectors using this inverse gives covectors in
\(N_+(L)\) evaluating to \(w^{\lfloor 2D\rfloor}\) times the coordinate vectors.
There are \(\sum_{e^+=L}m_e\ge(1-\kappa)\dim(V)\) such
coordinates. All source vectors retain their required extra
vanishing on the other outgoing edges. \proofstep{Independence at a source.}
The same proof at a fixed source uses the extra source vanishing to put the
off-block entries in \(w^cO\). Its combined matrix has full
rank over the fraction field, so a uniform number \(m\) that
can be recovered at every outgoing edge satisfies
\(gm\le\dim(V)\). This explains the divisor \(g\) in the
single-edge rank estimate, even when \(q_f\) grows with the
cycle.

\end{proof}

For the final integer depth \(D\), take \(c=8(D+1)\). This meets
the lower bound \(4(D+1)\) imposed on the neighbor test.

\paragraph{Attaining the discarded fraction.}
The source takes the infimum of the admissible discarded
fractions over the inequality parameter, with the required
rank test quantified over all targets, coefficient settings,
and allowed bounded modifications of the weights. This infimum
is attained. The fraction one is admissible by taking zero
directions. For each fixed test, if \(d_V=\dim V>0\) and the
infimum is \(u\), take an admissible \(u'\) sufficiently close
to \(u\) that
\[
 (1-u')d_V>\lceil(1-u)d_V\rceil-1.
\]

Integrality of the rank gives the test at \(u\). The zero-rank
case is immediate. The vectors realizing the required rank may be chosen separately for each test. Thus for \(p\ge p_0\) the same \(\kappa_v\) which bounds the local loss actually supplies the test dimensions used in the column index.

\subsection{Integral transfer from the tensor of two-letter spaces}
\label{sec:edge-transfer}
\hypertarget{source-edge-transfer}{}

Each top in the tensor model has independent vector and covector variables. Transferring its tests to the actual weight space requires an integral comparison with the irreducible summands.

Let \(T\) be the set of tops on the cycle and \(O\) the coefficient valuation ring. Put \(N=|T|\), \(k=m+1\), and \(K=\operatorname{Frac}(O)\).
Use \(A_s,B_s\) for the two nonnegative degrees at a top.
The tensor model \(U\) takes independent vector and covector
variables at each top. At an odd top their roles switch, so the
first-coordinate normalization of \(U\) is
\[
 a_U=\sum_{s\ {\rm even}}A_s+\sum_{s\ {\rm odd}}B_s.
\]
Use the component frames of the preceding subsection for the
coefficient lattice \(U(M)\), for a tuple \(M\) of \(O\)-lattices.
Apply the definitions of the edge tests to \(U\) with \(a_U\) in
place of \(A_1\), including the pairing \(p^{a_U}\phi(u)\).

In one component let \(P_{A,B}\) be the polynomials of bidegree
\((A,B)\) in \(x,y\), and let \(h=\sum_i x_i y_i\) and
\[
 \Delta=\sum_i\partial_{x_i}\partial_{y_i}.
\]

Writing \(\beta_{A,B}\) for the quotient by \(hP_{A-1,B-1}\),
this quotient is the dual Weyl model for
\((A,0,\ldots,0,-B)\), with its determinant twist.
The identification holds on coefficient lattices. To see the
normalization, send \(x\) to the first column and \(y\) to the
signed \((k-1)\)-row minors of the first \(k-1\) columns. Use
independent matrices for the two-letter space and one matrix
for the quotient. Under a change of the row frame, the column
transforms as a vector and the minors as a covector times the
determinant. The twist cancels the determinant.

\begin{proposition}[Transfer of a tensor test to the weight summands]\label{prop:edge-test-transfer}
Take \(k\ge3\), \(w^2=p\), and \(A_s+B_s+1\le8a\) at each top.
Put \(u=N\lceil(k+1)\log_p(8a+k-1)\rceil\). Suppose a test for
\(U\) has \(r_U\) directions at depth \(D_0\) at a fixed edge,
with the specified endpoint order and conductor. Then there are
tests with \(r_j\) directions on the summands \(V_j\) of weights
\((A_s-j_s,0,\ldots,0,-B_s+j_s)\),
\(0\le j_s\le\min(A_s,B_s)\), at depth \(D_0+3u\), with
the same endpoint order and conductor and total rank bounded below by
\[
 \sum_j r_j\ge r_U,\qquad D_0+3u\le D_1=D_0+4u.
\]
In particular the tests are available at depth \(D_1\).
\end{proposition}

\begin{proof}
\proofstep{The injective substitution and differential operators.}
In the independent model the substitution is injective over
every field, since each nonzero covector is the exterior
product of some frame of its kernel, with one frame variable
rescaled to give the specified exterior product. In the
common-matrix model the image is dense in the incidence
hypersurface \(h=0\). Over an algebraic closure, for any
nonzero \(x,y\) with \(y(x)=0\), extend \(x\) to a basis of
\(\ker(y)\) and adjust one of the other vectors to realize
\(y\). Here \(k\ge3\). The polynomial \(h\) is irreducible, by
viewing it as primitive of degree one in any \(y_i\); hence the ideal of relations in the common-matrix substitution is
\[
 \ker(\text{common-matrix substitution})=(h).
\]
These arguments also apply over the residue
field after base change to its algebraic closure.

Over \(O\)
one may divide by the monic leading monomial of \(h\), so the
bihomogeneous quotient is free. The injection after reduction
modulo \(w\) proves saturation of its image. In the standard
frame the integral polynomial realization of the dual Weyl
module is its span over \(K\) intersected with the coefficient
lattice. Equivalently one uses the integral standard monomial
basis in products of column minors \cite[Theorem II.2.16]{abw-schur}.
The Cartan product of the
first-column and cofactor weights generates that module, giving
the quotient identification. A change of frame gives the assertion for every \(M\).

\proofstep{Projectors on fixed bidegree.}
For \(a,b\ge0\), let \(H_{a,b}\) be the projector onto
\(\ker\Delta\) along \(hP_{a-1,b-1}\).
The commutator formula
\[
 \Delta(h^i u)=h^i\Delta u+i(k+c+d+i-1)h^{i-1}u
\]

for \(u\in P_{c,d}\) gives
\[
 H_{a,b}=\sum_{i=0}^{\min(a,b)}
    \frac{(-1)^i h^i\Delta^i}{i!(k+a+b-2)^{\underline i}},
 \qquad t^{\underline i}=t(t-1)\cdots(t-i+1).
\]
Indeed \(\Delta H=0\), and the constant summand gives the
identity modulo \(h\). If \(u=h^j v\) with \(\Delta v=0\),
\(\Delta^j u=j!(k+A+B-2j)^{\overline j}v\), where the bar
denotes a rising factorial. Induction on \(\min(A,B)\) gives
the unique decomposition into the \(h^j\ker\Delta\).
For \(a=A-j,b=B-j\), its projection and inclusion in
quotient coordinates are
\[
 \pi_j f=\beta_{a,b}\frac{H_{a,b}\Delta^j f}
                        {j!(k+a+b)^{\overline j}},
 \qquad
 \iota_j\beta_{a,b}(u)=h^jH_{a,b}u.
\]
The formula for \(\pi_j\) follows by applying \(\Delta^j\)
to each harmonic summand and killing the terms still divisible
by \(h\).

Both maps have denominator at most
\(p^{(k+1)\lfloor\log_p(A+B+k)\rfloor}\) in every component
frame. For \(\pi_j\), in the \(i\)-th summand of its formula
expand \(\Delta^{i+j}\) by multi-indices \(\nu\) with
\(|\nu|=i+j\). Apart from integral binomial coefficients
from the two differentiated monomials, the coefficient is,
up to sign,
\[
 \frac{\binom{i+j}{i}^2}
 {\binom{i+j}{\nu}\binom{k+a+b-2}{i}\binom{k+a+b+j-1}{j}}.
\]
The factorial formula for a valuation sums one term per \(p\)-power.
For the multinomial that term is at most \(k-1\); each binomial
contributes at most one. All factorial arguments are at most
\(A+B+k\). For \(\iota_j\) the projector formula has the same
estimate with \(j=0\), and an integral lift of the quotient
exists by its coefficient-lattice identification. Both maps
commute with the paired vector and covector action, so one
bound works for all frames.

\proofstep{Tensoring with control of denominators.}
Tensor these maps over the tops. Let \(j=(j_s)_s\),
\(0\le j_s\le\min(A_s,B_s)\), and let \(V_j\) be the weight
space with endpoints \(A_s-j_s,-(B_s-j_s)\). Write
\[
 \pi_j:U\to V_j,\ \iota_j:V_j\to U
\]
for the products of
the maps just constructed. Factoring the sum independently at each top gives
\[
 \sum_{\mathbf j}\iota_{\mathbf j}\pi_{\mathbf j}
 =\bigotimes_s\left(\sum_{j_s=0}^{\min(A_s,B_s)}
        \iota_{s,j_s}\pi_{s,j_s}\right)
 =\bigotimes_s 1=1_U ,
\]
where the map with first index \(s\) is the corresponding single-top
map.

When \(A_s+B_s+1\le8a\), the common lattice allowance is
\[
 u=N\left\lceil (k+1)\log_p(8a+k-1)\right\rceil.
\]
Thus \(p^u\pi_j U(M)\subset H_j(M)\) and
\(p^u\iota_j H_j(M)\subset U(M)\). Put \(a_j=A_1(V_j)\)
and \(\delta_j=a_U-a_j=\sum_s j_s\).
The maps with the normalizations of the edge tests are
\[
 \widetilde\pi_j=w^{\delta_j}\pi_j,\qquad
 \widetilde\iota_j=w^{-\delta_j}\iota_j .
\]
If \(x,\phi\) satisfy the two tests at an edge for \(U\), then
\begin{equation}\label{eq:edge-test-maps}
 p^u\widetilde\pi_j x\in C_e(V_j),\qquad
 p^{2u+\delta_j}\phi\circ\widetilde\iota_j\in C'_e(V_j).
\end{equation}
For source vectors the factor \(w^{\delta_j}\) changes the
edge scale \(w^{-a_U}\) into \(w^{-a_j}\), including at every
other outgoing edge.

At the source itself it only increases
integrality. For covectors use
\(p^u(\phi_0\circ\iota_j)\in H_j(M)^\vee\) whenever
\(\phi_0\in U(M)^\vee\), which follows by evaluating against
\(p^u\iota_j H_j(M)\). The uniform covector estimate in this argument allows
\(2u\) in this step. The factors
\(p^{\delta_j}w^{-\delta_j}=w^{\delta_j}\) give the same
edge and endpoint scales on the dual side. In the normalized
pairing the total effect is exactly \(p^{3u}\), since
\[
 p^{a_j}p^{2u+\delta_j}p^u\,
 \phi(\widetilde\iota_j\widetilde\pi_j x)
       =p^{3u}\,p^{a_U}\phi(\iota_j\pi_j x).
\]
Thus the strength of the extra vanishing at the other edges
is preserved. The exponent \(\delta_j\) may grow with the
weights and is canceled in the pairing.

Let \(zI_{r_U}\) be the matrix of the tensor test after its
integral right multiplication. Apply \eqref{eq:edge-test-maps}
and use the same right multiplier on every summand. This produces
integral matrices \(M_j\) of size \(r_U\) with
\[
                    \sum_j M_j=z_*I_{r_U},
             \qquad z_*=p^{3u}z\ne0 .
\]
For each \(j\), let \(r_j\) count the Smith diagonal entries
dividing \(z_*\), and let \(E_j\subset O^{r_U}\) be the span
of the corresponding left Smith vectors. Over the DVR all
other entries are multiples of \(wz_*\). Consequently the images
of the \(M_j\) in \(O^{r_U}/\sum_j E_j\) are contained in
\(w\) times the image of \(z_*O^{r_U}\). Their sum contains
the image of \(z_*O^{r_U}\), so that image is zero by Nakayama's
lemma. This gives
\[
  r_U=\operatorname{rk}_O\bigl(\textstyle\sum_j E_j\bigr)
                                       \le\textstyle\sum_j r_j .
\]

In the \(j\)-th pairing, take the retained Smith rows and columns
and multiply the right partners by \(z_*\) divided by their
diagonal entries. These factors lie in \(O\), so they preserve
all source and target membership and vanishing tests. The
resulting scalar identity has depth at most \(D_0+3u\).

\end{proof}

\subsection{Combining segment tests at one edge}

Transfer identifies the tensor model on one edge. The segment ranks multiply there only after their pairings and the associated denominator bounds have been combined.

\begin{proposition}[Assembly of the segment tests]\label{prop:segment-edge-assembly}
Use the two-letter tensor \(U\), the retained segment lattices
of ranks \(R_I\), and the segment tests of depths \(D_I\) and
ranks \(r_I\ge(1-\kappa_I)R_I\). Assume the supported
block-pairing congruence and order \(\lfloor t_0(A,B)\rfloor\)
from Sections~\ref{sec:block-pairing}--\ref{sec:symbol-correction},
and let \(D_0=\sum_I D_I+4\kappa_0\). If \(c>2D_0\), the
tensor edge test has depth \(D_0\) and at least \(\prod_I r_I\)
directions. Writing \(r_{\rm edge}\) for the rank of this test, the retained-product bound \eqref{eq:retained-product} gives
\[
 r_{\rm edge}\ge\prod_I r_I\ge\left(1-\kappa_P-\sum_I\kappa_I\right)\frac{\dim U}{g}.
\]

\end{proposition}

\begin{proof}
Choose Smith bases separately on each retained digit block. A
fractional pairing may first be multiplied by a uniformizer
power, which is then removed from the resulting diagonal entries.
If the Smith entries are \(s_i\), the low-valuation rank is
\[
 r_I=\#\{i:v(s_i)\le D_I\},\qquad v(0)=+\infty.
\]

The \(r_I\) left basis vectors span a saturated sublattice; a
primitive vector in it has a unit coordinate, so pairing against
the corresponding right basis vector gives a value of valuation
at most \(D_I\). Conversely, a saturated sublattice of larger
rank has a nonzero intersection with the rational span of the
remaining left vectors. Rescale a vector in the intersection to
be primitive in that sublattice. Every integral partner gives
valuation greater than \(D_I\). Thus the Smith count is exactly
the rank tested by primitive vectors.

Multiply the right partner for a retained entry \(s_i\) by
\[
 w^{\lfloor2D_I\rfloor}/s_i.
\]
The ratio is integral because
\(v(w)=1/2\) and \(2v(s_i)\) is an integer. We obtain diagonal
segment matrices with common value \(w^{\lfloor2D_I\rfloor}\)
and retain the original block support.

\proofstep{The rank after tensoring the segment tests.}
Let \(R_I\) be the rank of the retained lattice for segment \(I\)
and
\[
 r_I\ge(1-\kappa_I)R_I
\]
its successful rank at depth
\(D_I\). Sections~\ref{sec:block-pairing}--\ref{sec:symbol-correction} give the following identity
for a tuple of supported vectors and a tuple of supported
partners. If the support labels are \(\boldsymbol\beta\) and
\(\boldsymbol\beta'\), and the segment pairings are \(a_I\),
the normalized edge pairing is congruent to
\[
  {\bf 1}_{\boldsymbol\beta=\boldsymbol\beta'}\,
                z_{\boldsymbol\beta}\prod_I a_I \pmod{w^c},
   \qquad z_{\boldsymbol\beta}\ne0,\quad
                      v(z_{\boldsymbol\beta})\le4\kappa_0 .
\]
The lifted vectors must satisfy the tests \(C_e(U),C'_e(U)\)
at order \(\lfloor t_0(A,B)\rfloor\) and conductor \(c\).
 The scalar is allowed
to depend on the tuple of blocks. Because the segment
matrices are diagonal, the lifted matrix is diagonal modulo
\(w^c\), with nonzero diagonal values of valuation at most
\(D_0=\sum_I D_I+4\kappa_0\).

If \(c>2D_0\), write the matrix as the diagonal matrix times
an integral matrix congruent to the identity modulo \(w\).
The inverse of the latter is integral. Scaling the
diagonal values to \(w^{\lfloor 2D_0\rfloor}\) then gives a
tensor edge test on \(\prod_I r_I\) directions.

Equation~\eqref{eq:retained-product} supplies the ordinary rank count
on retained blocks.
Combining this with the tests gives
\[
 r_U\ge \left(1-\kappa_P-\sum_I\kappa_I\right)\dim U/g .
\]
In the nonvacuous case all the discarded fractions and their
sum are in \([0,1]\), so the inequality follows by multiplying
the rank bounds and using \(\prod(1-\kappa_I)\ge1-\sum\kappa_I\).
When the right side is nonpositive the zero test suffices.
This proves the required tensor rank at the chosen filtration.

\end{proof}

\subsection{Averaging the rank deficits}

The tensor space decomposes into summands of different dimensions. An estimate for the proportion discarded by the edge test must therefore average their deficits with dimension weights.

Let \(D(a)=\binom{a+k-1}{k-1}\) for \(a\ge0\), set
\(D(-1)=0\), and put
\(d(a,b)=D(a)D(b)-D(a-1)D(b-1)\).
The dimensions are \[
 \dim U=\prod_s D(A_s)D(B_s),\qquad \dim V_j=\prod_s d(A_s-j_s,B_s-j_s).
\]
Summing the latter over \(j\) telescopes to \(\dim U\).
The maximal successful rank at a fixed depth and order is
the same at every edge because \(\mathrm{GL}_k(\operatorname{Frac}W)\)
is transitive on these pairs of lattices. A frame change also
transports the neighboring edges and both endpoint tests.
Thus every summand of the transferred test has
\(r_j\le\dim V_j/g\) when \(4D_1<c\), by the previous
source rank argument.

\paragraph{Inputs to the dimension-weighted average.}
The dimension-weighted sum assumes rank bounds on its tensor constituents. We specify which sampled segment tests supply them on the enlarged box. For the enlarged weight box
below, the tensor edge test is required at depth \(D_0\)
and order \(\lfloor t_0(A,B)\rfloor\) with
\begin{equation}\label{eq:two-letter-tensor-rank}
 r_U\ge (1-\kappa_X(A,B))\dim U/g,\qquad
 \kappa_X(A,B)=\kappa_P+\sum_I\kappa_I(A,B).
\end{equation}
The sum is over segments, and \(\kappa_P\) accounts for the
digit blocks cut away by the margins. The two-letter construction supplies these tests at the base level.
Sections~\ref{sec:segment-window} and~\ref{sec:retained-rank} give \(\kappa_P\le C_s NN_1/Y\)
and, on the enlarged product box for fixed \(B\),
\[
 \mathbb E_A\kappa_I(A,B)
      \le C_A (s_I+1)N_1/Y,\qquad \sum_I(s_I+1)\le N.
\]
Here \(s_I+1\) is the number of sampled \(A\)-coordinates in
the segment, including its appended coordinate.

These coordinates have different indices in the sampling box, so their law is the
product marginal on that segment. In particular a uniform
bound for fixed \(B\) implies the same bound averaged over
both \(A\) and \(B\). This gives
\(\mathbb E\kappa_X\le(C_A+C_s)NN_1/Y\).

The function \(t_0\) is a constant independent of \(A,B\)
plus a sum of at most \(N\) terms of the form
\(\min(B_s,A_{s'}/p^e)\) with \(e\ge0\).
Let \(J=\lfloor J_0\rfloor\), where \(J_0\) is the cleaning
budget below. For \(0\le j_s\le J\), each term changes by
at most \(J_0\) under simultaneous subtraction of \(j\).
Thus the order
\[
 t_*(a',b')=\lceil t_0(a',b')+NJ_0\rceil
\]

for the lower weights is at least
\(\lfloor t_0(A,B)\rfloor\). Increasing the order only enlarges
the two endpoint lattices. Write \(F(\alpha,\beta)\) for
the least discarded fraction at depth \(D_1\) and order
\(t_*(\alpha,\beta)\) for the weight with these endpoints
and zero intermediate entries. When \(J\le\min_s(A_s,B_s)\),
the improved bound for the small \(j\)'s is
\begin{equation}\label{eq:transferred-small-deficit}
 r_j\le (1-F(A-j,B-j))\,\dim V_j/g.
\end{equation}
The comparison can now be stated entirely in terms of the
retained rank of the large tensor and the discarded fractions
of the transferred tests.

\begin{lemma}[Weighted rank deficit]\label{lem:transfer-rank-deficit}
Let \(J\le\min_s(A_s,B_s)\). Suppose the transferred ranks at
depth \(D_1\) satisfy \(\sum_j r_j\ge r_U\) and
\(r_j\le\dim V_j/g\) for every summand, with the sharper bound
\eqref{eq:transferred-small-deficit} for \(0\le j_s\le J\).
Assume the tensor lower bound \eqref{eq:two-letter-tensor-rank}. With the dimension weights of the indicated summands, the discarded fractions satisfy
\[
 \sum_{0\le j_s\le J\ (\forall s)}\frac{\dim V_j}{\dim U}F(A-j,B-j)\le\kappa_X(A,B).
\]

\end{lemma}

\begin{proof}
Using \eqref{eq:two-letter-tensor-rank} and
\(\sum_jr_j\ge r_U\) gives the weighted inequality
\begin{equation}\label{eq:edge-mean-weighted}
 \sum_{0\le j_s\le J} F(A-j,B-j)
    \prod_{s\in T}\frac{d(A_s-j_s,B_s-j_s)}{D(A_s)D(B_s)}
        \le \kappa_X(A,B).
\end{equation}
The upper rank bounds for the omitted summands are essential
to this inequality. The number \(g\) cancels before averaging.

\end{proof}

\begin{proposition}[From dimension weights to the product box]\label{prop:edge-box-deficit}
Let \(J=\lfloor J_0\rfloor\), \(e=J+c_{\rm cl}\) and
\(I_a^+,I_b^+\) be the respective upper enlargements by \(e\)
of the intervals \(I_a,I_b\) of lengths \(H_a,H_b\).
For \(0\le\tau\le c_{\rm cl}\), assume
\eqref{eq:edge-mean-weighted} on the enlarged box and the density
bound \eqref{eq:edge-marginal-density}, with \(\delta>0\), for every
coordinate pair in \((I_a+\tau)\times(I_b+\tau)\). Then
\[
 \mathbb E_{(I_a+\tau)^T\times(I_b+\tau)^T}F
 \le \delta^{-N}
    \left(\frac{(H_a+e)(H_b+e)}{H_aH_b}\right)^N
       \mathbb E_{(I_a^+)^T\times(I_b^+)^T}\kappa_X .
\]
The prefactor is less than two under \eqref{eq:edge-marginal-loss}. The assertion is uniform in the common shift \(\tau\).
\end{proposition}

\begin{proof}
Use integer intervals \(I_a,I_b\) of lengths \(H_a,H_b\)
for the original weight choice. Denote by \(I_a^+,I_b^+\)
their enlargements at the upper ends by \(e=J+c_{\rm cl}\),
where \(c_{\rm cl}\) bounds the permitted intermediate entries
and bounded shifts. For each \(0\le \tau\le c_{\rm cl}\)
and \((\alpha,\beta)\in(I_a+\tau)\times(I_b+\tau)\),
one has
\begin{equation}\label{eq:edge-marginal-density}
 \sum_{j=0}^J\frac{d(\alpha,\beta)}{D(\alpha+j)D(\beta+j)}
    \ge \delta,\qquad
 \delta=(1-(k-1)J_0/a)(1-\exp(-\Lambda_0/3)).
\end{equation}
Here \(a\) is the rounded lower endpoint of \(I_a\),
\(\Lambda_0=C_0N_1^2\log a\),
\(J_0=\Lambda_0 b_0/k\), and \(b_0\) is the reference
value for the second weight. Assume \(\alpha\ge a\),
\(\beta+J+1+(k-1)\le2b_0\), \(k\ge3\), and
\((k-1)J_0/a\le1\).

In fact
\(D(\alpha)/D(\alpha+j)\ge1-(k-1)J_0/a\)
by multiplying its reciprocal factors and applying Bernoulli's
inequality. Also
\[
 d(\alpha,\beta)/[D(\alpha)D(\beta)]\ge (k-1)/(\beta+k-1).
\]
Writing \(v_j=D(\beta)/D(\beta+j)\), the latter lower bound
times \(v_j\) dominates \(v_j-v_{j+1}\).
The final term \(v_{J+1}\) is at most
\(\exp(-(k-1)(J+1)/(2b_0))\le\exp(-\Lambda_0/3)\).
This proves the density for one top.

Write \(\mathcal B=(I_a+\tau)^T\times(I_b+\tau)^T\) and
\(\mathcal B^+=(I_a^+)^T\times(I_b^+)^T\). For every
\((\alpha,\beta)\in\mathcal B\) and \(0\le j_s\le J\), the source
\((A,B)=(\alpha+j,\beta+j)\) belongs to \(\mathcal B^+\), and each
such term occurs once for that \(j\) in the enlarged sum.
Keeping exactly those nonnegative terms in
\eqref{eq:edge-mean-weighted} gives
\[
 \sum_{(\alpha,\beta)\in\mathcal B} F(\alpha,\beta)
 \prod_{s\in T}\left(
 \sum_{j=0}^{J}\frac{d(\alpha_s,\beta_s)}
 {D(\alpha_s+j)D(\beta_s+j)}
 \right)
 \le
 \sum_{(A,B)\in\mathcal B^+}\kappa_X(A,B).
\]
This is a change of summation order over the product box and
uses no independence assumption on \(F\). Apply
\eqref{eq:edge-marginal-density} in each coordinate and divide
by \(\delta^N|\mathcal B|\) to obtain the stated prefactor. We use the
bounds below when
\begin{equation}\label{eq:edge-marginal-loss}
 N\left((k-1)J_0/a+\exp(-\Lambda_0/3)+e/H_a+e/H_b\right)
 \le 0.395 .
\end{equation}

\proofstep{Uniformity of the losses in the averaging.}
To check the constants, write \(\Lambda_*=C_0N_1^2\log a_{\rm real}\)
for the value before rounding, and
\(Z=\Lambda_* b_0/k\). The size hypotheses give
\(\Lambda_*\le\Lambda_0\le2\Lambda_*\), \(c_{\rm cl}\le Z\),
\(H_a,H_b\ge Y_bNZ/2\) and \(a\ge Y_b N\Lambda_* b_0\)
with \(Y_b\ge400\). These bound the first and last two
normalized terms by \(2/Y_b,6/Y_b,6/Y_b\).
Also
\[
 N\exp(-\Lambda_0/3)\le N(NY)^{-8/3}\le0.36
\]

for \(N\ge2,Y\ge1\).
Bernoulli bounds \(\delta^N\) below by \(1-s\), where
\(s=N((k-1)J_0/a+\exp(-\Lambda_0/3))\).
The enlargement factor is at most \(\exp(t)\), with
\(t=N(e/H_a+e/H_b)\). The numerical bounds give
\[
 e^t\le(1-t)^{-1},\qquad s+t\le0.395.
\]
Hence the prefactor
in the last mean is less than two.

\end{proof}

\subsection{The integral contraction for the intermediate weights}
\label{sec:edge-contraction}

The segment calculation uses weights with only the extreme entries varying. An integral contraction compares these tests to the prescribed weights with bounded intermediate entries. At one top put \(d=c_{\rm cl}\)
and \(r=k-2\). In this subsection write \(V_\nu\) for a Schur
space in \(k\)-row matrix polynomials with column degrees
\[
 \nu_0\ge\cdots\ge\nu_{k-1}\ge0.
\]
It consists of polynomials
transforming by \(\prod_i b_{ii}^{\nu_i}\) under right
multiplication of the matrix by an upper triangular matrix \(b\).
The flag minors have column sets consisting of initial columns.
The coordinate ring of \(\mathrm{GL}_k\) under left and right
translation identifies this space in characteristic zero with the
irreducible of highest weight \(\nu\).

The products of flag minors
of this shape span that irreducible, since their span is
row-stable and contains the highest vector. The integral standard
monomial theorem gives its basis over \(O\)
\cite[Theorem II.2.16]{abw-schur}. With a diagonal term order
on the flag minors the leading monomial of a standard tableau
records the entries of its rows. These monomials are distinct
because the tableau rows are weakly increasing. Thus the
span of the basis is saturated in the entire polynomial lattice.
For a weight \(\lambda\) whose last entry may be negative put
\[
 \nu=\lambda-\lambda_{k-1}{\bf 1}.
\]

The lattice in a frame with matrix \(g\) consists of
\[
 (\det g)^{-\lambda_{k-1}}f(g^{-1}X)
\]
with \(f\) in the
standard polynomial lattice. This is the determinant twist
used in \(H(L)\).

Take end weights \(A+d,-B-d\) and intermediate weight \(\mu\)
with \(|\mu_j|\le d\), and put
\begin{equation}\label{eq:edge-intermediate-weights}\begin{aligned}
 \lambda_-&=(A,0,\ldots,0,-B),&
 \lambda_+&=(A+2d,0,\ldots,0,-B-2d),\\
 S&=(d,\mu,-d),&
 S^*&=(d,-\mu_{\rm rev},-d).
\end{aligned}\end{equation}
We use a bilinear form on the \(r\)-row Schur lattices of
middle partitions
\(\alpha=B+d+\mu\) and \(\alpha'=d-\mu_{\rm rev}\).
Their entries satisfy
\(\alpha'_i=L-\alpha_{r+1-i}\) with \(L=B+2d\).
Put \(e=\alpha_r\) and \(\gamma=\alpha-e{\bf1}\).
Division by \(\det(Y)^e\) takes the first lattice to
\(V_\gamma(O)\), since each monomial of shape \(\alpha\)
has exactly \(e\) full-size minors.

For the second lattice,
with \(J_r\) the reversal matrix, use the map
\[
             f(Y)\longmapsto (\det Y)^{L-e} f(Y^{-\mathsf T}J_r).
\]
Jacobi's complementary-minor identity sends each initial-column
minor to an initial-column complementary minor divided by
\(\det Y\), up to sign. The determinant factor clears all
denominators and the output has shape \(\gamma\).
The same transformation in reverse, by its square on invertible
matrices, gives surjectivity on the integral lattice.

\begin{proposition}[Contraction across the intermediate weights]
\label{prop:integral-edge-contraction}
At each top take the weights of \eqref{eq:edge-intermediate-weights}
with \(|\mu_i|\le d\). The flat-weight Schur lattices admit
the multiplication and contraction maps through the integral
\(r\)-row Schur lattice of shape \(\gamma\) obtained from the
two complementary partitions \(\alpha,\alpha'\). For their tensor over \(N\) tops the exponent and the corresponding \(w\)-adic order are
\[
 d_T=Nd,\qquad \operatorname{ord}_w(p^{d_T})=2d_T\quad(w^2=p,\ \operatorname{ord}_w(w)=1).
\]
 After multiplication by \(p^{d_T}\), the pulled-back target
covector satisfies the target integral test and the normalized pairing
at the distinguished edge, with a loss of at most \(2d_T\) in the
\(w\)-adic conductor at the other incoming edges. The source test
holds simultaneously at all outgoing edges.
\end{proposition}

\begin{proof}
\proofstep{A uniform pairing for the intermediate representation.}
Here is a uniform integral pairing on \(V_\gamma\).
For the transpose adjoint, the rational contravariant dual has
the same highest weight as \(V_\gamma\). Choose the corresponding
isomorphism between the two rational modules to take the highest
polynomial to the highest dual functional evaluating to one
there. Both lattices have a primitive generator on that line
and are stable under divided powers. Write \(l=\gamma_1\)
and
\[
 h=\binom r2\lfloor\log_p(l+r-1)\rfloor
\]
for \(r\ge2\).
When \(r=1\) use \(h=0\).
In either lattice the submodule generated from the highest
vector by divided powers contains \(p^h\) times the full lattice.
Indeed take a basis \(u_j\) of the divided-power lattice in
that representation.

Apply the invariant Hom projector of the
preceding type calculation to \(u_j^\vee\otimes v\), where
\(v\) lies in the full lattice. The coordinate weights of
\(V_\gamma\) lie in \([0,l]\), as each symbol appears at most
once per column of a tableau. Each coordinate of a Hom weight is a difference of two such coordinates, so
\[
 0\le a,b\le l\quad\Longrightarrow\quad -l\le a-b\le l.
\]
Multiplying by \(p^h\) gives integral intertwiners,
which take the highest vector into the integral highest line
and its divided-power span into the divided-power span.
The trace identity for that projector reconstructs \(p^h v\)
as the sum of their values on \(u_j\), multiplied by
\(\dim V_\gamma\).

Let \(M\) and \(M^\sharp\) be the two full lattices and
\(L,L^\sharp\) the respective divided-power sublattices.
The contravariant isomorphism \(j\) sends \(L\) onto \(L^\sharp\).
The inclusions just proved in both lattices give
\[
 p^{2h}M^\sharp\subseteq p^h j(M)\subseteq M^\sharp .
\]
The right inclusion uses \(p^h M\subseteq L\), and the
left uses \(p^h M^\sharp\subseteq L^\sharp\).
Thus \(p^h j\) is an integral pairing map with cokernel
killed by \(p^{2h}\). Lifting \(p^{2h}\) times a dual basis
gives the integral partners with those diagonal values. This bound depends on the width
\(l\le2d\), even though the determinant power \(e\) can be
large. In particular for \(2d\le8a\) the depth is at most
\[
 \lceil (k-1)^2\log_p(8a+k-1)\rceil.
\]

Compose with determinant division and the complementary-minor
map to obtain a pairing \(J\) on shapes \(\alpha,\alpha'\),
with the same full set of partners. Simultaneous substitution
\(Y\mapsto gY\) in both inputs multiplies \(J\) by
\(\det(g)^L\). This follows from the transpose adjoint,
since the complementary-minor input then changes by
substitution by \(g^{-\mathsf T}\) and the factor
\(\det(g)^{L-e}\). The first input contributes the remaining
\(\det(g)^e\).

\proofstep{Contraction on the polynomial coefficients.}
For a \(k\)-row polynomial \(f(X)\), put
\(D(Y)=\operatorname{diag}(1,Y,1)\).
Apply the pairing \(J\) in the variable \(Y\), coefficientwise in \(X\), to define
\[
 \mathcal C(f,t)(X)=J_Y\bigl(f(XD(Y)),t(XD(Y))\bigr).
\]

Here the two \(k\)-row inputs have the normalized shapes
of \(\lambda=\lambda_-+S\) and \(S^*\).
Cauchy--Binet expands each flag minor in these substitutions
into middle flag minors of \(Y\) with coefficients in \(O[X]\)
for integral input. The \(Y\)-shapes are \(\alpha,\alpha'\)
because the first column of \(X\) is retained and the last is
unused. Thus contraction is integral. It commutes with left matrix substitution in \(X\).

For an invertible middle block \(g\) one has
\[
 D(g)D(Y)=D(gY),
\]
so the covariance of \(J\) gives
\[
 \mathcal C(f,t)(XD(g))=\det(g)^L\,\mathcal C(f,t)(X).
\]
This proves the middle diagonal degrees and invariance under
upper unipotent middle-column operations. A first-column
addition to a middle column is represented by a matrix \(u\)
whose only off-diagonal entry is in the first row. The matrix
\(D(Y)^{-1}uD(Y)\) is upper unipotent over \(K(Y)\), so
\[
 XuD(Y)=XD(Y)(D(Y)^{-1}uD(Y))
\]
has the same value in each input. Invariance there gives a
polynomial identity over \(K\). The last column is unused,
and the first-column degrees add. Together these operations
give precisely the right triangular character of the
normalized shape of \(\lambda_+\).

For any middle polynomial of shape \(d+\mu\), lift each
flag minor of size \(j\) to a flag minor including the first
row and first column, of size \(j+1\). If the product has \(a\) minors, the lifted product has first-row degree \(a\). Multiply by
\[
 X_{00}^{2d-a},\qquad 0\le a\le2d,
\]
to obtain first-row degree \(2d\).
Evaluation on \(D(Y)\) gives the original polynomial, and
the lift has shape \(S+d{\bf1}\) and last-row degree zero.
The same construction lifts the weight \(d-\mu_{\rm rev}\)
to the shape of \(S^*\). Since multiplication by
\((\det Y)^B\) identifies the shape \(d+\mu\) with \(\alpha\),
choose the lifted partners \(s_b,t_a\) so that
\[
 J((\det Y)^B s_b(D(Y)),t_a(D(Y)))
\]
is the chosen
scalar diagonal.

If \(F\) has the lower flat shape, its
substitution \(F(XD(Y))\) is \((\det Y)^B F(X)\),
by the first-column and cofactor formula.
It follows that \(\mathcal C(Fs_b,t_a)\) equals \(F\)
times a polynomial with the flat normalized shape
\((4d,2d,\ldots,2d,0)\), first-row degree \(4d\)
and last-row degree zero. This space is spanned by
\[
          Z=(X_{00}\det X_{[0,k-2],[0,k-2]})^{2d}.
\]
In its first-column and cofactor model, every variable factor
must have maximal first-row degree, and each cofactor must
omit the last row. Since \(Z(I_k)=1\), evaluation at the
identity matrix shows that
\[
 \mathcal C(Fs_b,t_a)= \delta_{ab}p^{2h}ZF.
\]

This constructs the one-component maps
\[
 i(F)=ZF,\ U_b(F)=s_b F,\ L_a(G)=\mathcal C(G,t_a).
\]

The number of pairs is \(\dim V_\gamma\), the dimension for
the middle weight. Tensor over all tops, reversing the two
end degrees at the odd tops, to obtain the maps on a cycle.

\proofstep{Verification on the edge lattices.}
We check the edge lattices for those maps. In a frame of an
edge the single changing row direction in the lattice has
scale \(1,w^{-1},w^{-2}\) at the source, intermediate
lattice and target. For first-row degree \(2d\) and last weight \(-d\) the determinant-twisted scaling factors are
\[
 w^d\quad\text{at the first change},\qquad w^{2d}\quad\text{at the second}.
\]
Thus the transported
\(s_b,t_a\) belong to the normalized intermediate lattice,
and \(w^{2d}t_a\) is integral at the target.
For \(Z\) the degree and the last weight are doubled.
It is primitive in the coefficient lattice of the intermediate
frame after normalization. To check saturation, reduce a product \(ZP\) in the coefficient
lattice modulo the uniformizer \(w\). Primitivity of \(Z\) and the
domain property give
\[
 \overline Z\,\overline P=0,\quad\overline Z\ne0
       \ \Longrightarrow\ \overline P=0 .
\]
The original coefficient lattices are saturated, so \(P/w\) then
belongs to the smaller lattice. This proves saturation of the image.

For any outgoing edge at the same source, a polynomial
integral at that source belongs to its normalized edge
lattice. Indeed in the respective adapted row coordinates
the transformed monomials have row weights between the
two end weights. In particular the normalization subtracts
the greatest one. This proves the outgoing tests for
\(i\) and \(U_b\), using integrality of multiplication
in every frame. For any incoming edge at the same target,
the target lattice for \(S^*\) likewise lies in its
normalized intermediate lattice. In the edge frame this uses
\[
 2j-(j-d)=j+d\ge0
\]
for every row weight \(j\in[-d,d]\).
For the tensor of \(t_a\)'s put \(d_T=Nd\).

At the target \(\mathcal C(\,\cdot\,,t_a)\) is integral after
multiplication by \(w^{2d_T}\), and it is integral at the
distinguished normalized edge with no multiplier.
At any other incoming normalized edge it loses at most
\(w^{2d_T}\), by the target inclusion just proved.
The normalized first weights of \(\lambda_+\) exceed those
of \(\lambda\) by \(d_T\). The pullback of a covector
multiplied by \(p^{d_T}\) therefore preserves the target
condition and the raw edge pairing, and decreases the
vanishing conductor at other incoming edges by at most
\(2d_T\). These are the properties needed in the next
subsection. They use the same integral section at all
neighbors of the fixed source or target.
\end{proof}


\subsection{Bounded changes of the full weight}

Apply the contraction to the two flat weights bracketing each prescribed top weight \((\alpha,\beta)\), with endpoints \((\alpha,\beta)\pm c_{\rm cl}\). Their deficits
are \(F_-\) and \(F_+\). The higher order \(t_*\) for the plus
weight can be used in both tests, since \(t_*\) is increasing in
each degree. The minus test is used when each shifted endpoint
is still in its original lower range. Write \(H_-\), \(H_+\)
and \(H_\lambda\) for the corresponding flat, enlarged flat,
and prescribed top weight spaces. Let \(d_-,d_+,d_\lambda\)
be their dimensions divided by \(g\).

The integral comparison needed here is stronger than a dimension
inequality. At each edge it uses a map \(i:H_-\to H_+\),
and maps \(U_b:H_-\to H_\lambda,\ L_a:H_\lambda\to H_+\),
such that
\[
 L_a U_b= \begin{cases}p^{q_b}i,&a=b,\\0,&a\ne b.\end{cases}
\]
The number \(b_0'\) of pairs of maps is the product, over the
tops, of the dimensions of the \(\mathrm{GL}_{k-2}\) modules
of the intermediate weights. The map \(i\) carries the
clean source lattice into the enlarged flat one and identifies
its normalized raw edge lattice with a saturated submodule
of the enlarged raw edge lattice.

Each \(U_b\) maps clean
sources to clean sources at the same order and conductor.
Pullback by each \(L_a\), with the correction
\(p^{A_1(H_+)-A_1(H_\lambda)}\), takes clean covectors on
\(H_+\) to clean covectors on \(H_\lambda\) with conductor
reduced by \(2Nc_{\rm cl}\). The denominator bound for these maps is
\(q_b\le N\lceil\kappa_1(p,a,k-1)\rceil\).
The contraction maps above have these properties.

\begin{proposition}[Edge tests with the restored bounded weights]\label{prop:bounded-weight-test-transfer}
Assume the maps \(i,U_b,L_a\) satisfying the indicated integral
properties, the two flat tests at depth \(D_1\) with deficits
\(F_-,F_+\), and \(4D_1<c\). On every nonboundary weight
where the two tests are available, denote the number of successful directions by \(r_\lambda\). In the prescribed size range,
\[
 r_\lambda\ge(1-F_--2F_+-\eta)d_\lambda,\qquad \eta=12Nkc_{\rm cl}/b_0.
\]

Under the following reserve and conductor bounds the depth is
at most \(D_p\) and the competitor conductor is \(c_\star\).
With the full loss assigned to the missing minus tests on the
lower boundary, the mean deficit over the original box is at most
\(6(C_A+C_s)NN_1/Y+2/Y\).
\end{proposition}

\begin{proof}
\proofstep{Recovering rank through the integral comparisons.}
Given these maps, the integral rank deduction goes as follows.
Use \(n_-,n_+\) successful directions in the respective flat
tests at depth \(D_1\), and make their two pairing matrices
scalar identities \(z_- I,z_+ I\). Transport the \(n_-\)
sources by \(i\), and pair them against the \(n_+\) clean
covectors in \(H_+\), giving a rectangular integral matrix
\(M\). Let \(r\) count its Smith factors dividing \(z_-z_+\),
and let \(s=n_--r\). In Smith coordinates the last \(s\) source rows therefore satisfy
\[
  M_{ij}\in w z_-z_+O\qquad(r<i\le n_-,\ 1\le j\le n_+).
\]
Here and below these are rows after the integral changes of bases.

Saturation gives an integral retraction of the raw edge
lattice in \(H_+\) onto that in \(H_-\), by Smith form.
Pull the \(n_-\) original covectors back along this retraction,
using the pairing scale \(p^{A_1(H_-)-A_1(H_+)}\).
They are integral on the raw edge lattice; they need no
additional vanishing at other edges. Take the column
combinations dual to the row operations of the preceding
paragraph and retain the \(s\) complementary columns.
Their evaluations on the \(s\) rows give \(z_- I_s\).
Now adjoin the \(n_+\) clean rows and covectors. The combined
matrix has blocks
\[
 \begin{pmatrix}z_- I_s & z_-z_+ E\\ B & z_+ I_{n_+}\end{pmatrix},
 \qquad E\in w\operatorname{Mat}(O),\quad B\in\operatorname{Mat}(O).
\]
The Schur complement of the second diagonal block is
\(z_- (I-EB)\).

For right-hand side \((a',b')\), solve explicitly:
\[
 x=(I-EB)^{-1}(z_-^{-1}a'-Eb'),\qquad
 y=z_+^{-1}(b'-Bx).
\]
Since \(E\) vanishes modulo \(w\), \(I-EB\) has an integral
inverse. Thus multiplication of the inverse of the combined
matrix by \(z_-z_+\) makes it integral. This gives \(s+n_+\)
directions at depth at most \(2D_1\), with clean source rows
and raw edge covectors. The outgoing-edge
rank bound used earlier needs extra vanishing only on the
sources, and therefore applies here when \(4D_1<c\).
Edge transitivity allows the same count at every outgoing edge.
It follows that
\[
 r\ge n_-+n_+-d_+.
\]


The low Smith factors of \(M\) supply \(r\) directions
with matrix \(z_-z_+ I_r\) using the clean covectors.
Apply \(U_b\) to these rows and the corrected pullbacks by
\(L_a\) to these covectors. The matrices are block diagonal
with blocks \(p^{q_b}z_-z_+ I_r\).
Scaling each block to the largest valuation, the test on
\(H_\lambda\) has at least
\[
 b_0'((1-F_-)d_- - F_+d_+)
\]
directions, at depth
\(2D_1+N\lceil\kappa_1(p,a,k-1)\rceil\).

\proofstep{Comparing the Weyl dimensions.}
The Weyl dimension formula compares these dimensions without
any loss that depends exponentially on the number of tops.
At a top write \(A=\alpha-c_{\rm cl}\), \(B=\beta-c_{\rm cl}\)
and \(\mu\) for the intermediate weight.
The internal \(\mu\)-factors form the \(\mathrm{GL}_{k-2}\)
dimension.

Each of the remaining \(2k-3\) factors changes by a
ratio between \(1\) and \(1+6c_{\rm cl}/b_0\) from the lower
flat weight to the prescribed weight, or from the lower flat
weight to the enlarged flat weight. This follows directly from
\(|\mu_i|\le c_{\rm cl}\) and \(A,B\ge b_0/3\). Hence
\[
 b_0' d_-\le d_\lambda\le {\cal E}b_0'd_-,
 \qquad d_+\le{\cal E}d_-,\qquad
 {\cal E}=(1+6c_{\rm cl}/b_0)^{(2k-3)N}.
\]
Put \(\eta=12Nk c_{\rm cl}/b_0\), with \(\eta\le1/Y\)
and \(\eta\le1/9\) in the size regime. Since
\({\cal E}\le\exp(\eta)\), one has \[
 1/{\cal E}\ge1-\eta,\qquad {\cal E}\le2.
\] Together with \(0\le F_-\le1\),
these estimates give at least
\[
 (1-F_--2F_+-\eta)d_\lambda
\]
directions.

At positions outside the tops all weights have absolute value
at most \(c_{\rm cl}\). Tensor with a pair of dual bases
in their normalized raw edge lattices. In an edge frame, the
three component lattices have a single coordinate at respective
scales \(1,w^{-1},w^{-2}\), so on the small-weight factors all
normalized orders differ by at most \(2n'c_{\rm cl}\) in
uniformizer units, where \(n'\) is the cycle length. In the
adapted basis let \(j\) be the total weight in that coordinate
and let \(a_S\) be the sum of the first weights of the small
factors. The two normalized exponents have opposite \(j\)-terms,
\[
 e_{\rm source}=j-a_S,\qquad e_{\rm dual}=-j-a_S.
\]

Their differences from the endpoint exponents \(0,-2j\)
are each between \(-2n'c_{\rm cl}\) and zero.
These same inclusions hold in the frames for all other
edges at either endpoint.
Thus the additional factors increase the endpoint order by
\(2n'c_{\rm cl}\), decrease the conductor by the same amount
and multiply the number of successful directions and the
total dimension by the same factor.
The conductor is first set to
\[
 c=c_*+2Nc_{\rm cl}+2n'c_{\rm cl}
\]
to allow these losses.

For the numerical comparison let \(a_r\) be the real lower
endpoint, \(a'=\lceil a_r\rceil\), \(L=\log_p a_r\) and
\(\ell=1+\log_p(8a'+k-1)\). The size hypotheses give
\(L\ge1\), \(\ell\le3L\), \(\log_p a'\le2L\), \(N_1^2\le k\),
\(2\le N\le k\) and \(Y\ge1\). Put \(P=NY^2k^3L\),
\(u_0=\lceil(k+1)\ell\rceil\), \(k_1=\lceil(k-1)^2\ell\rceil\)
and \(\kappa_0=2Nk\log_p a'\). There are \(n_s\le N\) segments.
Adding their depths and the two changes of lattices gives
\[
 D_1=n_s(C_A Y^2 N_1^2k^2\ell+2\kappa_0)+4\kappa_0+4Nu_0.
\]

In units of \(P\), the four terms on the right are bounded by
\(3C_A,8,16,16\), respectively, and \(Nk_1\le4P\).
Indeed \(u_0\le4kL\), \(k_1\le4k^2 L\), and
\(n_sN k\le Nk^2\). Thus the prescribed depth
\(D_p\ge C_H NY^2(N_1^5+k^3)L\), with \(C_H\ge8C_A+93\),
satisfies \(2D_1+Nk_1\le(6C_A+84)P\le D_p\).
As \(c_*\ge8(D_p+1)\), the strict inequalities used when
combining tests hold. The final conductor is exactly \(c_*\).

\proofstep{The depth and order after restoring the weights.}
For the order put \(\mathcal L=C_0N_1^2\log a_r\),
\(\mathcal L'=C_0N_1^2\log a'\), \(Q=p^{G'}\), and set
\[
 f(a,B)=a/p^2+b_0/p+Ba/(pk)+Bb_0/k.
\]
Here \(b_0\) is the lower reference parameter called \(b'\) in
the segment estimate. Let \(S\) be the sum over odd tops \(y\)
of \(\min(b_y,a_{\rm prev}/p^{d_{\rm prev}})\), with the
preceding even top and its distance taken cyclically. Each odd
top occurs once among all segment interiors. Under the plus
shift each summand increases by at most \(c_{\rm cl}\), giving
\(S'\le S+Nc_{\rm cl}\).

The pre-rounding order at these plus
degrees is \(t_0=S'+C_tN(f(a',\mathcal L')+kQ+Qc)\).
Adding \(N\mathcal L' b_0/k\), taking the ceiling and allowing
for the small factors gives
\begin{align*}
 t_*+2n'c_{\rm cl}\le{}&S+C NN_1 f(a_r,\mathcal L)+2n'c_{\rm cl}\\
 &+CN\big[Q(c_*+CNkL+4Nc_{\rm cl}+2n'c_{\rm cl})
                      +kQ+kL+c_{\rm cl}\big].
\end{align*}
Indeed \(a'\le2a_r\) and
\[
 \mathcal L'\le2\mathcal L,
\]
so the
four terms in \(f\) grow by factors at most \(2,1,4,2\).
The extra \(N\mathcal L'b_0/k\) uses the fourth term and the
ceiling costs at most \(1\le NkQ\). Taking \(C\ge4C_t+4\) proves
the bound, including the displayed conductor contributions.

\proofstep{Transport at the boundaries of the intervals.}
To transport a test, first identify the two degree-\(n'\)
unramified coefficient rings over \(\mathbb Z_p\). Their residue
field is \(\mathbb F_q\) with
\[
 q=p^{n'}.
\]
Each embeds in the
unramified ring of degree \(2n'\), with residue extension
\(\mathbb F_{q^2}/\mathbb F_q\). The specified isomorphism
on the smaller residue field extends to the larger ones, by
uniqueness of the quadratic finite field extension. Its Witt
lift extends the specified unramified-ring identification and
commutes with arithmetic Frobenius.

Adjoining the specified
root \(w^2=p\) to each larger ring gives the full valuation
ring by Eisenstein's criterion and total ramification. Extending
by
\[
 x+yw\mapsto\phi(x)+\phi(y)w'
\]
identifies both fraction
fields, the embedded smaller coefficient rings, the specified
uniformizers and the Frobenius on the smaller rings. Every
lattice membership and matrix pairing test therefore transports
in both directions, preserving its depth, conductor and integer
weight indices. Fixing one setting defines the same bound for
all settings at once. When there are no settings the uniform
assertion there is vacuous.

On the lower boundary assign \(F_-=1\) when a minus endpoint
falls outside its interval, and use the deficit bound
\(\kappa_{\rm text}=F_-+2F_++\eta\). For each nonboundary
point, shifting both families of degrees by \(-c_{\rm cl}\)
injects the point into the original box. Hence
\[
 \mathbb E F_-\le\mathbb E_{\rm box}F+
                     Nc_{\rm cl}(H_a^{-1}+H_b^{-1}),
 \qquad \mathbb E F_+=\mathbb E_{\rm plus\ box}F .
\]
The extra term bounds the union of the lower slices in the
\(2N\) intervals and is at most \(1/Y\); also \(\eta\le1/Y\).

Proposition~\ref{prop:edge-box-deficit} bounds both displayed
means of \(F\) by \(2(C_A+C_s)NN_1/Y\), using the same
enlargement and the fixed-\(B\) segment mean. Thus
\[
 \mathbb E\kappa_{\rm text}
 \le 6(C_A+C_s)NN_1/Y+2/Y .
\]
At \(Y=x^{13}\) this yields \(C N_vN_1/x^{13}\), where
\(N=N_v\), in Section~\ref{subsec:frobenius-selection}.
All allowances used for this estimate depend only on the
stated size data, rather than the edge or coefficient setting.

\end{proof}

\subsection{The coefficient rings and sizes on the selected cycles}
\label{sec:host-edge-tests}
\hypertarget{source-host-edge-tests}{}

The preceding transfer and averaging estimates are conditional on local size bounds. We now check those bounds on the good cycles supplied by the arithmetic choice of primes.

Fix the target character and \(0<\epsilon<1/52\). Fix a candidate prime \(p\) of even Frobenius order and a good cycle as in Section~\ref{subsec:frobenius-selection}. For its Frobenius step \(a_p\), the cycle length satisfies
\[
 n'=d_p=\operatorname{ord}(a_p),\qquad 2\mid n'.
\] Write \(N\) for the
number of tops, \(r_G=K+1\) for the fixed plane dimension,
\(G'=G\) for the length bound there, and \(Q=p^{G'}\).
Numbering from the chosen orbit representative identifies its
\(s\)-th place with the translate by \(s a_p\). Thus the top
positions, the signs by parity and the two endpoint weights agree
with those in the local test.

At a top the first and last entries
are \(a_s,-b_s\). The other entries there and all entries at the
remaining places have absolute value at most \(\lfloor c_G\rfloor\).
Here
\[
 c_G=4S_c=16k\Delta_{\rm host},
\]
where \(\Delta_{\rm host}\)
is the size parameter bounding the middle weights. These bounds specify the weights used in the local lattice estimate.

\paragraph{The local field.}
Write \(W=W(\mathbb F_{p^{n'}})\) and \(K_0=\operatorname{Frac}W\).
Take
\[
 O_s=W(\mathbb F_{p^{2n'}})[w]
\]
with \(w^2=p\). Fix a
place \(\mathfrak w\mid p\) of the number field of coefficients
\(F\). Embed \(F\) in an algebraic closure of
\(K_s=\operatorname{Frac}O_s\) by embedding its completion
\(F_{\mathfrak w}\) there over \(\mathbb Q_p\). The field
\[
 K'=K_s F
\]
in that closure is finite over \(K_s\), since a
primitive element of \(F/\mathbb Q\) is algebraic over \(K_s\).

The ring \(O'\) of integers of \(K'\) is the integral closure
of \(O_s\), a finite free module over the complete DVR \(O_s\).
Its valuation, normalized by \(v(p)=1\), restricts to the chosen
valuation on \(F\). This is the ring used for the place in the
host construction. The maps on \(W,K_0\) are the composites of the
standard embeddings with \(O_s\to O'\), \(K_s\to K'\), and the
square root of \(p\) is the image of the specified \(w\).

\begin{lemma}[Finite base change of an edge test]\label{lem:edge-test-host-base-change}
For \(O_s\subset O'\) at the chosen place of the coefficient extension, keep the
same vertices over \(W\) and the same competing edges.
Write \(o,c\) for the order and conductor of a depth-\(D\) test with \(m\) successful directions over \(O_s\). Extension to \(O'\) preserves these four parameters
\[
 (m',D',o',c')=(m,D,o,c).
\]
 The ambient dimension is unchanged.
The valuation normalization is \(v(p)=1\) in both rings.
\end{lemma}

\begin{proof}
\proofstep{The source and target over the actual local rings.}
The vertices are still lattices over \(W\) in \(K_0^k\).
In particular the same edges and all their competitors are
present after enlarging coefficients. At position \(s\) the
twisted lattice is spanned by the images of a \(W\)-basis under
the \(s\)-th Frobenius followed by the embedding. This span,
and the sum of the two spans with multipliers \(1,w\) at an
edge, commute with extension to \(O'\). One may evaluate the
Weyl lattice in any one frame of this sum. A second frame
differs by an integral invertible change of basis and acts by
an integral automorphism, including its determinant twist.

The products of flag minors generating the lattice in this
frame extend coefficient by coefficient. Thus for a lattice \(H\) in one frame its realization after extension of coefficients is
\[
 H_{O'}=\sum_{h\in H}O'\,\iota(h)\simeq O'\otimes_{O_s}H,
\]
where \(\iota\) denotes coefficient extension in that frame.
The same assertion holds for their tensor product on a cycle.
The tableau basis gives the same dimension over both fraction
fields. In the dual space extend each linear functional in that
basis. Integrality on the old lattice then implies integrality
on the new span. Multiplying the endpoint and edge lattices by
their defining powers of \(p,w\) commutes with the maps.

For a target evaluation map with chosen source vectors and
target covectors giving \(p^D I_m\), take their images at the
new place. If \(\iota:O_s\to O'\) is the coefficient map, the evaluation on the chosen test vectors has matrix
\[
 \bigl(\iota(p^D\delta_{ij})\bigr)_{i,j}=p^D I_m.
\]
This follows by extending the bilinear evaluation on each frame basis. The new
lattice conditions at the same vertices are given by those
basis spans, so extending each vector and covector preserves
them. In particular the preimages of the scaled coordinate
vectors still give a right inverse after the same scaling. All required memberships at the
other edges are preserved. The number \(m\) and its denominator
\(\dim V\) are identical. This proves finite base change even
when \(O'/O_s\) is ramified.  \label{cor:host-edge-tests-1}

\end{proof}

\paragraph{Counting tops and bounding the parameters of the auxiliary construction.}
Summing condition (iii) defining a good cycle over its \(n'\)
possible interval starts counts each top \(G'\) times. Also every
gap between consecutive even tops is at most \(G'\). Thus
\[
 n'/G'\le N\le n'N_1/G'\le 40n'\delta=:N_+ .
\]
For the last step use
\(N_1\le2+\max(6eG'\delta,(h+3)\log_2 x)\), with
\(G'\delta\ge (6r_G+1)(h+3)10^{6r_G+1}\log x\).
In particular \(N\gg_\epsilon x^{\epsilon/8}/\log x\),
as \(n'\ge x^{\epsilon/4}\), and \(N_+\le20N_*\) with
\(N_*=2n\delta\). The same marked blocks give at least two
segments and their length and interior counts are bounded by
\(G',N_1\).

We have \(3G'\le n'\) and
\(2L_1(p)\le L_2<\lceil1/\delta\rceil\).
In both ranges use
\[
 \Lambda=C_0 N_1^2\log a,\quad a=x\exp(x^\epsilon),
 \quad b'=a/X',\quad X'=\exp(x^{\epsilon/3}).
\]
Here and below the constants in the local estimates are fixed
before \(x\). From the rank bounds of
Section~\ref{subsec:frobenius-selection}, \(n\ll x^{2+8\epsilon}\),
\(x^{9/2+10\epsilon-\epsilon/100}\le k\le2x^{9/2+10\epsilon}\)
and \(N_1=O_\epsilon(\log x)\). The discriminant bound
\(\log\Delta_{\rm host}\le x^{\epsilon/4}/2\) gives
\(\log(1+c_G)\ll x^{\epsilon/4}+\log x\).
All comparisons below are uniform in the candidate prime, the
cycle and the weights within the indicated intervals.

\paragraph{Large primes.}
Put \(Y=x^{13},\ Y_b=x^\epsilon/20\) and let \(I_b\)
denote the length \(N_* x^\epsilon\Lambda b'/k\) of the
second weight interval. Then
\[
 k^2Y^2\le p<X',\qquad
 Y_bN\Lambda b'/k\le I_b\le b'/2,\qquad
 2Y_bN\Lambda+C N\Lambda^2=o(k)
\]
for each fixed \(C\). Indeed \(N\le40n\delta\) and both
\(x^\epsilon\Lambda\) and \(\Lambda^2\) are
\(x^{2\epsilon+o(1)}\).
We have \[
 \begin{gathered}
 2N_1+2\le k-1, \\ 8\log(NY)\le\Lambda, \\ c_G\le\Lambda b'/k, \\ kY^2Y_bN\Lambda\log a\le X',
\end{gathered}
\] the last left side being
polynomial in \(x\). These verify the spacing and width
conditions even at the ends of the integer intervals.

\paragraph{The hypotheses for the large-prime test.}
To check that the weight ranges remain separated, use the width parameter
\(\Delta=C_{\Delta}N_1\mathcal L/k+2(b_h-b_l)/b'\),
where the local reference \(a\) may be rounded up by one and
\(\mathcal L\) is then computed there. The original real interval
has width \(I_b\). A shift by \(\lfloor c_G\rfloor\), the lower
boundary allowance and the enlarged box add \(O(\Lambda b'/k)\)
to the widths on each side of this interval
(the shift from the harmonic decomposition is bounded by
\(\Lambda b'/k\)). Thus a band of relative half-width
\[
          4I_b/b'+C_\rho N_1\Lambda/k
\]
covers the local band of half-width \(2\Delta\), taking the
fixed \(C_\rho\) large enough for these shifts and rounding
the center upward by at most \(1/b'\).

The possible increase
relative to the center \(a/b'\) is \(O(1/a)\).
This is the band included in the exceptional set of
Section~\ref{subsec:frobenius-selection}. We may take
\(C_\Delta\) large enough that the enlarged weight deviation
is at most \(\Delta/2-18N_1\mathcal L/k\) and
\(40N_1\mathcal L/k\le\Delta\). We have \(\Delta<1/6\)
eventually by the bound on the exceptional width.

\paragraph{Depth and conductor on the selected Frobenius cycle.}
The depth for this cycle is
\[
 D_{\rm I}=\left\lceil C_D N_+Y^2(N_1^5+k^3)\log_p a\right\rceil .
\]
\begin{proposition}[Comparison of the selected lattices for large primes]
\label{prop:host-large-prime-test}
In the stated large-prime range, the rank estimate applies
on the selected lattices at depth \(D_{\rm I}\) and conductor
\(8(D_{\rm I}+1)\). If \(\bar\delta_{\rm I}\) denotes the resulting mean deficit, the depth, conductor and deficit satisfy
\[
 D=D_{\rm I},\qquad c=8(D_{\rm I}+1),\qquad\bar\delta_{\rm I}=O(NN_1/x^{13}).
\]

\end{proposition}

\begin{proof}
The defining inequality makes \(D_{\rm I}\) dominate the local
depth, and the size bounds make it polynomial in \(x\).
Since \(G'\ll_\epsilon x^{\epsilon/8}\log x\) and \(p<X'\),
\(\log Q\ll_\epsilon x^{11\epsilon/24}\log x=o(x^\epsilon)\).
Consequently
\[
 C k^3NQ^4\big(8(D_{\rm I}+1)+(N+n')c_G\big)Y=o(b')
\]
for every fixed \(C\), by taking logarithms. The larger internal
conductor in Section~\ref{sec:edge-transfer}, which includes
\(O(Nk\log_p a)\), satisfies the same inequality.

In particular
\(kQ\le N_1\mathcal L b'/k\). Dividing the expression for \(E_0\)
in Section~\ref{sec:symbol-correction} by \(a\), using
\(p\ge k^2Y^2\), \(b'/a=1/X'\) and the same bound on
\(Q^2c/a\), gives \((E_0+Q)/a=o(N_1\mathcal L/k)\).
The large-prime estimate therefore applies at the prescribed depth
and conductor \(8(D_{\rm I}+1)\). From \(p<a/b'\) we obtain both comparisons needed in the order estimate,
\[
 \frac{b'}p<\frac a{p^2},\qquad \frac{\mathcal L b'}k<\frac{\mathcal L a}{pk}.
\]

The local deficit bound is unchanged by the finite scalar
extension in Lemma~\ref{lem:edge-test-host-base-change}: both
the test rank and its denominator are dimensions of the same
extended vector spaces. Thus it gives \(O(NN_1/x^{13})\)
on the host.
\end{proof}

\paragraph{Smaller primes.}
Proposition~\ref{prop:residue-plane-neighbor-test} gives the
tests in this range. We check its size conditions on the
same cycles and level choices as in the larger-prime case.

\begin{proposition}[Comparison of the selected lattices for smaller primes]
\label{prop:host-small-prime-test}
For \(13\le p<p_0=k^2x^{26}\) on the same selected cycles,
the residue-plane comparison applies at order
\(\lceil t_{\rm II}\rceil\) and conductor \(8(D_{\rm II}+1)\),
with the depth and order in \eqref{eq:small-prime-depth-order}.
Writing \(\kappa\) for its deficit, the resolution and the bound are
\[
 o=\lceil t_{\rm II}\rceil,\qquad c=8(D_{\rm II}+1),\qquad\kappa=O(Nx^{-12}).
\]

\end{proposition}

\begin{proof}
Set \(Y_2=x^{26}\),
\(c_0=2b'+2c_G\), \(A=a-c_G\), and
\(L_1=\lceil\log_p((kx^{26})^4)\rceil\). The interval conditions
give \(A+\lfloor c_G\rfloor\le a_s\le 2A-\lfloor c_G\rfloor\)
and \(b_s+2\lfloor c_G\rfloor\le c_0\), as \(c_G=o(b')\).
For large \(x\) the lower bound on \(N\) gives
\(N_1^2\le r_G N\). Also \(2\le r_G\le Y_2\),
\(N_1+3r_G\le k\), and \(C N_1^2\log a\le k\). The spacing
condition follows from
\[
 p^{L_1-1}\ge k^4 x^{104}/p
 \ge k^2 x^{78}\ge256k^2 N Y_2^2 .
\]
Here the last inequality uses \(N\le n\le x^3\).
The prescribed depth and real order are
\begin{equation}\label{eq:small-prime-depth-order}\begin{aligned}
 D_{\rm II}={}&\left\lceil C_D' N_+(N_1^3+r_G^5+k^2)\log_p a
                 +2r_Gn' c_G+r_G N_+c_0\right\rceil,\\
 t_{\rm II}={}&C_t r_G N\big(\Lambda a/(pk)+a p^{-L_1}
                         +c_0+kQ\big)+2r_G n'c_G .
\end{aligned}\end{equation}

Now \(\log Q\ll_\epsilon x^{\epsilon/8}(\log x)^2
=o(x^{\epsilon/3})\), since \(p\) is polynomially bounded.
Thus \(D_{\rm II}\le x^{O(1)} b'\) and
\[
 C k^3r_G^2NQ^5\big(8(D_{\rm II}+1)+r_G N c_0+
                 2r_G n'c_G\big)Y_2=o(a)
\]
by \(a/b'=X'\). These are all the numerical hypotheses for
the smaller-prime test at the host. Section~\ref{sec:plane-lifts}
proves the estimate on edges with residue quotient of dimension
\(r_G\), at order \(\lceil t_{\rm II}\rceil\) and conductor
\(8(D_{\rm II}+1)\). With the bounds on \(r_G,k\) and
\(Y_2=x^{26}\), the deficit
\(\min(1,CNr_G^2 k/Y_2)\) is \(O(Nx^{-12})\).
Lemma~\ref{lem:edge-test-host-base-change} preserves both
ranks in this fraction. \label{cor:host-edge-tests-2}
\end{proof}

\subsection{Separating and testing residue planes}
\label{sec:plane-lifts}
\hypertarget{source-plane-lifts}{}

In the smaller-prime range of Section~\ref{sec:host-edge-tests}, the incoming lattices are indexed by residue planes. To add their ranks, each test must vanish to higher order on competing planes. Put \(r=r_G\).
An edge of type \(r\) has \(L_+/L_-\) of dimension \(r\) over
\(\mathbb F_q,\ q=p^{n'}\). For an edge of type \(r\), the exponent in the normalized pairing is
\[
                   A_r=\sum_s\sum_{i=1}^r\mu_{s,i}.
\] For a fixed source, competitors are
\(r\)-planes in \(p^{-1}L_-/L_-\). Choose coordinates taking one to the plane \(I_0=\{1,\ldots,r\}\) by an integral frame change.
At this plane the substitution multiplies the even vector
coordinates outside \(I_0\) and the odd covector coordinates
inside \(I_0\) by \(w\). Call these coordinates normal and the
others tangential. At a different plane use the corresponding
twisted change of basis before measuring divisibility.

\begin{proposition}[Neighbor tests on the smaller-prime range]\label{prop:residue-plane-neighbor-test}
Use the selected good cycle and the smaller-prime bounds of
Section~\ref{sec:host-edge-tests}, with its prescribed order
\(t_{\rm II}\), depth \(D_{\rm II}\), and cutoff
\(Y_2=x^{26}\). At each target let \(V_{\rm rec}\subseteq V\) be the subspace recovered from the incoming edges of type \(r\). Then
\[
 \dim V_{\rm rec}\ge(1-\kappa)\dim V,\qquad\kappa=O(Nx^{-12}).
\]

Here both endpoint orders are at most
\(\lceil t_{\rm II}\rceil\), and the evaluation on the recovered
directions has an integral right inverse after multiplication
by \(p^{D_{\rm II}}\).
\end{proposition}

The proof first works in a fixed plane. It separates that plane
from its competitors, lifts a space of residue symbols, and
constructs detecting covectors with the required vanishing.
Packing the resulting per-plane spaces at the same target
then gives the rank assertion.

\begin{lemma}[Frames, separation and normal directions on a good cycle]
The selected cycle admits \(r\)-row frames with all the required
Frobenius powers at most \(Q=p^{G'}\). In the adapted coordinates the separating polynomial vanishes on the other residue planes and is nonzero on the chosen plane, as specified in the proof. The bound for the Frobenius exponents \(p^{a_i}\) required in these frames is
\[
 p^{a_i}\le Q=p^{G'}\quad\text{for every required }i.
\]
 At a selected position the normal
coordinates on the chosen plane have integral polynomial lifts
after multiplication by the frame determinant and by \(w\),
to first order modulo \(w^2\).
\end{lemma}
\begin{proof}
In each marked block of Section~\ref{subsec:frobenius-selection}
choose the middle \(r\) of its \(3r\) isolated selected even
positions. These are consecutive tops. On either side there are
another \(r\) even tops in that block. Between the last position
of one frame and the first of the next there are consequently at
least \(2r\) tops. The first and last \(r\) have preceding gaps
greater than \(L_1\), by their isolation and the empty cells.
The first position of a frame has such a gap as well. From the
beginning of one frame to the end of the next the distance,
including both enclosing blocks, lies within \(G'\), even across
the partial block at wrap-around. Indeed marked block starts
are at distance at most \((M+2)l\), leaving two full block
lengths of slack in
\[
 G'=(M+5)l.
\]
The number of tops in
this segment is thus at most \(N_1\).

For any top after the
incoming frame, the \(r\)-th preceding even top is present in
the segment. Its distance is greater than \(L_1\), since
\(r=K+1\) and no cyclic interval of length \(L_2\ge2L_1\)
can contain \(r\) tops. In particular all powers taking any of
these \(r\) rows up to the tested top are at most \(Q=p^{G'}\).
These properties give the frame conditions and the twisting bounds used below.

\proofstep{Separate Pl\"ucker coordinates along the cycle.}
Write the frame positions as
\[
 a_{j,1}<\cdots< a_{j,r}=z_j
\]
in cyclic order, lifted to integers so that
\(z_{j+J}=z_j+n'\). For each \(j\), form the \(r\)-by-\(k\)
matrix of vector rows at \(a_{j,l}\), raising the \(l\)-th row's
entries to \(p^{z_j-a_{j,l}}\). Write \(P_{j,I}\) for the
minor on a subset \(I\) of size \(r\), \(D_j=P_{j,I_0}\) and
\[
 Q_j=p^{z_{j+1}-z_j}.
\]
There are \(r(k-r)\) subsets \(I\)
differing from \(I_0\) by one exchange. Use all of them in
\[
 R=D_0\prod_I\left(\prod_{j=0}^{J-1}D_j^{Q_j-1}
             -\prod_{j=0}^{J-1}P_{j,I}^{Q_j-1}\right),\quad
 R_0=\prod_j D_j^{\alpha_j},\quad
 \alpha_j=\mathbf1_{j=0}+r(k-r)(Q_j-1).
\]
Modulo \(w\) at a competing plane, let \(B\) be its rank-\(r\)
matrix of rows over \(\mathbb F_q\) in the distinguished frame
and \(b_I=\det B_I\). The \(j\)-th row matrix is a product of
an \(r\)-square matrix in the free variables with \(B^{[p^{z_j}]}\).

Indeed at a position \(a\) the surviving row is a linear
combination of the plane rows twisted by Frobenius at \(a\).
Powering moves them all to \(z_j\). Write Frobenius exponents
modulo \(n'\) in this residue calculation. Every minor thus has
the same polynomial factor at \(j\), times \(b_I^{p^{z_j}}\).
For \(b_I\ne0\) the factors telescope:
\[
 \prod_j (b_I^{p^{z_j}})^{Q_j-1}
    =b_I^{p^{z_J}-p^{z_0}}=1,\qquad z_J=z_0+n' .
\]
The final equality uses \(b_I\in\mathbb F_{p^{n'}}^\times\). If \(b_{I_0}=0\),
the initial \(D_0\) vanishes. Otherwise row reduce \(B\) in
the first \(r\) columns. In the resulting matrix \([I_r\mid C]\), a plane different
from the coordinate plane has some \(C_{ih}\ne0\). Exchanging column
\(i\) of \(I_r\) with column \(h\) of \(C\) gives
\[
 \det(e_1,\ldots,e_{i-1},C_{\bullet h},e_{i+1},\ldots,e_r)=C_{ih}\ne0 .
\]
Reordering the columns can only change the sign. The corresponding
difference factor therefore vanishes. This proves \(R\in wO[\mathbf X]\) after
substitution at every competitor.

At the chosen plane \(D_j\) depends only on tangential
coordinates. In every monomial of a one-exchange minor one
entry comes from a normal column. After substitution its
coefficient has \(w\)-order equal to its total normal degree.
Expanding \(R-R_0\), or \(R^c-R_0^c\), each time a one-exchange
term is chosen it costs at least \(Q_j-1\) normal degrees at
each \(j\). Thus, before collecting terms, each error term of normal degree
\(e\) has an explicit factor \(w^e\) and retains the tangential
factor
\[
 D_j^{\max(0,c\alpha_j-e)}
\]
at each frame for \(R^c\).
Here the word monomial refers to expansion after the undamaged
determinant factors have been set aside.

\proofstep{Recover a normal coordinate at the chosen plane.}
Use the \(r\) preceding selected even positions at each selected
\(s\), and power their vector rows to the index \(s\). Denote the
first \(r\) columns by \(U_s\) and put
\[
 T_s=\det U_s.
\]

Every power comes from a strictly earlier position. For even
\(s\) and \(h>r\), append column \(h\) and the current row to
form an \((r+1)\)-minor. At the distinguished plane its leading
expression is \(wT_s X_{s,h}\), with error divisible by \(w^2\).
Any other term using column \(h\) from an earlier row has at
least \(p\) powers of \(w\). At every plane the enlarged
matrix has rank at most \(r\) modulo \(w\), by the same
Frobenius factorization.

For odd \(s\), contract each powered
row with the current covector. At every plane these contractions
vanish modulo \(w\), since the surviving covector annihilates the
surviving row space. The leading terms at the distinguished plane
are the entries of
\[
 wU_s(X_{s,i})_{i\le r},
\]
again with errors
in \(w^2\). The adjugate of \(U_s\) expresses \(T_s X_{s,i}\)
as a tangential linear combination of their leading terms
divided by \(w\). This gives an order-one polynomial for each
required term. All letters are at selected positions.

\end{proof}

\paragraph{Correction to the prescribed precision.}

Separating one residue plane gives a congruence at the other planes. A correction is needed to reach the prescribed precision there while preserving the initial symbol.

\begin{lemma}[Correction on all residue planes]\label{lem:plane-normal-correction}
In this paragraph a factor
of order \(d\le2Q\) means an integral polynomial having order
at least \(d\) in every plane frame, with distinguished leading
expression \(S\) divisible by \(w^d\), homogeneous in the
normal coordinates of degree \(d\), and accurate modulo \(w^{d+1}\).
Use finitely many such factors \(K_l\) with expressions \(S_l\),
a power \(w^i\), \(i\ge0\), and any integral tangential polynomial
\(H\). Use only letters at tops. The level is \(L=i+\sum_l d_l\).

At any \(b\) with \(0\le b\le c\), consider the desired expression
\[
  w^i\prod_l S_l\, R_0^c\prod_jD_j^{b}
                  \prod_{s\ {\rm top}}T_s^{(2Q+3)b}\,H.
\]
In this expression all determinants are in the distinguished
coordinates. There is an integral polynomial whose substitution
agrees with it modulo \(w^{L+b}\) and whose order at competitors
is at least \(L+b\). For an initial level \(L\ge o\), taking
\(b=c\) gives the precision \(o+c\) needed before division by
\(w^o\). The construction can preserve any prescribed full
multidegree of the desired expression.
\end{lemma}

\begin{proof}
To prove the lemma, form the raw product by using the true
factors \(K_l,R\) and the true determinants and tangential
polynomial in the distinguished frame. The latter determinants
and polynomial restrict exactly there, since their entries are
tangential. Every term of the raw product has the required
competitor divisibility as a polynomial, by the universal
orders and the factor \(R^c\). The case \(b=0\) also follows
by taking zero. For \(b>0\), expand its error at the designated
plane. A replaced factor \(K_l\) contributes a monomial of
normal degree \(m\) with coefficient divisible by
\(w^{d_l+e_l}\), where
\[
 e_l=\max(1,m-d_l).
\]

Indeed the substitution of an integral polynomial first incurs
\(w^m\) on that monomial, and its difference from \(S_l\)
is divisible by \(w^{d_l+1}\). If \(R^c\) is replaced, let
\(e_R\) be its order as computed above, otherwise set \(e_R=0\).
The total extra order is
\[
 e=e_R+\sum_{\rm replaced}e_l>0.
\]

Discard terms with \(e\ge b\). For the others we recover all
\(R_0^c\) by spending at most \(e\) of the \(b\) spare
determinants at each \(j\). Put the unused excess above
\(b-e\) into the tangential coefficient.

\proofstep{The remaining normal order after correction.}
The newly exposed normal multiplicity \(m_*\) satisfies
\[
 m_*\le e_R+\sum_{\rm replaced}(d_l+e_l)
       \le (2Q+1)e .
\]
At each normal variable use one copy of the corresponding
\(T_s\) and the preceding coordinate identity. This replaces it
with a sum of order-one leading expressions divided by \(w\)
and multiplied by tangential coefficients. There are enough
copies at every position to leave \((2Q+3)(b-e)\) copies, with
the surplus also tangential. After absorbing \(w^{m_*}\) into the order-one leading
expressions, the external power is
\(i'=i+\sum_{\rm replaced}d_l+e-m_*\ge0\).

Adding the degrees of the untouched factors and the
\(m_*\) order-one factors yields precisely
\[
 i'+\sum_{\rm untouched} d_l+m_*=L+e .
\]
The residual determinant reserves now have parameter \(b-e\),
so each error term is in the induction at the smaller parameter. Induction lifts them
to precision \((L+e)+(b-e)=L+b\) at the chosen plane and with
at least that order at competitors. Subtract those lifts from
the raw product. Finally one can project onto the prescribed
multidegree, since every substitution is linear within each
true letter and preserves its total degree.  Its use requires generating symbols
with the displayed determinant powers in the available degrees.

\end{proof}

\paragraph{Extracting determinants from a prefix.}

The correction must also preserve the orders at the cuts of a path. Extracting powers of the determinant from a prefix measures how much of its order remains.

\begin{lemma}[Determinant divisibility from a prefix order]\label{lem:plane-prefix-determinant}
In the plane variables
before a cut, let \(A\) be a list of \(r\) distinct vector
letters, let \(Z\) be the square matrix of their rows, and
let \(B\) be the greatest total degree in all other vector
letters before the cut in any monomial of an integral polynomial
\(P\). Coefficients may themselves be polynomials in auxiliary
variables over the DVR. Suppose that under a generic conjugate
of \(\operatorname{diag}(1,\ldots,1,T)\) acting on the prefix
vectors and contragrediently on its covectors every Laurent
power in \(P\) is at least \(\theta>B\). Then
\[
                    P=(\det Z)^{\theta-B}P_1
\]
integrally. Dividing by this factor lowers the minimum Laurent
power on any other subset of letters by at most \(\theta-B\).
\end{lemma}

\begin{proof}
\proofstep{The determinant identity on the prefix chart.}
For the first assertion localize the polynomial ring at the
prime of the generic determinant \(\pi=\det Z\). This ring
is a DVR. Indeed the generic determinant is irreducible over
the fraction field of the coefficients and primitive over the
coefficient UFD. Irreducibility follows by expanding in the
last row, whose complementary minors have no common factor.
For \(r>1\) two distinct such minors are nonassociate irreducibles
by induction, and the case \(r=1\) is immediate. Only the finitely
many variables occurring in the calculation are needed.

In the
residue field of this localization \(Z\) has rank \(r-1\)
since a nonzero minor of size \(r-1\) is a unit there.
Choose \(g\in\mathrm{GL}_r\) over the localization whose last
column spans its kernel modulo \(\pi\), by lifting a basis of
the residue space. Write
\[
 H_T=g\operatorname{diag} (1,\ldots,1,T)g^{-1}.
\]
Specialize the generic conjugating
matrix to this \(g\). All coefficients remain integral and the
lower Laurent-power bound persists, as the only matrix denominator
is a power of the unit \(\det g\).

In the transformed polynomial evaluate the input prefix rows at
\(xH_\pi^{-1}\) and the input covector columns at \(H_\pi y\).
At \(T=\pi\) the result is exactly the original \(P\).
All covectors are integral. Each row of \(Zg\) has last entry
divisible by \(\pi\), so the \(r\) chosen input rows are
integral as well. Each remaining prefix vector can cost one
power of \(\pi\) per coordinate factor. Their combined
degree is still at most \(B\) under a linear change within
each letter. For a term \(A_l\) of Laurent degree \(l\ge\theta\), the possible denominator loss is therefore at most \(B\), giving
\[
 \operatorname{ord}_\pi(A_l|_{T=\pi})\ge l-B\ge\theta-B.
\]

This gives the claimed divisibility in the localization and
hence in the original UFD. For the last assertion apply the
generic conjugate to any subset of the determinant rows. Their
wedge is transformed by an exterior power of \(H_T\), in which
the largest Laurent exponent is at most one. The determinant
has minimum exponent at most one there. Minimum exponents
add on nonzero products over a domain, giving the asserted
bound for \(P_1\). The case \(P=0\) uses \(P_1=0\).
The other cuts are controlled simultaneously.

\end{proof}

\begin{lemma}[Allocating the determinant reserves]\label{lem:plane-determinant-reserve}
Assume the signed-prefix order and the one-sided coarse degree
bound of \eqref{eq:plane-prefix-degree-assumptions}, with
\(p\ge13\), \(k\ge3\), \(Y\ge2\), and the bounds on size and precision in Section~\ref{sec:host-edge-tests}.
The residue-tangential product has a factor
\(F_*\prod_{s\ {\rm top}}T_s^{kQ}\) to precision \(p^{c'}\)
after substituting the original radices. Here
\[
 F_*=R_0^{c'}\prod_jD_j^{c'}
                         \prod_{s\ {\rm top}}T_s^{(2Q+3)c'} .
\]
These factors are available in the stated degrees for
Lemma~\ref{lem:plane-normal-correction}.
\end{lemma}

\begin{proof}
\proofstep{Allocating the determinant factors.}
For an application at the coarse degrees, let \(z\) be the last
position of the outgoing frame in a segment. At a top \(s\)
after the incoming frame take \(A\) to be the preceding \(r\)
even tops and \(u\) their earliest position. Write
\(B_o=\deg(u)/p^{z-u}\), with the true top degrees between
\(a_0\) and \(2a_0\), where \(a_0\) is the lower bound for
the plane model (at the host it is \(a-c_G\), up to integer
rounding). Put \(Y=Y_2\). Let the minimum of the Laurent exponents at the
prefix strictly before \(s\) be \(\operatorname{ord}_{<s}\).
The conditions for this prefix and the coarse degrees on the
incoming frame and intervening tops are
\begin{equation}\label{eq:plane-prefix-degree-assumptions}
 \operatorname{ord}_{<s}\ge\theta
      =\left\lceil B_o/2-\left\lfloor B_o/(Yk^2)\right\rfloor\right\rceil,
 \qquad
 d_y\le1.01\deg(y)/p^{z-y}.
\end{equation}
The order bound holds on sums because it excludes every lower
Laurent power.

In particular the stated one-percent two-sided hypothesis suffices.
Every earlier even top outside \(A\) belongs to these lists,
because the outgoing frame ends the segment and has only
\(r\) members. Such tops have \(y\le u-2\), so
\[
 B\le \sum_{y<u,\ y\ {\rm even}} d_y
       \le 2.02 B_o\sum_{j\ge1}p^{-2j}
       =2.02 B_o/(p^2-1).
\]
For \(p\ge13,\ k\ge3,\ Y\ge2\), this leaves
\(\theta-B> 2B_o/5\). In particular these are large
reserves even for the earliest row, since \(B_o\ge a_0/Q\).
For example let \(c'=8(D_{\rm II}+1)+rN c_0+
2rn'\lfloor c_G\rfloor\) and set
\[
 m_*=(N_1+1)\big((2Q+3)c'+kQ+c'\max_j\alpha_j+3c'\big).
\]
As \(\max\alpha_j\le1+rkQ\), \(N_1^2\le rN\), and
\(k^3r^2NQ^5 c'Y=o(a_0)\) by
Section~\ref{sec:host-edge-tests}, we have \(N_1m_*=o(a_0/Q)\).

At each such cut extract \(m_1=(2Q+3)c'+kQ+c'\) copies of the
ordinary determinant on the chosen rows. From the incoming frame
extract \(m_2=c'\max_j\alpha_j+2c'\) copies of its determinant.
This frame consists of the preceding \(r\) rows at the first cut.
The total number extracted in a segment is at most
\[
 N_1m_1+m_2\le m_*.
\]
Each copy costs at most one in every
prefix minimum and cannot increase the degree in any subset of
the other vector variables, since highest degrees in a subset
add on nonzero products. The absolute lower bound
\(\theta-B>2B_o/5\) permits these successive extractions.

The
original plane-versus-complement inequalities have margin
\(m_*\), which absorbs their loss from the extractions at each cut,
including the cuts that end at a covector.
The determinants continue to give the required tangential
factors after radix substitution, to the precision needed for
the symbol. Indeed put \(P_s=p^{z-s}\). The image of a
determinant on the \(r\) preceding rows is \(\det(U_s^{[P_s]})\).
Over \(O\) this is
\[
 T_s^{P_s}+pE
\]
with \(E\) integral,
since Frobenius commutes with determinants in characteristic
\(p\). In a power of \(h+c'\) such factors, terms using more
than \(c'\) copies of \(pE\) vanish modulo \(w^{c'}\).
The remaining terms have the common factor \(T_s^h\), with
integral tangential quotient.

Take \(h=(2Q+3)c'+kQ\).
The coarse determinant of the incoming frame has image
\(D_{j-1}^{p^{z-z_{j-1}}}+pE'\). The same argument gives a
factor \(D_{j-1}^{(1+\alpha_{j-1})c'}\) from the \(m_2\) copies.
Every top occurs after the incoming frame in exactly one segment
per cycle. Around the cycle the factored expression is therefore
\[
 F_*\prod_{s\ {\rm top}}T_s^{kQ}\,H,\qquad
 F_* = R_0^{c'}\prod_jD_j^{c'}
                 \prod_{s\ {\rm top}} T_s^{(2Q+3)c'},
\]
for an integral tangential polynomial \(H\), and the discrepancy
has coefficients in \(p^{c'}O\).  The factor \(\prod_s T_s^{kQ}\) will be used
on the digit labels.

\end{proof}

\paragraph{Integral generators after determinant extraction.}
\hypertarget{source-plane-symbols}{}
This part of the proof of Proposition~\ref{prop:residue-plane-neighbor-test}
gives the integral source polynomial on a single plane.
The determinant factors extracted above remain part of the
symbol. The result of the lifting is as follows.

\begin{lemma}[Lifting a plane symbol]\label{lem:lift-plane-symbol}
Fix a selected plane and the determinant factors of
Lemma~\ref{lem:plane-determinant-reserve}. Suppose that a residue symbol
\(\sigma\) has normal degree \(o\) and belongs to the prescribed
span of digit monomials times products of segment symbols,
with their prefix and multidegree conditions. Write \(\Pi\) for the selected plane. At precision \(c'\) there is one integral polynomial \(P\) with the following two restrictions
\[
 P|_\Pi\equiv w^o\sigma\pmod{w^{o+c'}},\qquad P|_{\Pi'}\equiv0\pmod{w^{o+c'}}\quad(\Pi'\text{ competing}).
\]
 It may be projected to the required
true multidegree. The column factors used in the lift have
prescribed order at most \(2Q\).
\end{lemma}

\begin{proof}
\proofstep{A paired path after the determinant extraction.}
Split the coordinate module into its plane and complement,
\[
 O^k=E\oplus C,\qquad E=O^r,\quad C=O^K,\quad K=k-r.
\]
At a top indexed by \(s\)
the even letters are vectors and the odd letters are covectors.
Fix one segment with end index \(z\). Its prescribed coarse
degree in a letter is the sum of its \(E\)-degree and its
\(C\)-degree. A segment symbol belongs to the span of the products \(A(E)G(C)\) satisfying the prefix bounds. To put the summands at the prescribed coarse
degrees, decompose \(A\) and \(G\) into homogeneous components
in every letter and keep the pairs whose degrees sum to the
prescribed ones. If \(\pi_\alpha\) extracts a multihomogeneous component, each prefix substitution \(g\) is linear within its letters. Hence
\[
 \pi_\alpha(gF)=g\,\pi_\alpha(F).
\]
Each component retains the same bounds. The
normal-degree cutoff is a condition on the original segment
symbol used in the rank argument. The realization here
applies to this whole span after the multidegree projection.

Apply the determinant extraction to \(A\) in each summand.
The determinants have fixed degrees in each coarse letter, so
the quotient \(A'\) also has fixed degrees. For every top \(s\)
after the incoming frame, let \(a_{<s},a_{\le s}\) denote the
minimum Laurent exponents of \(A'\) under the two prefix actions
ending before and at \(s\). Let \(g_{<s},g_{\le s}\) denote the
corresponding maximum exponents of \(G\). The remaining
separation says
\[
  g_{\le s}\le a_{<s}\quad(s\ {\rm even}),\qquad
  g_{<s}\le a_{\le s}\quad(s\ {\rm odd}).
\]
The complemented factor has degree zero in the incoming frame.

\proofstep{Integral basis after removing the determinant factors.}
Use the integral Pieri basis from
Section~\ref{sec:column-lifts} for the two factors separately.
The prefix filtration used there also works with a lower
bound on the last entry of the intermediate weight. In each
column the signed subset weight on one coordinate lies
between the contributions from the smallest and largest
entries of the intermediate weight. In the product it lies
between those two entries. This gives both bounds upon
conjugating the one-coordinate torus.

Conversely, over the
fraction field, the requirement that every conjugate torus
act within the specified range selects exactly the isotypic
summands whose first and last entries lie within that range.
Indeed the permitted subspace is invariant under the prefix
group. On each copy of \(V_\lambda\) in \(V_\lambda\otimes M_\lambda\), the extrema of the coordinate weights \(\mu\) are
\[
 \max\mu=\lambda_1,\qquad \min\mu=\begin{cases}\lambda_r,&\text{plane},\\\lambda_K,&\text{complement}.\end{cases}
\]
Complete reducibility gives the assertion. Splitting
successively by Pieri at the prefixes as in that section proves
the dimension formula for the intersection. The allowed basis
words span the intersection over \(O\) because the full basis is
unimodular.

Use the actual lower prefix exponents of \(A'\) as its lower
bounds and the actual upper prefix exponents of \(G\) as its
upper bounds, discarding zero summands. If \(\nu\) and
\(\kappa\) are the paths in any resulting product of basis
words, then for even \(s\) one has
\[
 \kappa_1(\le s)\le\nu_r(<s).
\]
For odd \(s\) one has
\(\kappa_1(<s)\le\nu_r(\le s)\).
These inequalities give exactly the interlacing across the
boundary between the \(r\)-th and \((r+1)\)-st rows of the
concatenated path \(\lambda=(\nu,\kappa)\). Each separate path already interlaces within its own rows.

For the boundary, if \(\lambda_i(s)\) denotes the combined
row weight after \(s\), the inequalities say
\[
 \lambda_r(s-1)\ge\lambda_{r+1}(s)\quad(s\ \mathrm{even}),
 \qquad
 \lambda_r(s)\ge\lambda_{r+1}(s-1)\quad(s\ \mathrm{odd}).
\]
The two row weights are nondecreasing at vector steps and
nonincreasing at covector steps. Thus the displayed boundary
inequality and the separate interlacings give dominance
both before and after each step. In the incoming frame
the complement stays at zero and plane entries are nonnegative
and increasing, so interlacing also holds there. Each strip's
size is the sum of the two sizes, and the concatenated path
starts at zero.

\paragraph{Full order of the two-block columns.}
\label{par:plane-column-orders}
Fix a column \(b\) of a concatenated path, and write \(h\) for
its height at a cut. A positive column uses the exterior
degree \(h\); a nonpositive column uses dual exterior degree
\(k-h\), as in Section~\ref{sec:column-lifts}. All changes of
height are by one, rising at a vector and falling at a
covector. For a subset \(I\) of size \(h\) in a positive column,
\[
 \#(I\cap C)=h-\#(I\cap E)\ge h-\min(r,h).
\]
Here the intersections count coordinate indices in the two summands. Equality characterizes the subsets of smallest normal cost. For a
nonpositive column
\[
 h>r
\]
at every cut. Indeed it starts at
\(k\) and loses at most one per letter, and the segment length
plus \(r\) is at most \(K\). The subsets of smallest cost
in the dual tensor use only complement indices.

The costly transitions have precisely the two forms
\[
 h\longrightarrow h+1\quad(h\ge r,\ \text{vector}),\qquad h\longrightarrow h-1\quad(h\le r,\ \text{covector}).
\]
In a nonpositive
column only the first possibility occurs. This is the normal
coordinate selected when both exterior subsets have their
prescribed plane occupancy. For a positive column, define the
defect of a subset as \(\min(r,h)\) minus its plane occupancy.
For a nonpositive column its defect is its plane occupancy.
In either case, along a sequence of wedge and contraction
terms, write \(n_i=1\) if the \(i\)-th letter
uses a normal coordinate and \(c_i=1\) if that step is
costly, taking both indicators to be zero in the other
case. The defect \(\delta_i\) then satisfies
\[
                    \delta_i-\delta_{i-1}=n_i-c_i .
\]
This follows from the change in occupancy in the wedge
or contraction operation and the change in the target
occupancy \(\min(r,h)\).

For example, at a positive rise from \(h<r\), adding a
complement index increases the defect by one. At a positive
fall from \(h\le r\), removing a complement index decreases
it by one. When \(h>r\) a rise on a plane index decreases the
defect and a fall on that index increases it. The dual
operations wedge on a covector and contract on a vector,
which increase and decrease its plane occupancy when the
plane index is used.

\proofstep{The allowed path of the restored column.}
The segment letters are the increasing enumeration of a set of
distinct top positions \(s_i\). Their radices
\[
 \omega_i=p^{z-s_i}>0
\]
are strictly decreasing. Let
\(\delta_i\) be the defect just after letter
\(i\), with \(\delta_0=0\). The excess of weighted normal cost
over the weighted count of costly steps is
\[
 \sum_{i<M}\delta_i(\omega_i-\omega_{i+1})+\delta_M\omega_M
\]
for a segment of \(M\) letters. Thus the lowest part is
computed by keeping defect zero throughout. In positive
columns the plane operations then take place with no
complement indices in the subset. All remaining operations
act on the complement past the plane indices.

In the
nonpositive columns the plane exterior degree in the dual
tensor stays zero. With the right wedge and contraction
conventions, an operation on a complement index counts only
larger complement indices for its sign, and every active
operation on a plane index counts only larger plane indices.
The defect-zero recursion therefore splits without a sign
change into the two column recursions for \(\nu\) and
\(\kappa\). For terminal coordinates take their union in the
corresponding blocks. Writing \(\operatorname{in}\) for the component of smallest normal degree, we obtain for the column polynomials \(C_b\)
\[
 \operatorname{in}\Bigl(\prod_b C_b\Bigr)=\prod_b\operatorname{in}(C_b).
\]
The defect-zero recursion identifies the right side with the prescribed paired word.

For the true lift, use the same exterior recursion in the
designated twisted frame. Between successive used indices
\(s<t\), raise every entry to \(p^{t-s}\), then perform the
operation at \(t\). After the last operation power up to \(z\).
Take \(\omega_i=p^{z-s_i}\), where \(s_i\) is the index of
letter \(i\). The column has degree \(\omega_i\) in each used
true letter. Its prescribed order is
\[
             d_b=\sum_{i\ {\rm costly}}\omega_i\le 2Q .
\]
The bound follows from the strict spacing of the used integer
indices and the geometric sum starting at the largest power,
which is at most \(Q\).

\proofstep{The order in a neighboring plane.}
For any tested plane, measure the running exterior tensor in
its own twisted frame. Suppose at a stage the common order is
at least \(d\), and every entry with positive defect has order
at least \(d+1\). This holds before the first operation.
For the passage \(s<t\) put \(P=p^{t-s}\) and write the
adapted coordinates at \(s\) as \(w^d U\). If \(L_s,L_t\)
are the changes of exterior coordinates from the designated
frame to the tested frame at the two times, then
\[
       (w^d U L_s^{-1})^{[P]} L_t
                 = w^{Pd}\big(U^{[P]}+pV\big)
\]
with \(V\) integral. Entrywise Frobenius is additive modulo
\(p\), and \((L_s^{-1})^{[P]}\equiv L_t^{-1}\pmod p\)
because the matrices and their inverse minors come from the
unramified ring and its successive Frobenius twists.

Since
\(p=w^2\) and \(P\ge2\), the terms in positive-defect entries
now have order at least \(Pd+2\). At a free step the free
coordinates take defect-zero subsets to defect-zero subsets,
and the normal coordinates contribute a factor \(w\).
At a costly step the free part kills the entries at defect
zero. The normal part carries them to the new defect-zero
subsets and contributes \(w\). Use as a common lower bound immediately before the linear step the order \(Pd\).
After a costly step the required order is \(d'=Pd+1\).
Thus the guaranteed lower bounds for the two contributions are
\[
 \begin{aligned}
 \text{defect zero, normal factor }w &: \quad Pd+1=d',\\
 \text{positive defect, integral factor} &: \quad Pd+2\ge d'+1 .
 \end{aligned}
\]
Only the first contribution can affect the leading coefficient.
At a free step the required order stays \(Pd\), and a normal
factor again raises it. These inequalities apply in every frame.

At the designated frame the change-of-frame matrices are
identities. After division by the prescribed power of \(w\)
the leading recursion simply keeps the defect-zero part. The
powering between steps and up to \(z\) commutes with its
additions and multiplications modulo \(w\), where the
characteristic is \(p\). Its result is therefore the
defect-zero coarse column with the substitution
\[
 X_{i,a}\mapsto X_{s_i,a}^{p^{z-s_i}}.
\]

Consequently the integral true column is a factor of order
\(d_b\) for the correction lemma. The product of its
normalized leading expressions is the desired paired symbol
after this substitution, with the full normal degree equal
to the sum of the column orders.  \label{cor:plane-symbols-1}

\paragraph{Digits, normal degree, and the final projection.}
Let \(X^\ell\) be the common digit monomial. At every top
each digit is less than the letter's radix, which is at
most \(Q\). Thus its total normal degree at that top is at
most \(kQ\). Spend one \(T_s\) for each normal digit factor
and use the leading order-one formulas for that coordinate.
If \(d_s\) normal digit factors are used at this frame, the reserve of \(T_s\) becomes
\[
 T_s^{kQ-d_s},\qquad kQ-d_s\ge0.
\]
It and the tangential digit variables are integral tangential coefficients. This
expresses \(X^\ell\prod_s T_s^{kQ}\) as a sum of products of
normalized leading expressions of universal order-one
factors with integral tangential coefficients. Every term has
the full normal degree of the digit monomial.

\proofstep{Recovering the true letter degree.}
Now start with a symbol \(\sigma\) of normal degree \(o\),
presented as \(X^\ell\) times a sum of products of the
transformed segment symbols. The extracted determinant
product is common to these summands. In the distinguished
coordinates it gives the identity
\[
  \sigma =
  F_*\,H\,X^\ell\!\prod_s T_s^{kQ}\,S + p^{c'}Y',
\]
where \(S\) is the sum of products of transformed residual
segments and \(Y'\) is integral. After expanding residual
segments in paired words, every summand on the right outside
the error has normal degree equal to the sum of the
appropriate column and digit orders. Since \(F_*,H\) are
tangential, take the component of normal degree \(o\).
The resulting main terms, after multiplication by \(w^o\),
are products of leading expressions of combined order \(o\),
with the reserves for precision \(c'\) contained in \(F_*\).
The error acquires order at least \(o+2c'\).

Apply the correction lemma to each surviving term, then add.
For the resulting integral true polynomial \(P\), the two tests read
\[
 \begin{aligned}P|_{\rm designated}&\equiv w^o\sigma\pmod{w^{o+c'}},\\ P|_{\rm competing}&\equiv0\pmod{w^{o+c'}}\qquad(\text{each competitor}).\end{aligned}
\]
If \(\sigma\) has a prescribed true multidegree,
project the resulting polynomial on that multidegree.
Every plane substitution is linear within each true
letter, so the projection commutes with all the substitutions
and preserves all divisibilities.  Its uniform order
bound on a key factor is \(2Q\), and its precision is the
prescribed \(c'\) at every plane. The determinant estimate above
also proves the one-sided version used for the dual test in
Section~\ref{par:plane-reassembly}. \label{cor:plane-symbols-2}
\end{proof}

\begin{paramdefinition}[Frames, coarsening and sizes for the plane lattice]\label{hyp:coarsened-plane-frames}\label{par:plane-rank}
\hypertarget{source-plane-rank}{}
In the plane construction let \(r\) be the plane dimension and \(K\) its codimension in dimension \(k=r+K\). Write \(a\) for the lower bound on the boundary degrees \(b_x\).
Let \(J\) be the number of frames, \(N\) the number of tops in a
cycle, and \(N_s\) an upper bound for the number of letters in a
segment. In the count use the lifted frame endpoints \(s_j,z_j\)
with \(s_j\) the first position of frame \(j\). Segment \(j\)
starts at \(s_{j-1}\) and ends at \(z_j\). Write \(Q\) for the
original bound on \(p^{z_j-x}\). The coarsening uses
\[
 A=\frac{a}{8Yk^2N_s},\qquad L=\lfloor\log_p A\rfloor,\qquad
 P_{j,x}=p^{L+s_{j-1}-x},\qquad H_j=p^{s_j-s_{j-1}}.
\]
Indeed the extra power at segment \(j\) is the largest power of
\(p\) at most \(\lfloor A/p^{z_j-s_{j-1}}\rfloor\).

Since \(A\ge Q\),
this equals \(p^{L-(z_j-s_{j-1})}\). In particular \(P_{j,x}\le A\),
\(P_{j+1,x}=P_{j,x}H_j\) on frame \(j\), and \(\prod_j H_j=q\).
Use the separation bound \(p^{L_1-1}\ge128 k^2N_sY^2\) and the
size bound
\begin{equation}\label{eq:plane-count-size}
  a\ge32pYk^2N_sQ^2(E_0+m_*+1).
\end{equation}
These are the separation and size inequalities for the plane model.
The argument below uses them with \(p\ge13\), \(r\ge2\),
\(Y\ge1\), \(N_s+r\le K\), \(a\le b_x\le 2a\), and the lifted
frame assumptions of Section~\ref{sec:plane-lifts}.
Put \(D_k(b)=\binom{b+k-1}{k-1}\).
\end{paramdefinition}
For the selected cycle, the separation and size comparisons in the proof of Proposition~\ref{prop:host-small-prime-test} use \(Y=x^{26}\), \(p^{L_1-1}\ge256k^2NY^2\) and \(\log Q=o(\log(a/b'))\). They hold uniformly over the admissible frames.


\begin{lemma}[Frame and digit degrees in the coarsened lattice]\label{lem:plane-digit-frame-lattice}
For the coarsened lattice of Definition~\ref{hyp:coarsened-plane-frames},
the late degrees satisfy
\[
 a'_x=b_x-E'_x-E_x>0.
\]

The digit and outgoing choices produce the same total degree
\(a'_x\) in the true coordinate at each top, with incoming
degree \(b_{j+1,x}\ge W_{j+1,x}\ge kN_s+1\).
The corresponding segment word families are independent over
the residue field in the fixed digit and degree choices.
\end{lemma}

\begin{proof}
\proofstep{The coarse degrees at the later tops.}
At each late top \(x\) of segment \(j\), meaning one after the
incoming frame, form the ordinary plane determinant of the last
\(r\) preceding even rows and give it the exponent
\[
 g_{j,x}=\left\lfloor\frac{b_u}{k^2Yp^{z_j-u}}\right\rfloor,
 \qquad u=\text{the earliest of those rows}.
\]
Write \(\mathcal D\) for the product over all these determinants
with their original radices \(p^{z_j-y}\).

At a top \(y\) its
true degree \(E_y\) is at most \(4N_s a/(k^2Y)\). A determinant
contributes at most \(g_{j,x}p^{z_j-y}\le 2a/(k^2Y)\), since
\(y\ge u\). Each top occurs in at most two segments. Use coarsened
true degrees \(a'_x=b_x-E'_x-E_x\), where \(E'_x\le E_0\) is the
prior cleaning budget. In particular
\[
 a'_x\ge a/2,
\]
since
\(4N_s/k^2\le 1/4\) follows from \(N_s+2r\le k\),
\(r\ge2\), and \(E_0\le a/32\).

At each top there is a unique occurrence after the incoming
frame. Choose a coordinate digit vector \(\ell_x\in[0,P_{j,x})^k\)
at this occurrence, with total \(|\ell_x|\equiv a'_x\pmod{P_{j,x}}\).
All digit choices have \(|\ell_x|<kP_{j,x}\le a'_x\). Write
\[
 n_x=(a'_x-|\ell_x|)/P_{j,x}.
\]
At a top strictly between frames
the degree in the segment is \(n_x\).

At an outgoing frame row
of segment \(j\) choose a degree \(e_x\in[0,rH_j)\) with
\(e_x\equiv n_x\pmod{H_j}\). Use degree \(b_{j+1,x}=(n_x-e_x)/H_j\)
when the same top reappears in the incoming frame. In fact
\[
 b_{j+1,x}>\frac{a'_x}{P_{j+1,x}}-k-r,
 \qquad
 W_{j+1,x}=\left\lfloor\frac{a'_x}{P_{j+1,x}}\right\rfloor-k-r-1,
 \qquad b_{j+1,x}\ge W_{j+1,x}.
\]
The strict inequality uses \(H_j\ge1\) and \(e_x<rH_j\). Our size
bounds give \(a'_x/P_{j+1,x}\ge4Yk^2N_s\), so
\[
 W_{j+1,x}\ge m_0=kN_s+1
\]

for \(m_0=kN_s+1\). The digit and frame degrees thus add to exactly
\(a'_x\) at every top, with nonnegative incoming padding.

\paragraph{The segment words.}
Fix the digits of the intervening tops in segment \(j\). The pre-outgoing path degree at \(x\) is
\[
 B_x=\begin{cases}n_x,&\text{intervening top},\\ W_{j,x}-m_0,&\text{incoming row}.\end{cases}
\]
Write \(S=\sum_x B_x\)
for this sum before the outgoing frame. Consider paths \(\mu\) in
\(k\) rows with these degrees. The last \(r\) steps are vector
steps of variable size. Just after the \(i\)-th outgoing step
require that the first \(i\) entries are all \(S\), with indices
starting at one.

Write \(\nu,\kappa\) for the restrictions to the
first \(r\) and the last \(K\) entries. Each is a signed Pieri
path with sizes determined by its weight changes. Each covector step creates at most one further negative entry. Thus
\[
 \#\{i:\mu_i<0\}\le N_s\le K-r,
\]
so the first \(r\) entries of the dominant weight \(\mu\) stay nonnegative. They stay at most \(S\) before the outgoing
frame by the sum of previous vector degrees. They stay at most
\(S\) afterward by the head condition and dominance. The
complement has zero degrees in the incoming frame because the
first \(r\) steps are vectors.

Impose the admissible path inequalities. At every
late top \(x\), \(\nu_r\) before the step is at least \(2B_u/3\)
for the earliest \(u\) of the preceding \(r\) vectors. The
complement increase at a vector and the plane decrease at a
covector are at most
\[
 \frac{\mathcal L a}{pkP_{j,x}}+
   \mathbf 1_{x\text{ even}}\frac{4a}{P_{j,u}},
 \qquad \mathcal L=C_{\mathcal L} N_s^2\log a.
\]
Let \(\mathcal T_j\) be the set of such paths with a choice of
standard terminal tableau for \(\kappa\). Its size depends only on
the intervening degrees and the fixed data.

For the plane word use
the tableau of its terminal rectangle \((S,\ldots,S)\), which is
unique. Let \(d_i\) be the full strip size at outgoing row \(i\).
It is at most \(2S\). At a vector step from \(\mu\) to \(\mu'\),
interlacing gives \(\sum_h(\mu'_h-\mu_h)\le\mu'_1-\mu_k\).
Here \(\mu'_1=S\), and \(\mu_k\ge-S\) since the total of all
covector degrees is at most \(S\).

\proofstep{Counting the outgoing decorations.}
For each member of \(\mathcal T_j\), choose outgoing decoration
exponents \(v_{i,h}\ge0\) for the \(r\) plane coordinates of
outgoing row \(i\), with
\begin{equation}\label{eq:plane-outgoing-decoration}
 v_{i,h}+2S<H_j,\qquad d_i+\sum_h v_{i,h}=e_i.
\end{equation}
Use the product of the plane and complement words, the outgoing
monomial with these exponents, and the same incoming factor for
all paths at the chosen degrees. The incoming factor is the
plane determinant of the incoming frame raised to \(m_0\), times
one coordinate of each incoming row to degree \(b_{j,x}-W_{j,x}\).
The product is in the indicated lattice \(M_j\). The signed strip
interlacing at the plane boundary supplies the two inequalities
between plane and complement at each late cut. The incoming determinant power shifts every plane coordinate of the weight by \(m_0\).

If the old minimum is \(\nu_r\), then
\[
 \min_i(\nu_i+m_0)=\nu_r+m_0.
\]
The other added vector factors can only increase this minimum.
The path bound on \(\nu_r\) before the step gives the absolute
lower cutoff independently. The normal degree of the product
at a vector is the complement strip size and at a covector it
is the plane strip size. The path inequalities give the normal
cutoff. Finally the exponents in each outgoing plane coordinate
are less than \(H_j\) by \eqref{eq:plane-outgoing-decoration} and
the bound
\[
 d_i\le2S.
\]


\proofstep{Independence in fixed degrees.}
These products are independent at each fixed set of degrees,
also over the residue field. First omit their common nonzero
incoming factor. In the \(i\)-th outgoing plane row specialize
all but coordinate \(r+1-i\) to zero. For each column
\(c\in[1,S]\) of the plane word, the height increases by zero
or one at each outgoing step. A rise wedges the sole surviving
coordinate there. Reading back from the final full exterior
subset, remove each of the coordinates at its rise. This leaves
precisely the terminal subset for the column at the pre-outgoing
cut, the complement of its rise coordinates. The coefficient is
\(1\) or \(-1\) by the wedge rule.

If \(R_c\) is the set of outgoing rises in column \(c\),
its pre-outgoing terminal subset is
\[
       [1,r]\setminus\{r+1-i:i\in R_c\}.
\]
Thus the specialization of the product word is, up to sign,
\[
 \prod_{i=1}^r x_{i,r+1-i}^{\,\#\{c:i\in R_c\}}
\]
times the
truncated plane word with those terminal subsets and the whole
complement word. The variables \(x_{i,r+1-i}\) are the surviving
outgoing coordinates.
The terminal subsets of the truncated plane word are standard.
At every outgoing cut the height on the left is at least the
height on the right.

If their current heights are \(h_L\ge h_R\), the remaining rises to the common terminal height obey
\[
 r-h_L\le r-h_R.
\]
Let \(D_L,D_R\) be the sets of removed coordinates in two adjacent columns. The surviving coordinates run in reverse order, so for every initial interval \([1,b]\), 
\[
 |D_L\cap[1,b]|\le|D_R\cap[1,b]|.
\]
 The retained
subsets therefore have weakly increasing entries across a row.
The truncated plane path and its subsets recover all
plane heights, since each missing coordinate specifies the time
of one rise. The two paths and their tableaux are therefore
recovered from the standard words in the variables other than
the outgoing plane variables.

To include the decorations in this argument, take the terms of
least total degree in the discarded coordinates in a putative
relation. For each tuple of exponents in those coordinates and
in the surviving outgoing plane coordinates, the displayed
specializations are independent by the standard word basis
in the remaining variables. Once the word indices are known, let \(a_i\) be their fixed antidiagonal exponents. The observed exponents \(b_i\) then recover each decoration by
\[
 d_i=b_i-a_i.
\]
This proves
independence and also explains why it holds with the same
coefficient ring as the residue rank calculation.

\end{proof}

\paragraph{Weight of paths at the segment cutoffs.}
\label{par:plane-typical}
\begin{lemma}[Weight of the paths obeying the cutoffs]
\label{lem:good-plane-path-count}
For segment \(j\) of Definition~\ref{hyp:coarsened-plane-frames}, let \(\mathcal T_j\) be the paths satisfying the degree cutoffs, with a choice of standard terminal tableau for the complement weight \(\kappa\). With \(D_k(b)=\binom{b+k-1}{k-1}\), their number satisfies
\begin{equation}\label{eq:plane-good-path-count}
 |\mathcal T_j|\ \ge\ (1-a^{-N_s^2})
    \prod_{x\text{ between frames}}D_k(n_x)
    \prod_{x\text{ incoming}}D_k(W_{j,x}-m_0).
\end{equation}
\end{lemma}

\begin{proof}
We prove this at the simultaneous cutoffs above. Let \(m\le N_s\)
be the number of letters in the segment and number the first
\(m-r\) steps before the outgoing frame. Use the hypotheses
\(k,\mathcal L\ge1088N_s^2\log a\). From the degree-loss bound
\(E'_x+E_x\le5b_x/(kY)\) we have \(a'_x\ge9b_x/10\).
For either kind of pre-outgoing row,
\[
 \frac{4a}{5P_{j,x}}\le B_x\le\frac{2a}{P_{j,x}},\qquad B_x\ge k.
\]
In fact the floor and padding lose at most
\[
 4kN_s\le b_x/(10P_{j,x}),
\]

since \(b_x/P_{j,x}\ge8Yk^2N_s\). Thus degrees of successive
rows of the same parity grow by at least
\[
 \rho=p^2/5>13.
\]

For any letter \(x\), each preceding pre-outgoing degree is at
most \(4\beta_x/p\), where \(\beta_x=a/P_{j,x}\), by strict
increase of the top positions.

First count all paths with the outgoing head conditions,
weighted by the number of tableaux of the terminal complement.
Before the outgoing frame the signed Pieri rules decompose the
tensor product of the indicated symmetric powers of the standard
module and its dual. Each path to a weight \(\lambda\) has weight
\(\dim V_\lambda\). After \(t\) outgoing steps remove the first \(t\)
entries, all equal to \(S\), and call the remaining row list the
tail.

At the next step its old and new tails \(\eta,\eta'\) obey
\[
 \eta_i\ge\eta'_i\ge\eta_{i+1}.
\]
Conversely every such
interlacing tail reinstates a vector strip on inserting the
extra head equal to \(S\), because all entries were at most
\(S\). Thus at the outgoing steps the ordinary branching rule
from \(\mathrm{GL}_{k-t}\) to \(\mathrm{GL}_{k-t-1}\) preserves the
total weight. These multiplicity-one rules and the Weyl dimension
formula apply to rational weights by a determinant twist and
duality. \cite{fulton-harris}
The last tail is the complement, whose tableau count is its
Weyl dimension. The total count before imposing the good path
inequalities is exactly
\[
 R=\prod_{x\text{ pre-outgoing}}D_k(B_x).
\]


\proofstep{The constraints on a path and its two signs.}
For a given path and sign \(\epsilon\in\{+,-\}\), let
\(\alpha^\epsilon\) be the nonincreasing list of positive parts
or of absolute negative parts of its tail. Up to the outgoing
frame the tail is the full list of \(k\) rows. Let \(\mathcal B^\epsilon\)
be the list of previously used pre-outgoing degrees of the
corresponding sign, with vectors positive and covectors negative.
For every \(T\ge0\) put
\[
 \delta_\epsilon(T)=\sum_{B\in\mathcal B^\epsilon}\min(B,T)
                    -\sum_h\min(\alpha^\epsilon_h,T).
\]
Each of these deficits starts at zero and is nondecreasing along
the path.

At a pre-outgoing step of the corresponding sign let
\(b\) be the strip size and \(\alpha,\alpha'\) the old and
new lists for this sign. Monotonicity and interlacing give
\[
 0\le\sum_h\bigl(\min(\alpha'_h,T)-\min(\alpha_h,T)\bigr)
           \le \min(b,T).
\]
The \(b\) bound sums the coordinate increments. For the
\(T\) bound use \(\alpha'_{h+1}\le\alpha_h\) and telescope. At a step of the opposite sign the
second sum can only decrease. During the outgoing frame both
lists of parts decrease entry by entry by the tail interlacing,
and the reference lists stay fixed.

Write
\[
 \Xi_x=\mathcal L\beta_x/(pk).
\]
We bound the weight of
paths violating any of the following conditions at a late top
\(x\): a covector decreases positive mass by at most \(\Xi_x\),
a vector decreases negative mass by at most \(\Xi_x\), and the
last plane entry before either step is at least \(2B_u/3\).
Here \(u\) is the earliest of the preceding \(r\) vectors, which
is before the outgoing frame. These conditions imply the good
path inequalities. At a covector all plane entries are
nonnegative, and the complement's positive mass cannot
increase, so the plane loss is at most the positive mass loss.

For a complement weight \(\kappa\), denote the positive and negative masses by \(m^+=\sum_i(\kappa_i)_+\) and \(m^-=\sum_i(-\kappa_i)_+\). At a vector step the change in the sum of the entries is 
\[
 \Delta\!\sum_i\kappa_i=\Delta m^+-\Delta m^-,
\]
 the positive mass gain plus the negative mass loss. Interlacing bounds its
positive mass gain by the last plane entry before the step.
Reading that entry backwards through the last \(r-1\) vector
steps bounds it by the first entry just after \(u\). Geometric growth of the vector degrees gives
\[
 \text{first entry just after }u\ \le\!\sum_{x\le u,\ x\text{ a vector}} B_x\le 2B_u.
\]
Using the bound
\[
 2B_u\le4a/P_{j,u},
\]
the complement
gain is at most
\[
       \text{negative mass loss}+2B_u
                      \le \Xi_x+4a/P_{j,u}.
\]

\proofstep{The loss from the mass test.}
For a violation of the mass condition let \(U\) be the sum
of preceding degrees of the opposite sign. If there is any
mass of that sign, at least one such pre-outgoing degree has
appeared. Its largest degree is at least \(k\) and at most
\(4\beta_x/p\). The geometric sum gives
\[
 U\le5\beta_x/p.
\]

The total mass of the sign before this step is at most \(U\)
by the deficit inequality at \(T\ge S\). Hence for \(T\ge U\)
the deficit of that sign grows by more than \(\Xi_x\). This
includes negative mass losses at outgoing steps, since removing
an entry equal to \(S\) preserves the negative mass.

\proofstep{The loss from the last-entry test.}
For a violation of the last-plane-entry condition, put \(B=B_u\).
If \(t\) outgoing steps have passed at this cut, \(u\) is the
oldest of the
\[
 h=r-t
\]
most recent pre-outgoing vectors.
The sum of positive rows after the first \(h\) in the tail is
at most the sum of all vector degrees strictly before \(u\).
At each of those \(h\) vector additions, interlacing bounds the
positive sum below row \(i+1\) after the addition by the sum
below row \(i\) just before. Covectors and the later outgoing
trims can only decrease these sums. Thus the positive sum after
row \(h\) is at most \(B/(\rho-1)\). Row \(h\) has positive part
less than \(2B/3\).

For
\[
 B\le T\le13B,
\]
the \(h\) latest pre-outgoing
vectors contribute at least \((h-1)T+B\) to the reference sum.
In the clipped tail, the first \(h-1\) rows contribute at
most \((h-1)T\), row \(h\) at most \(2B/3\), and the rest at
most \(B/(\rho-1)\). Subtracting gives
\[
 \delta_+(T)\ge \big((h-1)T+B\big)
                 -\big((h-1)T+2B/3+B/(\rho-1)\big)\ge B/4
\]
by the stipulated growth of the vector degrees.

\proofstep{Converting exceptional events into lost weight.}
To convert these deficits into a weight loss, set
\(J_k(T)=\sum_{l=1}^{k-1}(T+l)^{-2}\). For \(b\ge0\),
\[
 \log D_k(b)=\int_0^\infty\min(b,T)J_k(T)\,dT,
\]
where the left side for real \(b\) denotes the logarithm of the
binomial product. The equality follows by integrating each
\((T+l)^{-2}\) separately.

At the terminal complement weight
\(\lambda\) in dimension \(K\), the positive and negative lists
bound its dimension by
\[
 \dim V_\lambda\le\prod_h D_k(\alpha_h^+)D_k(\alpha_h^-).
\]
Indeed each Weyl factor is bounded by
\[
 1+\frac{\lambda_i-\lambda_j}{j-i}
 \le\left(1+\frac{\max(\lambda_i,0)}{j-i}\right)
     \left(1+\frac{\max(-\lambda_j,0)}{j-i}\right).
\]
Multiplying gives the bound with \(D_K\), which is at most
\(D_k\) for each nonnegative argument. Therefore
\(\log(R/\dim V_\lambda)\) is at least the integral of
\((\delta_++\delta_-)J_k\) at the last cut. If the mass bound
fails, this integral is at least
\[
 \Xi_x\sum_{l=1}^{k-1}\frac1{U+l}
 \ge \frac{\mathcal L(k-1)}{9k}\ge\frac{\mathcal L}{16},
\]
using \(k\le4\beta_x/p\). If the plane cutoff fails, it is at
least
\[
 \frac B4\sum_{l=1}^{k-1}
     \left(\frac1{B+l}-\frac1{13B+l}\right)
 \ge\frac{11(k-1)}{104}\ge\frac{k}{16}.
\]
Here \(B\ge k\), and both estimates are independent of the
remaining path and its strip degrees.

At each cut the two lists of nonzero parts have lengths at
most \(m\), since each pre-outgoing strip introduces at most
one new part of its sign. The outgoing interlacing keeps that
bound after removing the heads. Each part is an integer in
\([0,S]\). The lists padded to length \(m\), and the time of
the cut, determine the whole shape. There are at most
\((1+S)^{2m^2}\) shape paths because the zero start is fixed.
Thus the exceptional weight is at most
\[
 R\exp\bigl(2m^2\log(1+2aN_s)-\min(k,\mathcal L)/16\bigr)
 \le R a^{-N_s^2}.
\]
The last inequality uses \(1+2aN_s\le a^2\), supplied by
\eqref{eq:plane-count-size}, and the bounds on \(k,\mathcal L\).
This proves \eqref{eq:plane-good-path-count} for each choice
of intervening digits with the same constants.
\end{proof}


\paragraph{Sum over outgoing rows and digits.}
The last gap before the first outgoing position exceeds \(L_1\).
Thus each pre-outgoing \(x\) has
\(P_{j,x}\ge p^L p^{L_1+1}/H_j\). The positions are distinct, and
\(B_x\le 2a/P_{j,x}\). Summing the geometric series and using
\(a/p^L<8pYk^2N_s\) gives
\[
 S\le \frac{4aH_j}{p^L p^{L_1+1}},\qquad
 4YS\le H_j .
\]

For each outgoing row of a fixed path, the congruence
\(d_i+\sum_h v_{i,h}\equiv n_i\pmod{H_j}\) has exactly
\(H_j^{r-1}\) solutions with \(0\le v_{i,h}<H_j\). Each coordinate
is uniform among these solutions since \(r\ge2\). The union
bound removes at most \(2rS H_j^{r-2}\) solutions on requiring
\(v_{i,h}<H_j-2S\) for all \(h\). Each surviving solution has
\(d_i+\sum_h v_{i,h}<rH_j\), and so supplies a permitted outgoing
degree. The count per row is at least
\(H_j^{r-1}(1-r/(2Y))\) when this factor is nonnegative.

\begin{proposition}[Rank after summing the coarsened families]\label{prop:coarse-plane-rank}
With the admissible frame degrees and path count
\eqref{eq:plane-good-path-count}, the coarsened generators
span a residue space \(\overline\Lambda\) with
\[
 \dim\overline\Lambda\ge q^{-rK}(1-Nr^2 k/Y)\prod_xD_k(b_x).
\]

\end{proposition}

\begin{proof}
\proofstep{Summation using the degree at the outgoing row.}
For fixed digits the number of decorated words depends only on the outgoing degree
of its own segment. The incoming padding changes with adjacent
outgoing degrees, but the independence and membership just
proved hold at every such degree. Hence the sum of products
of word counts over the outgoing choices factors across segments.
For the digit sum at a top between frames, dividing each of the
\(k\) exponents of a monomial of degree \(a'_x\) by \(P_{j,x}\)
gives the identity \(\sum_{\ell_x} D_k(n_x)=D_k(a'_x)\).

At a frame top there are exactly \(P_{j,x}^{k-1}\) digit choices
by the sum congruence. Its incoming path factor in the next
segment is \(D_k(W_{j+1,x}-m_0)\), so its full contribution is
at least
\[
 \frac{1-r/(2Y)}{H_j^K}\,
 (P_{j,x}H_j)^{k-1}D_k(W_{j+1,x}-m_0).
\]
This is where the complement dimension \(K\) in the exponent
arises. The product of the factors \(H_j^{-K}\) over all \(r\)
rows in each frame is \(q^{-rK}\).

For the other factors, \(E'_x+E_x\le5b_x/(kY)\).
The binomial product formula and Bernoulli's inequality give
\(D_k(a'_x)\ge(1-5/Y)D_k(b_x)\). For an incoming radix
\(P'=P_{j,x}H_j\), put \(B'=W_{j+1,x}-m_0\). The floor and
padding lose at most \(4kN_s\) in the coarse degree.

Since
\(a'_x/P'\ge 4Yk^2N_s\),
\[
 P'B'\ge a'_x-4kN_sP'\ge(1-1/(kY))a'_x,\qquad
 (P')^{k-1}D_k(B')\ge(1-1/Y)D_k(a'_x).
\]
For the latter comparison use \(P'(B'+i)\ge(1-1/(kY))(a'_x+i)\)
for each \(1\le i\le k-1\). If \(1-Nr^2k/Y\le0\), the
claimed lower bound is automatic. Otherwise every loss factor
is in \([0,1]\). Since \(a\ge32Y\), \(J\le N\), and \(rJ\le N\),
the product of all loss factors is at least
\[
 1-N(5+1+r/2+1/32)/Y\ge 1-Nr^2k/Y.
\]

The next bound counts independent residue symbols in the
coarse radices. The following paragraph then checks that the
same families are admitted in the original radices, as needed
for Lemma~\ref{lem:lift-plane-symbol}.

To see that the dimension counts add, recover the choices from
the exponent at each top. For a coordinate \(h\) in a shared
frame row, write its true exponent as
\[
 \ell_{x,h}+P_{j,x}\bigl(a_{x,h}+H_j b_{x,h}\bigr).
\]
The digit \(\ell_{x,h}\) is the remainder modulo \(P_{j,x}\).
The outgoing exponent \(a_{x,h}\) is less than \(H_j\)
by the coordinate condition, so division of the quotient by
\(H_j\) recovers \(a_{x,h}\) and \(b_{x,h}\). In a complement
coordinate \(b_{x,h}=0\). The sum of the outgoing exponents
recovers \(e_x\). At unshared tops there is no second quotient.
Hence multiplication induces an injective map on the monomial
tuples, and distinct choices of digit and outgoing degree give
disjoint monomial supports. By
\eqref{eq:plane-good-path-count} and the count just proved,
\begin{equation}\label{eq:plane-coarse-rank}
 \dim \overline\Lambda\ \ge\ q^{-rK}(1-Nr^2k/Y)\prod_xD_k(b_x).
\end{equation}

\end{proof}

\paragraph{Return to the original radices.}
\label{par:plane-rank-transfer}

The count so far used coarsened digit classes. Returning to the original radices determines the rank available for the actual symbols.

\begin{proposition}[Rank of symbols at the original radices]\label{prop:plane-original-symbol-rank}
In the setting of Proposition~\ref{prop:coarse-plane-rank},
uncoarsening and multiplication by the reserved determinants
give integral original symbols of degrees \(b_x-E'_x\)
satisfying the required prefix tests. Their span contains a direct summand of rank \(r_{\rm orig}\) satisfying the bound in \eqref{eq:plane-coarse-rank},
\[
 r_{\rm orig}\ge q^{-rK}(1-Nr^2k/Y)\prod_x D_k(b_x).
\]

\end{proposition}

\begin{proof}
It suffices to return the products of the counted words to
the original symbol module. Use \(P^0_{j,x}=p^{z_j-x}\),
\(F_j=P_{j,x}/P^0_{j,x}=p^{R_j}\). Euclidean division in the unique late segment gives
\[
 \ell_{x,h}=\ell^0_{x,h}+P^0_{j,x}u_{j,x,h},\qquad
 0\le\ell^0_{x,h}<P^0_{j,x},\quad 0\le u_{j,x,h}<F_j .
\]

To
uncoarsen a counted segment product, raise its plane and
complement factors to \(F_j\). Multiply each factor by the
corresponding coordinate monomial of exponents \(u_{j,x,h}\)
in the late rows. Multiply the plane factor also by the
ordinary determinants of that segment with powers \(g_{j,x}\).
All coefficients are integers. Indeed, for an integral factor \(A\) with integer coefficients, Frobenius gives
\[
 \overline{A(X)^{F_j}}=\overline{A}(X^{F_j})
\]
in the residue field, with coordinatewise powers on the right. Applying the original radices in every
segment and multiplying by the remaining digits \(\ell^0\)
therefore gives \(\overline{\mathcal D}\) times the coarsened
product. In the integral polynomial the total degrees are
the coarsened true degrees plus the exact determinant degrees
\(E_x\), hence equal to \(b_x-E'_x\).

\proofstep{The inequalities for the restored products.}
We check the original inequalities on these integral products.
The plane and complement letter supports are unchanged, since
the digits are placed only in late rows and the determinants
use earlier plane vectors. There are at most \(kN_s\) coordinate digits in a prefix and each has size at most \(F_j-1\). Their combined adverse change in separation is therefore at most
\[
 \sum_{x,h}u_{j,x,h}\le kN_s(F_j-1).
\]
An odd plane coordinate can lower the plane
minimum by at most its exponent, and an even complement
coordinate can raise the complement maximum by at most its
exponent.

The other coordinate types do not decrease the
separation. The determinant associated to a given cut lies
wholly before that cut, so it contributes \(g_{j,x}\) to the
minimum. All other determinants have nonnegative minimum
since they contain only vectors. Adding the determinant gain to the worst digit loss gives, at this cut,
\[
 \text{separation}\ge F_jm_0-kN_s(F_j-1)+g_{j,x}>m_*+g_{j,x}.
\]
For the strict inequality use the separation margin in \(m_0\) and
\(F_j>a/(8pYk^2N_sQ)>m_*\) from \eqref{eq:plane-count-size} and
the definition of \(L\).

\proofstep{The rank lower bound at a plane.}
For the absolute plane lower bound put \(M=b_u/P_{j,u}\) at
the earliest preceding vector \(u\). The count assumptions
include \(k\ge C_kN_s^2\log a\) with \(C_k\ge1088\). In
particular \(a'_u\ge9b_u/10\) by the bound on \(E_u+E'_u\).
Since \(u\) is pre-outgoing, the degree \(B_u\) used by the
count satisfies \(B_u\ge 9M/10-4kN_s\) for either an
incoming or an intervening row. Also \(M\ge8Yk^2N_s\ge
40kN_s\). The plane minimum before the cut in the uncoarsened
word is at least
\[
 F_j(2B_u/3-kN_s)\ge F_j M/2=\frac{b_u}{2P^0_{j,u}}.
\]

For the normal cutoff at \(x\), the extra digit degree is at
most \(kF_j\). The gap \(x-u>L_1\) and \(F_jP^0_{j,u}\le A\)
give \(kF_j\le a/(p^{L_1}P^0_{j,x})\). The term at an even
\(x\) from the coarsened normal bound is
\[
 4a/P^0_{j,u}\le a/(p^{L_1}P^0_{j,x}).
\]
Thus the total
above \(\mathcal L a/(pkP^0_{j,x})\) is bounded by
\(a/(p^{L_1}P^0_{j,x})\) at an odd top and by twice that at
an even top. These imply both of the original order bounds.

The original symbols produced this way have the thresholds,
normal cutoffs and degrees needed by the precision induction.
After the original radix substitutions, the total normal
degree at each top is bounded by
\(\mathcal L a/(pk)+10a p^{-L_1}+kQ\). Indeed there is one
late occurrence, the incoming contribution is tangential,
and the final digit costs at most \(kQ\). The same upper bound
holds on supports of sums, so the total over all tops is at
most
\[
 N(\mathcal L a/(pk)+10a p^{-L_1}+kQ)
\]
on the whole
retained lattice.

The determinant product has nonzero reduction. For each of
its factors the preceding \(r\) even positions are distinct
modulo \(n'\). An interval of length \(n'\) starting immediately
after the earliest of those positions contains all \(r\) even
positions of any given frame modulo \(n'\), each once. If the
last position were at least that far from the earliest, there
would be at least \(r\) later preceding positions in the same
segment. Therefore the determinant monomials remain distinct
under the positive-power substitutions and have coefficients
\(\pm1\).

Since the residue polynomial ring is a domain, multiplication by the nonzero determinant product has zero kernel,
\[
 \overline{\mathcal D}\sum_i a_iW_i=0\quad\Longrightarrow\quad\sum_i a_iW_i=0.
\]
Hence the
products counted in \eqref{eq:plane-coarse-rank} remain independent
after multiplying by \(\overline{\mathcal D}\). Complete their
independent reductions to a basis of the ambient fixed-degree
residue space. Lifting the additional basis vectors to the
integral monomial lattice gives a basis by a change of basis
with unit determinant. This gives an integral direct summand
of the indicated rank generated by original symbols.
\end{proof}

The preceding rank and lifting statements concern the symbol
with the reduced weights. In Proposition~\ref{prop:residue-plane-neighbor-test}
the source must have the full weight. The replacement below
retains the rank of the model family and supplies the integral
contractions used on the target side.

\paragraph{Restoring the bounded weights.}
\label{par:plane-middle-weights}
\hypertarget{source-plane-middle-weights}{}
At a top let \(\lambda=(\lambda_1,\ldots,\lambda_k)\) be its
true dominant weight and let \(b=\lambda_1-\lambda_2\). Remove
\((b,0,\ldots,0)\) to form the remaining dominant weight
\[
 \mu=(\lambda_2,\lambda_2,\lambda_3,\ldots,\lambda_k).
\]

Away from the tops take the removed weight to be zero and
\(\mu=\lambda\). At odd indices use the corresponding reversed
negative weights in the polynomial model. Write \(d_l(\gamma)\)
for the dimension of highest weight \(\gamma\) in dimension \(l\).
The local factor of the count for restoring weights is
\(d_{k-1}(\lambda_2,\ldots,\lambda_k)\) at a top and
\(d_k(\lambda)\) elsewhere. We construct integral polynomials
realizing exactly this count.

\begin{auxlemma}[Counting the tableaux at one index]\label{lem:bounded-tableau-count}
With the top weights \(\mu\) defined at the start of
paragraph~\ref{par:plane-middle-weights}, use tableaux of the
determinant-shifted shape. At an even top take tableaux with every entry in the first row
equal to zero. At an odd top take the reversed negative shape and exclude zero
from its alphabet. Writing \(n_s\) for the number admitted at position \(s\),
\[
 n_s=\begin{cases}d_{k-1}(\lambda_2,\ldots,\lambda_k),&s\text{ a top},\\d_k(\lambda),&s\text{ away from the tops}.\end{cases}
\]
Away from the tops every tableau is allowed. Distinct selected tableaux have distinct leading monomials over
every residue field.
\end{auxlemma}
\begin{proof}
Use a matrix of variables \(X=(X_{i,h})\), with rows and columns
numbered from zero. For a dominant weight \(\gamma\) the
polynomial realization after its determinant twist has shape
\(\gamma-\gamma_k(1,\ldots,1)\). Its integral basis consists of
products of initial-column minors indexed by semistandard
tableaux of this shape on the alphabet \(0,\ldots,k-1\).
For each tableau its leading monomial uses the diagonal of each
minor. Order the variables by column first and by increasing
row within a column, giving the earlier variable precedence.

In a minor the largest monomial uses the smallest available row
in the first column and then the smallest remaining row in each
successive column. For rows \(i_0<\cdots<i_{s-1}\), this gives
\[
 \operatorname{in}\det(X_{i_a,b})_{0\le a,b<s}
    =\prod_{b=0}^{s-1}X_{i_b,b}.
\]
The coefficient is one. If \(M_T\) is the diagonal monomial of the product for a tableau \(T\), then
\[
 \deg_{X_{i,b}}M_T=\#\{c:T(b,c)=i\}.
\]
Its weakly
increasing order reconstructs that row from these counts. These are the integral
standard monomial coordinates used for the component lattices.
\cite[Theorem II.2.16]{abw-schur}

Away from the tops take every tableau of \(\mu\) in the
corresponding signed realization. At an even top restrict to
those with first tableau row all zero. Removing that row
leaves the tableau of the smaller shape with alphabet
\(1,\ldots,k-1\), since every remaining box had a box above it.
This restriction counts
\[
 d_{k-1}(\lambda_2,\ldots,\lambda_k),
\]

by subtracting the common determinant shift. At an odd top
restrict to tableaux with no entry zero. The reversed negative
weight of \(\mu\) has two equal last entries. The determinant
shift leaves at least its last shape row empty. Dropping this
row and using the alphabet \(1,\ldots,k-1\) counts the dimension
of the reversed negative of \((\lambda_2,\ldots,\lambda_k)\),
which equals \(d_{k-1}(\lambda_2,\ldots,\lambda_k)\) by duality.
\end{proof}
Let \(n_2\) be the product of these counts and \(H_i\) the
resulting basis products in the independent block variables.

\begin{proposition}[Restoration of the bounded highest weights]\label{prop:restore-plane-model-rank}
Use the removed flat degrees and the \(n_2\) tableau products
\(H_i\) counted in Lemma~\ref{lem:bounded-tableau-count}. The products \(\theta(Q)H_i\) retain
independence modulo the uniformizer for every independent
family of flat polynomials \(Q\) of the prescribed degrees.
With the common factor \(C\) of \eqref{eq:plane-clean-factor}
and precision \(c'\ge\max(2(D_1+1),c_\star+o_S)\), the products
admit the stated clean approximations at the distinguished edge. With the bounded spread \(c_0\) and lower
degree \(a\), let \(f_{\rm true}\) and \(f_{\rm flat}\) be the retained fractions in the true and flat dimensions, respectively. In the indicated size range,
\[
 f_{\rm true}\ge(1-kc_0/a)^N f_{\rm flat}.
\]

\end{proposition}

\begin{proof}
\proofstep{Independence after restoring the removed weight.}
To prove independence after restoring the removed weight, let
\(\theta\) insert the first column at an even index and the
signed \((k-1)\)-row minors of the first \(k-1\) columns at an
odd index. This is the vector and covector polynomial model of
Section~\ref{sec:edge-transfer}. At an even top with removed
degree \(b\), the leading monomial of a degree-\(b\) monomial
in the flat variables inserts exponents \(\alpha_c\) at
\(X_{c,0}\). The exponent matrix of a chosen tableau has no
\(X_{c,0}\) for \(c>0\).

At each \(c>0\) the exponent of \(X_{c,0}\) in the leading product
is therefore precisely \(\alpha_c\). Recover the missing exponent from
the homogeneous degree:
\[
 \alpha_0=b-\sum_{c>0}\alpha_c.
\] At an odd top the leading
monomial of the cofactor omitting row \(c\) uses \(X_{t,t}\)
for \(0\le t<c\) and \(X_{t+1,t}\) for \(c\le t<k-1\).
In the selected tableau the entry in row \(t\) is at least
\(t+1\), so its exponent at \(X_{t,t}\) is zero. The diagonal exponents of the product read
\(d_t=\sum_{c>t}\alpha_c\) for \(0\le t<k-1\).

If \(b\) is the fixed flat degree, invert this triangular
reading by
\[
     \alpha_0=b-d_0,\qquad
     \alpha_c=d_{c-1}-d_c\ (1\le c\le k-2),\qquad
     \alpha_{k-1}=d_{k-2}.
\]
Thus the diagonal exponents and the degree recover the
full flat monomial. After recovering the flat exponents, subtract their contribution
to recover the leading exponent matrix of the tableau. At an
off-top position there is no contribution to subtract. Thus the
pairs of a flat monomial and a selected tableau product have
distinct leading monomials. Each has leading coefficient a unit,
including in the residue field. In a proposed nonzero relation
\[
 \sum_i\theta(Q_i)H_i=0,
\]
choose the largest of those leading
monomials among the nonzero coefficients. Its coefficient cannot
cancel with that of any other pair. This proves the assertion
over the residue field for arbitrary \(Q_i\) of the fixed degrees.

\proofstep{The three normalized edge lattices.}
Use the normalized edge coordinates on all three weights
\(\lambda,\mu,\lambda-\mu\). The normalization consists of
pullback to the component frames, the last-entry determinant
powers in the signed realizations, and the common edge powers
\(w^{A_r}\) for the respective weights. In each frame the common powers on the two factors have product
\[
 w^{A_r(\mu)}w^{A_r(\lambda-\mu)}=w^{A_r(\lambda)},
\]
by additivity of the edge exponent in the weight.

On a flat polynomial this change to edge
coordinates is the letter substitution followed by \(\theta\).
Indeed a vector transforms by the frame matrix \(M\), and the
cofactors by \(\det(M)M^{-t}\), which follows by multiplying
the adjugate identity by \(M\). The last-entry twist on the
covector cancels the determinant. The flat edge exponent is
one copy of the degree for each even block and zero for each
odd block. Write
\[
 D=\operatorname{diag}(w^{-1}I_r,I_{k-r})
\]
for the diagonal
part. After the common determinant cancels, the two letter maps are
\[
 wD=\operatorname{diag}(I_r,wI_{k-r})\quad\text{on vectors},\qquad
 D^{-t}=\operatorname{diag}(wI_r,I_{k-r})\quad\text{on covectors}.
\] The
other frame factors act by the same matrix identities, with
Frobenius at the corresponding index. Thus this normalization
pulls the edge coefficient lattices back to their integral
standard monomial models, including their reductions.

In the distinguished frame the model vectors \(S_b\) have
normalized polynomials \(B_b C\) with one common factor \(C\).
This is \(\operatorname{subst}(\mathrm{SPtrue})\) in the symbol
identity, with residue the product of \(F_*\) and the separate
determinant reserve on the digits. Each determinant uses distinct
rows and has a permutation monomial with unit coefficient, so
\(\overline C\ne0\). To obtain the clean source, expand \(B_b\)
by its segment presentation. Its nonnegative separation margin
allows each homogeneous pair to be expanded in the paired paths
of Section~\ref{par:plane-column-orders}. The common factor
\begin{equation}\label{eq:plane-clean-factor}
                 C=F_*\prod_s T_s^{kQ}
\end{equation}
supplies the precision reserve and the factors for the digits. Write \(S_{b,o}\) for the standard-edge pullback of the
normal degree-\(o\) component of \(C B_b\). Apply the path
and digit constructions to this component.

In the flat source
lattice \(R_e\) on the chosen edge, the resulting vector
\(F_{b,o}\) has source order at most \(o\) and satisfies
\[
                 F_{b,o}-w^o S_{b,o}\in w^{o+c'}R_e .
\]
The insertions commute with
the projection to the full true multidegree because \(C\)
has the reserved degrees. Divide by \(w^o\) in the displayed bound,
then sum over the finitely many \(o\) bounded in
Section~\ref{par:plane-rank}. Thus
\[
  F_b:=\sum_o w^{-o}F_{b,o}\in S_b+w^{c'}R_e .
\]
Here \(S_b=\sum_o S_{b,o}\). The competing-edge conditions
of the construction give \(F_b\in w^{c'}R_{e'}\) on every
other edge out of the same source. Restore the weights by multiplying
with the inverse edge transports of the \(H_i\). They are integral
on the chosen edge, so the same congruence holds in the restored
lattice there.

If a combination with coefficients \(c_{b,i}\in O\) is
zero modulo \(w\) in the normalized edge lattice, then in the
residue polynomial model
\[
 0=\sum_i\theta\left(\overline C\sum_b\overline{c_{b,i}}\overline{B_b}\right)
                    \overline{H_i}.
\]
All polynomials inside \(\theta\) have the full common flat
multidegree. By block independence, for each \(i\) the coefficient polynomial
vanishes. Cancellation takes place in the residue polynomial ring,
a domain, and then independence of the \(\overline{B_b}\) gives
\[
 \overline C\sum_b\overline{c_{b,i}}\overline{B_b}=0
 \ \Longrightarrow\ \sum_b\overline{c_{b,i}}\overline{B_b}=0
 \ \Longrightarrow\ c_{b,i}\in wO\quad\text{for every }b .
\] In particular primitive combinations remain primitive
in the actual source lattice.

\proofstep{Pairing the restored model family.}
For the pairing one can therefore test the restored model family.
An integral target test evaluates the difference between the clean
and model combinations in \(w^{c'}O\). If \(c'\ge 2(D_1+1)\), this
is contained in \(p^{D_1+1}O\), since \(w^2=p\). A model evaluation
of valuation at most \(D_1\) has the same valuation on the clean
combination. On competing source edges restoration costs at most
\(o_S\) powers of \(w\), by the source vertex bound for the remaining
weights. Thus \(c'\ge c_\star+o_S\) gives the source conductor there.

At each top the spread \(\lambda_2-\lambda_j\) for \(j\ge2\)
is at most \(c_0\), with \(b\ge a\) in the plane count.
The Weyl factors involving the first entry give
\[
 d_k(\lambda)\le d_{k-1}(\lambda_2,\ldots,\lambda_k)D_k(b)
                  (1+c_0/a)^{k-1}.
\]
This follows term by term from
\(\lambda_1-\lambda_j+j-1\le (b+j-1)(1+c_0/a)\).
The retained fraction of the true dimension is therefore at
least the flat fraction multiplied by \((1-kc_0/a)^N\), using
\((1+t)^{-(k-1)}\ge1-(k-1)t\) for \(t\ge0\). The product of
the off-top dimensions is included exactly.
\end{proof}

The pairing of two lifted symbols is computed from the unit parts of factorials in the prescribed digit classes. The identities below give the
precision for that pairing in the proof of
Proposition~\ref{prop:residue-plane-neighbor-test}.
\paragraph{Factorials and divided derivatives in a digit class.}
\label{par:plane-digits}
\hypertarget{source-plane-digits}{}
We record two uniform identities used in constructing the signed-chain
pairing. Write \(P=p^g\) and \(0\le\rho<P\), and use the residue
class \(\rho+P\beta\) for integers \(\beta\ge0\). Congruences in
this paragraph are coefficientwise over \(\mathbb Z_{(p)}\) and
therefore hold over the coefficient DVR as well.

For \(M\ge0\), let \(U(M)=\prod_{1\le v\le M,\,p\nmid v}v\),
with \(U(0)=1\). Removing the multiples of \(p\) from a factorial
and repeating \(g\) times gives
\begin{align*}
 (\rho+P\beta)!&=\beta!\,p^{A_\rho+\beta B_g}
       \prod_{h=0}^{g-1} U(b_h+P_h\beta),\\
 b_h&=\left\lfloor\frac\rho{p^h}\right\rfloor,\quad
 P_h=p^{g-h},\quad A_\rho=\sum_{h=0}^{g-1}\left\lfloor\frac{b_h}{p}\right\rfloor,
 \quad B_g=\frac{P-1}{p-1}.
\end{align*}
For each \(h\) put \(C_h=U(P_h)\). Since \(b_h<P_h\), splitting
into full blocks followed by the last partial block gives
\begin{align*}
 U(b_h+P_h\beta)={}&U(b_h)C_h^\beta\,
 \prod_{j=0}^{\beta-1}\prod_{\substack{1\le v\le P_h\\p\nmid v}}
                (1+P_hj/v)\,
 \prod_{\substack{1\le v\le b_h\\p\nmid v}}(1+P_h\beta/v).
\end{align*}
The two products on the right together form a unit. In a group
of coordinates with the same \(g\), the powers of \(p\) and \(C_h\)
contribute a constant depending on the total \(\sum_a\beta_a\)
and the fixed residues. The remaining unit is the product of
the displayed products of \(1\) plus a multiple of \(p\).
This gives the factorial factorization with the exact degree
of the quotient monomials as the argument of the constant.

\begin{lemma}[Factorial units and divided derivatives at fixed digits]\label{lem:plane-digit-units-and-cuts}
On a fixed digit class \(\rho+P\beta\), with \(P=p^g\),
\(g\ge0\), \(0\le\rho<P\), \(\beta\ge0\), and \(c\ge1\), either
sign of the variable unit factor in the factorial is congruent
modulo \(p^c\) to a product-binomial polynomial of total degree
at most \(2(c-1)\). A system of divided derivatives of
orders \(M_a\) admits the quotient expansion
\eqref{eq:plane-digit-derivative-expansion} with the following bound when the common radix is \(P\),
\[
 \sum_a\kappa_a\le\frac1P\sum_a M_a+c-1.
\]
 In a signed cut a divided derivative of order
\(|M|\) loses at most \(|M|\) on a vector letter and zero
on a covector letter. The factorial formula and expansions
include the radix-one case.
\end{lemma}

\begin{proof}
\proofstep{Uniform polynomial descriptions of units and inverse units.}
For \(c\ge1\), the remaining unit and its inverse both have binomial
polynomial approximations to precision \(p^c\) of total
degree at most \(2(c-1)\). To prove it, replace the leading
factor \(p\) in each \(P_hj/v\) and \(P_h\beta/v\) by a variable
\(z\). For either sign, expand the factors to degree \(c-1\)
as polynomials in \(z\) modulo \(z^c\).

In a full block at \(j\)
the coefficient of \(z^d\) is a scalar in \(\mathbb Z_{(p)}\)
times \(j^d\). Multiplying over \(j\) expresses the coefficient
of \(z^d\) as a combination of sums
\[
 \sum_{0\le j_1<\cdots<j_l<\beta}j_1^{d_1}\cdots j_l^{d_l},
 \qquad d_i\ge1,\quad \sum_i d_i=d.
\]
Induct on the number \(l\) of indices. For the induction step,
fix the last index \(j_l\). The sum over its predecessors is
an integer-valued polynomial in \(j_l\) of degree at most
\((d-d_l)+(l-1)\). Multiplication by \(j_l^{d_l}\) gives
degree at most \(d+l-1\). Summing from zero to \(\beta-1\)
raises the degree by at most one, since
\[
 \sum_{j<\beta}\binom jt=\binom{\beta}{t+1}.
\]

This proves the degree bound \(d+l\) and integrality of
the values at nonnegative integers. A partial block contributes an ordinary
polynomial of degree \(d\) in \(\beta\) at \(z^d\). The coefficient
at total order \(d\) over all blocks and coordinates is thus a
polynomial of total degree at most \(2d\), taking values in
\(\mathbb Z_{(p)}\) at nonnegative integer tuples. Write \(P(\boldsymbol\beta)\) for this coefficient polynomial and
\(\Delta_i\) for forward difference in coordinate \(i\). Its Newton
expansion is
\[
 P(\boldsymbol\beta)=\sum_{|\mathbf h|\le2d}
    (\Delta_1^{h_1}\cdots\Delta_m^{h_m}P)(0)
    \prod_{i=1}^{m}\binom{\beta_i}{h_i},
\]
where \(m\) is the number of coordinates. Each coefficient is an
integral linear combination of values of \(P\), hence lies in
\(\mathbb Z_{(p)}\). Now
put \(z=p\). This proves the approximation uniformly in the
number of coordinates and quotient exponents, also when the
fixed residues and powers \(g\) differ by group. If \(g=0\)
its product is empty and the same formulas give unit one.

\proofstep{Divided derivatives and their factorial exponents.}
For the divided derivative, write \(\partial^{[M]}X^n=\binom nM
X^{n-M}\), with the derivative zero for \(n<M\). Use
\[
 (1+T)^P=1+T^P+pE(T),\qquad \deg E\le P-1,
 \qquad E\in\mathbb Z[T].
\]
Then \(\binom{\rho+P\beta}{M}\) is the coefficient of \(T^M\) in
\((1+T)^\rho(1+T^P+pE(T))^\beta\). At order \(\ell<c\) in \(pE\),
the part with \(T^{Pj}\) from the other factor has coefficient
\[
 p^\ell[T^{M-Pj}](1+T)^\rho E(T)^\ell\,
       \binom{\ell+j}{\ell}\binom\beta{\ell+j}.
\]
This uses the binomial identity for selecting \(\ell\) marked
factors and then \(j\) others. The support of this expression
requires \(j\le M/P\) and \(M-Pj\le\rho+\ell(P-1)\).

Put \(t=\lceil(M-\rho)/P\rceil\) and \(\rho'=\rho-M+Pt\), so
\(0\le\rho'<P\). With \(\kappa=\ell+j\), the support condition
gives \(\kappa\ge t\). Therefore the corresponding divided
derivative term has the form
\begin{equation}\label{eq:plane-digit-derivative-expansion}
 \begin{aligned}
 &c_{\ell,j}X^{\rho'}
       \left(X^{\kappa-t}\partial^{[\kappa]}X^\beta\right)(X^P),\\
 &c_{\ell,j}=p^\ell[T^{M-Pj}](1+T)^\rho E(T)^\ell
                                      \binom{\ell+j}{\ell}.
 \end{aligned}
\end{equation}
All constants are integral and independent of \(\beta\).
Terms with \(\beta<\kappa\) vanish. For multiple coordinates
multiply these expansions and keep \(\sum_a\ell_a<c\), so
\[
 \sum_a\kappa_a\le\sum_a M_a/P+c-1.
\]
The discarded coefficients
are divisible by \(p^c\), also after grouping equal monomials.
This is the required derivative expansion inside the old residue
class, with the new residues \(\rho'_a\) fixed by \(M_a\).
The case \(P=1\) has \(E=0\) and gives the ordinary divided
monomial derivative exactly.

\proofstep{Effect of a derivative on a signed cut.}
In the same argument a prefix lower bound is the least Laurent
exponent for the conjugates of \(\operatorname{diag}(T,1,\ldots,1)\)
acting on the letters in the cut. For any rank-one projection
\(P_1\), write this matrix as
\[
 A_T=1+(T-1)P_1.
\]

The inverse is \(1+(T^{-1}-1)P_1\). A vector coordinate in the
cut has least order at least zero and a covector coordinate
has least order at least \(-1\). A divided derivative of
order \(|M|\) in one letter loses at most \(|M|\) in the vector
case and zero in the covector case. To see this, write the
substitution of that letter as \(x\mapsto B_Tx\). Expand the
polynomial after substitution at \(x+B_T^{-1}h\) and take the
coefficient of \(h^M\). Thus the Laurent orders of the entries satisfy
\[
 \operatorname{ord}_T\!\big((B_T^{-1})_{ij}\big)\ge\begin{cases}-1,&B_T^{-1}=A_T^{-1}\quad(\text{vector}),\\0,&B_T^{-1}=A_T^t\quad(\text{covector}).\end{cases}
\]
For a letter
outside the cut the Jacobian is the identity. The coefficient
estimate holds for each conjugate and hence for the cut bound.

Adjoining vector letters to a cut composes its substitution
with polynomial substitutions on those separate letters.
It can only increase the least order. Removing them composes
with inverse substitutions and loses at most the sum of their
degrees by homogeneity in each removed letter.  They give the local cutoff losses
when the binomial approximation is implemented by divided
derivatives and multiplication by monomials of the same degree.
\end{proof}

To obtain the covector in Proposition~\ref{prop:residue-plane-neighbor-test},
we compare the actual pairing with its model on prescribed cut
subspaces. The corrected pairing estimate below uses the same original multidegree as the source lift in
Lemma~\ref{lem:lift-plane-symbol}.
\paragraph{The Fock metric and the signed cuts.}
\label{par:plane-gram}
\hypertarget{source-plane-gram}{}
We give the pairing argument in its signed polynomial model.
Normalize the valuation by \(v(p)=1\). For a monomial with
exponent \(\alpha\), write \(\alpha!=\prod_x\alpha_x!\) and
\(h_\alpha=v(\alpha!)/2\). Use the bilinear form and its
associated coefficient bound
\[
  [F,G]=\sum_\alpha F_\alpha G_\alpha\,\alpha!,
  \qquad \nu(F)=\min_\alpha \bigl(v(F_\alpha)+h_\alpha\bigr).
\]
Then \(v([F,G])\ge \nu(F)+\nu(G)\). In true multidegree
\(d=(d_i)\), set
\[
 n_0=\frac12 v\!\left(\prod_i d_i!\right),\qquad
 K_d=\sum_i r\bigl(1+\log_p(d_i+r)\bigr),\qquad
 e_d=r^2\bigl(1+\log_p(2\textstyle\sum_i d_i+r)\bigr).
\]

In this multidegree \(2n_0-K_d\le v(\alpha!)\le2n_0\).
Indeed the difference is the valuation of a product of
multinomial coefficients. Each multinomial is a product of
at most \(r\) binomial coefficients with upper entries at most
the letter degree. A binomial with upper entry \(n\ge1\) has
valuation at most \(1+\log_p n\), by the base-\(p\) carry
formula, which follows from Legendre's formula for factorials.
The same argument for a split \(\alpha=A+M\) gives
\[
   v(A!)+v(M!)\ge 2n_0-2K_d,
   \qquad \text{when \(A+M\) has multidegree \(d\)}.
\]

\begin{lemma}[Projectors for signed cuts in the factorial norm]\label{lem:plane-signed-projectors}
For every prescribed cut and threshold in fixed multidegree
\(d\), there is a Fock-self-adjoint rational projector onto
the signed cutoff subspace, whose loss in \(\nu\) is at most
\(e_d\). For nested or disjoint cuts write \(P_i\) for their projectors and \(P\) for the stipulated product. Then
\[
 P^2=P,\qquad \operatorname{im}P=\bigcap_i\operatorname{im}P_i,
\]
with loss at most \(e_d\) per cut.
It preserves multidegree and the monomial exponents outside
the union of those cuts.
\end{lemma}

\begin{proof}
\proofstep{The projector for a signed cut.}
Using the signed, generically conjugated Laurent order, the space attached to a cut \(S\) is
\[
 U(S,b)=\{F:\operatorname{ord}_S(F)\ge b\}.
\]
In fixed
multidegree it is the sum of the \(\mathrm{GL}_r\)-types
for the action on \(S\) whose least weight \(\lambda_r\) is
\(\ge b\). Indeed all torus exponents in a type are at least
\(\lambda_r\) in the selected coordinate, and the conjugates
of that torus test detect the lowest weight in every nonzero
copy. More explicitly the intersection of the conjugate
weight-cut spaces is invariant, and the copies it contains
are exactly the stated types. The type sum has a Fock
self-adjoint projector \(\Pi\), defined over \(\mathbb Q\).

By the trace formula and divided-power estimate for type
projections proved above, its matrix on the coefficient lattice satisfies
\[
 v_p(\Pi_{\beta\alpha})\ge-e_d\qquad(\text{all }\alpha,\beta).
\]
In applying the
estimate the coordinate span on the cut is at most \(\sum_i d_i\),
and it suffices to use \(2\sum_i d_i\) in the Hom weight bound.
Zero degree and \(r=1\) give scalar projections on the affected
monomial spaces. The trace formula works for the sum of types
with the same bound by the ultrametric inequality.

For the matrix entry \(\Pi_{\beta\alpha}\), self-adjointness
gives
\[
 \beta!\Pi_{\beta\alpha}=\alpha!\Pi_{\alpha\beta}.
\]

Consequently the two numbers
\(v(\Pi_{\beta\alpha})+h_\beta-h_\alpha\) and
\(v(\Pi_{\alpha\beta})+h_\alpha-h_\beta\) are equal.
Their common value is the mean of the two unweighted
entry valuations, each bounded below by \(-e_d\).
For the coefficient of \(x^\beta\) in \(\Pi F\) the
ultrametric inequality therefore gives
\[
 v((\Pi F)_\beta)+h_\beta
    \ge\min_\alpha\{v(F_\alpha)+h_\alpha-e_d\}.
\]
Taking the minimum in \(\beta\) yields the claimed norm bound.

For nested cuts the smaller-cut projector commutes with
the action on that cut and with the action on its complement.
It therefore preserves every isotypic component for the
larger cut and commutes with its type projector. For disjoint
cuts the projectors act on separate factors. A product over
such cuts is self-adjoint, fixes their intersection and loses
at most \(e_d\) per cut. It preserves multidegree and each
exponent outside the union of the cuts. This proves the projection and gives the product bound used below.

\end{proof}

\begin{paramdefinition}[The true and model lattices, cuts and signs]\label{hyp:plane-gram-data}
The block letters \(B_j\) in each frame \(j\) have specified
signs and total block degree \(b_j\). The shared letters \(X_t\),
indexed by \(t=(j,m)\), \(0\le m<r\), are vectors. Let
\(P_t=p^{g_t}\) be their radices. The model has the same
block letters and three vector letters \((L_t,W_t,T_t)\)
over each shared letter. Its collapse \(\phi\) sends them
respectively to \(X_t,X_t^{P_t},X_t\), coordinate by coordinate.
The exponents on \(L_t,W_t,T_t\) will be denoted by
\(\rho_t,\beta_t,M_t\). An admissible monomial has
\[
 \begin{gathered}
 0\le\rho_{t,a}<P_t, \\ M_{j,0}=0, \\ |M_{j,m}|\le b_j,
\end{gathered}
\]
and exact collapsed multidegree \(d\).

\smallskip
The cuts are defined as follows.
A true link cut contains \(B_j\) and the shared letters
\(X_{j,m'}\) for \(m'>m\). Denote it by \(A_{jm}\);
we use \(m+2\le r\). The corresponding model link cut
\(U_{jm}\) contains \(B_j\), the letters \(T_{j,m'}\)
for \(m'>m\), and all the \(W_{j-1,m'}\), where the frame
indices are cyclic.

A model chain cut \(S_{ji}\) contains a
nested subset of \(B_j\) and all the \(W_{j-1,m'}\).
Thus the link cuts alone, the chain cuts alone, and the
union of the chain cuts with the model link cuts are each
families of nested or disjoint cuts. Set
\[
 D_l=J(r-1)e_d,\qquad D_m=e_d\sum_j n_j,
\]
where \(J\ge1\) is the number of frames and \(n_j\) the
number of chain cuts at \(j\). The model total degree is
at most \(\sum_i d_i\) on admissible monomials because
each \(P_t\ge1\).

Use a diagonal unit \(\gamma(\alpha)\) on the true
monomials. Its two signs admit integral binomial polynomial
approximations of total degree at most \(6c\) modulo \(p^c\).

Suppose that \(\nu(F)=0\) in true multidegree \(d\), that
the true link orders of \(F\) are at least \(\ell_{jm}\),
and that each digit part of \(F\) has chain orders at least
\(\theta_{ji}\). The digit part for a fixed residue keeps the
block exponents and reads on \(W_t\) the quotient exponent
\[
 m_t=\rho_t+P_t\beta_t\quad\longmapsto\quad \beta_t=(m_t-\rho_t)/P_t,
\]
with no \(L_t,T_t\) exponents.
The pairing under consideration is \([F,\gamma\,\phi(Z)]\),
where \(\gamma\) acts coefficientwise. Use the inequalities
\begin{equation}\label{eq:plane-gram-reserves}
 \begin{split}
    c&>4(D_m+D_l+3K_d),\\
    s_{ji}&\ge r\, b_{j-1}/P_{j-1,m}+13rc
                     \qquad(0\le m<r).
 \end{split}
\end{equation}
\end{paramdefinition}
The family of test polynomials is indexed by primitive vectors of the true lattice. The two signs of the diagonal unit are treated by Lemma~\ref{lem:plane-digit-units-and-cuts}. The reserves are inequalities on the integer conductor chosen for the small-prime calculation; its specialization is made in Section~\ref{sec:host-edge-tests}.


\begin{proposition}[An integral polynomial with prescribed orders at the model cuts]\label{prop:plane-corrected-gram-polynomial}
Use the lattices, orders and sign approximations in
Definition~\ref{hyp:plane-gram-data}, including \(\nu(F)=0\)
and the reserves \eqref{eq:plane-gram-reserves}.
Then there is an integral model polynomial \(Z_*\) of collapsed
degree \(d\), with model link orders at least
\(\ell_{jm}-30c\) and chain orders at least
\(\theta_{ji}-s_{ji}-12c\), such that
\[
 v([F,\gamma\phi Z_*])-n_0+K_d
                               \le D_l+D_m+4K_d.
\]
The order and valuation bounds are uniform over the coefficients,
digit classes and radices with the data of Definition~\ref{hyp:plane-gram-data}.
\end{proposition}

\begin{proof}
\proofstep{Projecting a monomial to the link cuts.}
Take a monomial \(x^\alpha\) attaining \(\nu(F)=0\) and
project it to the true link cuts at levels \(\ell_{jm}\).
The result \(Y\) has \(\nu(Y)\ge h_\alpha-D_l\), and
\[
 [F,Y]=F_\alpha\,\alpha!
\]
has valuation \(h_\alpha\).
Apply an integral binomial approximation \(g\) of \(\gamma^{-1}\)
of degree at most \(6c\) to its coefficients. A factor
\(\binom{\alpha_x}{a}\) acts as
\(x^a\partial_x^{[a]}\). Its loss in a signed cut is at
most \(a\), by the cut estimates for the derivative and
multiplication. This bounds the total link loss by \(6c\).
The change in the pairing from \([F,Y]\) to
\([F,\gamma(gY)]\) has valuation at least
\[
 c+h_\alpha-D_l>h_\alpha.
\]

Ordinary coefficients of \(gY\) have valuation at least
\(h_\alpha-D_l-n_0\). Multiplication by
\[
 p^{\lceil n_0+D_l-h_\alpha\rceil}
\]
therefore makes \(gY\)
an integral polynomial \(G\) in the original variables. It has link orders at least
\(\ell_{jm}-6c\) and its twisted pairing has valuation at
most \(n_0+D_l+1\).

\proofstep{Lifting \(G\) in the prescribed plane.}
To lift \(G\), first put each \(X_t\) entirely in \(L_t\).
Within frame \(j\), order its block letters first and then
the shared letters in decreasing order of \(m\). Each
\(A_{jm}\) is a prefix of this order. Use the integral basis
of signed path words proved in Section~\ref{sec:column-lifts},
now testing the least torus weight at the prefixes.
In a positive column the least exponent selects a
subset avoiding the torus coordinate whenever \(h<r\).
In a nonpositive column it selects a dual subset
containing the coordinate whenever \(h<r\). Thus if
\(h_b\) are the prefix heights,
\[
 \lambda_r=
 \sum_{b\ {\rm positive}}\mathbf1_{h_b=r}
       -\sum_{b\ {\rm nonpositive}}\mathbf1_{h_b<r}.
\]
The Pieri decomposition at the successive prefixes
selects exactly the paths whose displayed sums obey all
the lower bounds. The associated words span the
intersection of the cuts. Since they are a
subset of the integral basis, the expansion of an integral
polynomial in this intersection uses integral coefficients.

In a path word, a column meeting at least one block
letter moves all its shared occurrences with \(m>0\) from
\(L_{j,m}\) to \(T_{j,m}\). Keep \(L_{j,0}\) and every
occurrence in a column disjoint from the blocks. Each
column uses a given shared letter at most once. Count the columns by one of their block occurrences. Each chosen occurrence lies in only one column, so
\[
 \#\{\text{moved columns}\}\le\#\{\text{block occurrences}\}\le b_j.
\]
This gives the tangent degree bounds.
On a moved column the new cut \(B_j\cup\{T_{j,m'}:m'>m\}\)
acts on exactly the old cut occurrences. An unmoved
column has zero order there. Its contribution to the
prescribed lower bound was at most zero.

In fact for a
positive column the height at the cut is bounded by the
number of its occupied vector letters before the cut,
which is \(<r\) since the column meets only shared letters.
For a prefix, a moved column retains its signed cut
order and an unmoved column contributes zero, at least
its previous order. Minimum Laurent orders add on products, giving the required cut
bound for the whole word. For the coefficient valuation the relevant
identity is
\[
      \nu_{\rm coeff}(AB)=\nu_{\rm coeff}(A)+\nu_{\rm coeff}(B)
                 \qquad(A,B\ne0).
\]
Indeed, divide each factor by its minimum uniformizer power. Their
reductions are nonzero polynomials over the residue field, so their
product has nonzero reduction. The injective renaming on each
column preserves that column's \(\nu_{\rm coeff}\), and
hence preserves the valuation of the product word.
Collapsing \(T_{j,m}\) back to the corresponding shared
letter restores each column separately.

Perform the operation coefficientwise outside each frame.
Since the renamed letter sets are disjoint, that operation
commutes with the cut substitutions in the other frames.
For a remaining shared variable, divide its exponent by its radix.
The lift and collapse act on the monomial by
\[
 a_t=P_tq_t+r_t,\quad 0\le r_t<P_t,\qquad
 L_t^{r_t}W_t^{q_t}\:\longmapsto\: X_t^{r_t+P_tq_t}=X_t^{a_t}.
\]
The tangent exponents stay fixed throughout. To see that it preserves
the link bounds, group by the \(L\)-exponents. These letters
are outside the new link cuts, so coefficient extraction
preserves the bounds. The added monomial in the \(W\) letters
has nonnegative order at every model link cut, since those
letters are vectors. The integral model polynomial \(Z\)
therefore has admissible support, \(\phi(Z)=G\), and
\(\operatorname{ord}_{U_{jm}}Z\ge\ell_{jm}-6c\).

\proofstep{Correction to impose the chain cuts.}
For an admissible model exponent \(\zeta\), let \(A\) be
its collapsed exponent without \(T\), and let \(M\) be its
tangent exponent on the true shared letters. Thus
\(\phi(x^\zeta)=x^{A+M}\). The digit factorial identity
gives
\[
       A! M! = \kappa_B\,u_\zeta\,\zeta!,
       \qquad
       \eta_\zeta=\gamma(A)u_\zeta ,
\]
where \(u_\zeta\) is the product of the digit units.
The constant \(\kappa_B\) depends on the individual label and
tangent exponents and on the model multidegrees, which we fix
in each block of the next projection. Explicitly it is the
product of the digit constants for each \(W_t\)-degree divided
by the product of the label factorials. The letter \(T_t\)
contributes the same factorial on both sides.

\proofstep{Diagonal correction factors.}
Choose diagonal multipliers \(g_+\) and \(g_-\) approximating
\(\eta_\zeta\) and its inverse modulo \(p^c\), with
integral values and cut cost \(12c\) on cuts excluding \(L\).
For fixed label exponents, substitute
\(\rho_{t,a}+P_t\beta_{t,a}\) into the binomial approximations
to \(\gamma^{\pm1}\). Apply forward differences in \(\beta\) to the integer-valued polynomial \(A(\beta)=\binom{\rho+P\beta}{k}\). Its expansion is
\[
 A(\beta)=\sum_{j=0}^k(\Delta^j A)(0)\binom{\beta}{j},\qquad (\Delta^j A)(0)\in\mathbb Z.
\]

Multiplying this approximation by the degree-\(6c\) digit
unit approximation yields the cost \(12c\). Extracting
coefficients of \(L\) leaves every such cut intact, so the
estimate holds after adding the label classes.

In each model multidegree let \(\Pi_m\) be the product
projector on chain orders at least \(\theta_{ji}-s_{ji}\).
Use its linear extension over all multidegrees. It fixes
the exponents on \(L,T\), preserves the admissible support
conditions, and commutes with all model link projections.
Each multidegree component of a polynomial satisfying a
link bound still satisfies it, by the homogeneous cut action.
 Set
\[
 \Omega=g_-\Pi_m g_+.
\]

Its chain orders are \(\ge\theta_{ji}-s_{ji}-12c\).
Any model link bound is lowered by at most \(24c\).
For \(\zeta,\zeta'\) in the same fixed block of \(\Pi_m\),
\[
 \lvert v(\zeta!)-v(\zeta'!)\rvert\le K_d.
\]

Only the block and \(W\) exponents can vary. Their degrees
are bounded by the corresponding true degrees, so the
multinomial estimate applies without a second factor for
the fixed \(L,T\) exponents. The factorial-weighted projection
bound proves that \(p^{N_2}\Omega\) preserves integrality
on admissible support, with
\(N_2=\lceil D_m+K_d/2\rceil\).

\proofstep{Precision of the corrected pairing.}
For a fixed tangent exponent \(M\) put
\[
         F_M=\gamma^{-1}\,\partial^{[M]}(\gamma F).
\]
Every coefficient of \(F_M\) has valuation \(\ge -n_0\):
the unit multipliers and the integral divided derivative
preserve the original coefficient bound. At a fixed label
\(\rho\), its digit part has a decomposition \(Q+E\) such
that
\[
 v(E_\beta)\ge c-n_0
\]
and
\(\operatorname{ord}_{S_{ji}}Q\ge\theta_{ji}-s_{ji}\)
for all \(j,i\). To obtain it, start with the digit parts
of \(F\). In each of the two twist operations replace the
unit multiplier by its degree-\(6c\) approximation. Each
replacement error has coefficient valuation at least
\(c-n_0\), and each approximation loses at most \(6c\) in
the chain orders, using the affine binomial expansion.
Between the two multipliers use the divided derivative
expansion of the preceding subsection on each shared letter.

The divided-derivative expansion takes the input label to the residue of \(\rho-M\)
coordinatewise. The error coefficients have valuation at
least \(c-n_0\). The approximate derivatives and all
multiplying monomials have integral coefficients, so errors
already made retain this valuation. Only derivatives on
the \(W\) letters in the cut can lower its order. Their
total loss is bounded by
\[
    \sum_{m=0}^{r-1}
       \left(\frac{|M_{j-1,m}|}{P_{j-1,m}}+c\right)
    \ \le\ s_{ji}-12c
\]
by \eqref{eq:plane-gram-reserves} and \(r\ge1\).
This produces \(Q,E\) with the asserted relation.

\proofstep{Comparing factorial terms in the corrected pairing.}
For an admissible monomial the factorial identity and the
coefficient formula for the divided derivative give
\[
 [F,\gamma\,\phi(x^\zeta)]
       = \kappa_B\,\eta_\zeta\, (F_M)_A\,\zeta! .
\]
In the model multidegree of \(\zeta\), multiply \(Q\) by
the fixed label and tangent monomial, then keep that
multidegree. The resulting polynomial \(Q_B\) still has
the required chain orders because \(L,T\) are outside
those cuts. On that block the difference
\[
    g_-(\zeta')\,\eta_{\zeta'}(F_M)_{A'}-(Q_B)_{\zeta'}
\]
has valuation at least \(c-n_0\). The pairing with
\(Q_B\) is unchanged by \(\Pi_m\).

To estimate the remaining
error from a coefficient of \(\Pi_m x^\zeta\), use its bound
\(h_\zeta-h_{\zeta'}-D_m\). From the split factorial
bound applied to \(\zeta\) and \(\zeta'\), and the equality
of their \(\kappa_B\), we have
\[
 v(\kappa_B)+h_\zeta+h_{\zeta'}\ \ge\ 2n_0-2K_d .
\]
The error terms in comparing the evaluation before and
after \(\Pi_m\) thus have valuation at least
\(c+n_0-2K_d-D_m\), including the original monomial term.
Applying \(g_+\) at the input preserves this error estimate.
The factor \(g_+(\zeta)g_-(\zeta)-1\) has valuation at least
\(c\), and the original evaluation has valuation at least
\(n_0-K_d/2\) by \(\nu(F)=0\) and the true factorial bound.
Linearity and integral coefficients of \(Z\) now give
\begin{equation}\label{eq:plane-gram-stability}
 v\bigl([F,\gamma\phi(\Omega Z)]-[F,\gamma\phi(Z)]\bigr)
               \ \ge\ c+n_0-2K_d-D_m .
\end{equation}

\proofstep{The available precision after the conductor loss.}
The strict conductor inequality in \eqref{eq:plane-gram-reserves}
makes \eqref{eq:plane-gram-stability} larger than
\(n_0+D_l+1\). Hence the integral test
\(Z_*=p^{N_2}\Omega Z\) has valuation of its pairing at
most \(n_0+D_l+1+N_2\). Its model link and chain orders are
\(\ge\ell_{jm}-30c\) and \(\ge\theta_{ji}-s_{ji}-12c\),
respectively. Its collapse has true multidegree \(d\) and
integral coefficients, hence \(\nu(\phi Z_*)\ge n_0-K_d\).
In particular
\[
 v([F,\gamma\phi Z_*])-n_0+K_d
                 \le D_l+D_m+4K_d .
\]
Here \(K_d\ge1\) absorbs the two roundings. The estimate is uniform and needs no error polynomial.
The chain and link losses and the valuation gain are
uniform over all input coefficients, label classes, and
radices at the specified degrees and cuts.
\end{proof}

The Gram test must apply to the actual span of allowed
operators. We verify the prefix bounds on that span before
using the chosen-plane projection in
Proposition~\ref{prop:residue-plane-neighbor-test}.
\paragraph{Prefix orders on the exact operator span.}
\label{par:plane-orbit}
\hypertarget{source-plane-orbit}{}
In the distinguished edge coordinates the flat generators of the model
are exactly \(B_b\,\mathrm{SPR}\). Indeed pulling back each row of the
true frame determinant by the edge substitution gives its designated
plane row, so the product \(\mathrm{SPtrue}\) pulls back to
\(\mathrm{SPR}\). Every monomial in this common factor uses only
even plane coordinates. We work over the fraction field in their
linear span and apply words in \(X_{sa}\partial_{sb}\), using each
top \(s\) at most \(c_0\) times. Conjugation by an invertible linear
substitution within each letter takes a row operator to a linear
combination of row operators at the same letter. In matrix-unit notation the substitution by \(g\) acts on a row operator through
\[
 gE_{ab}g^{-1}=\sum_{c,d}g_{ca}(g^{-1})_{bd}E_{cd}.
\]
Each term acts at the same top, so the word allowance persists in edge coordinates.

List the tops after the incoming frame of segment \(j\) in increasing
order as \(x_0,\ldots,x_{h+r-1}\), where \(h\) is the number strictly
between the frames. The last \(r\) positions are the outgoing frame,
with shared index
\[
 m=r-1-(u-h)
\]
at position \(x_u\).
Write \(B_u\) for the true degree of the oldest of the \(r\) even
predecessors of \(x_u\), divided by its radix in this segment. Put
\[
 s_u=\lfloor B_u/(Yk^2)\rfloor.
\]

For a joint complement type,
let \(k_u\) denote its largest weight at the prefix of the first
\(u\) of the listed tops. In this paragraph frame subscripts on
\(B_u,s_u,k_u\) are implicit. Define \(T_u\) as the largest of
the following applicable quantities:
\[
 \frac{B_u}{2},\quad k_{u+1}+m_*+s_u\quad(x_u\text{ even}),\quad
 k_{u-1}+m_*+s_u\quad(u>0,\ x_{u-1}\text{ odd}).
\]
The chain and link thresholds for the plane Gram argument are
\begin{equation}\label{eq:plane-chain-link-thresholds}
   T_u-s_u/8\quad(0\le u\le h),\qquad
   T_u-3s_u/8\quad(h<u<h+r),
\end{equation}
respectively. Its separate chain
slack is \(3s_u/8\). We show that these thresholds work for
every choice of low digits, quotient multidegrees, joint
complement type and coefficient of that type.

\begin{proposition}[Prefix thresholds on the full permitted operator span]\label{prop:plane-operator-prefix-orders}
On the fraction-field span of the flat generators \(B_b\,\mathrm{SPR}\)
and the allowed words in \(X_{sa}\partial_{sb}\) with at most
\(c_0\) operators at each top, the chain and link thresholds from \eqref{eq:plane-chain-link-thresholds} are
\[
 \begin{cases}T_u-s_u/8,&0\le u\le h,\\T_u-3s_u/8,&h<u<h+r.\end{cases}
\]
 They hold after taking any indicated digit class,
quotient multidegree and joint complement type, with the same
allowance \(s_u\) for the selected coefficient.
\end{proposition}

\begin{proof}
\proofstep{Expansion into segment generators.}
First expand a generator in its presentation by products of
segment pairs \(A_jG_j\) in the plane and complement
coordinates. Do this before selecting a coefficient. The external
label \(l\) at each top is strictly below the segment radix there.
Take homogeneous components at the prescribed true degrees of
\(B_b\) in this expansion. Renamed products of homogeneous segment
factors have definite true degrees, and the sum of the components
in other degrees is zero. Thus the degree bound below applies on
each retained product. Expand \(\mathrm{SPR}\) into monomials
\(X^\gamma\).
In a coordinate of radix \(P\) and prescribed residue \(\rho\), the term can survive extraction only if
\[
 l+\gamma\equiv\rho\pmod P,\qquad (l+\gamma-\rho)/P\in\mathbb Z_{\ge0}.
\]
Apply this test in every coordinate.
In the matching class the quotient after multidegree projection
is a tensor product. Its complement factor is the projected
product of the \(G_j\) with their positions renamed.

The
complement is absent on incoming rows and the complement label
has zero quotient. Its plane factor, before multidegree projection,
is
\[
        X^b\,\phi\!\left(\prod_j \chi_j A_j\right).
\]
Here \(b\) is the plane quotient of \(l+\gamma\) by the low
radices, coordinate by coordinate. It vanishes on odd letters.
The variable map \(\chi_j\) sends intervening, outgoing, and
incoming letters respectively to the block letters,
\(L_{j,m}\), and \(W_{j-1,m}\). If \(P_{\rm low}\) is the low
radix on the preceding frame, the exponent multiplier at an incoming
row is
\[
           P_{\rm coarse}=P_{\rm low}\,P_{j-1,m},
                 \qquad P_{j-1,m}=p^{z_j-z_{j-1}}.
\]
This proves the displayed formula after collapse, since \(\phi\)
raises \(W_{j-1,m}\) to the \(P_{j-1,m}\)-th power.
These maps use the original frame radices after the transport in
Section~\ref{par:plane-rank}.

If the complement projection vanishes, the extracted plane
coefficient is zero. Otherwise its type at every prefix has
largest weight at most the upper Laurent order of the corresponding
prefix of \(G_j\). The other complement factors are outside the
cut, and the multidegree projection commutes with its action.
An isotypic projection with larger first weight would annihilate
the vector by the conjugate torus criterion. The original segment
inequalities now give
\[
 \operatorname{ord}_{<x_u} A_j\ge T_u.
\]

For an even letter use the complement through that letter. For a
preceding odd letter use its complement before it and the plane
through it. On segment support this last plane prefix is precisely
the prefix before \(x_u\). Its oldest even predecessor and hence
the available margin are unchanged in traversing the odd letter.
The absolute plane lower bound gives the first term of the maximum.
The extracted coefficient is a scalar multiple of the projected
plane factor. All these conditions describe linear spaces.

\proofstep{Transport of the prefix cuts by relabeling.}
For the chain cut \(S_{ju}\), the renaming \(\chi_j\) carries the
prefix before \(x_u\) to all incoming \(W\)'s and the first \(u\)
intervening letters with their original signs. Other segments
have no variables in this cut. Thus
\[
 \Pi=\prod_j\chi_j A_j
\]

has order at least \(T_u\). Restrict to a desired true quotient
multidegree of \(X^b\phi\Pi\) before applying the shared digit
extraction. On \(\Pi\) this restriction selects the model
multidegrees for which the weighted sum of letter degrees after
collapse has the prescribed value. Their projections commute with
the chain action. Expand the resulting polynomial by monomials
\(L^M\) of the label letters, which are outside the cut. Each
coefficient, a polynomial in block and \(W\) letters, retains the
order.

For each shared coordinate \(t\), divide the exponent in
\(X^b\phi(L^M)\) by \(P_t\):
\[
     b_t+M_t=\rho_t+P_t a_t,\qquad 0\le\rho_t<P_t,\quad a_t\ge0 .
\]
Thus the coefficient contributes only to the residue \(\rho_t\)
and acquires the fixed carry \(W_t^{a_t}\). For a tested residue \(\rho'_t\), the carry factor after extraction is therefore
\[
 \mathbf1_{\rho'_t=\rho_t}W_t^{a_t}.
\]
Apart from these carries the extra monomial
is the block part of \(X^b\). All affected coordinates of this monomial are vectors.
They have nonnegative least weight and multiplication preserves
the lower bound. This proves the required chain order for every
shared residue.

A true link at \(x_u\) with \(u>h\) sees the intervening letters
and the earlier outgoing rows. The original prefix also sees the
incoming rows. In a homogeneous component of \(A_j\), let their
total degree be \(d_W\). Let \(D(T)\) be the rank-one one-parameter torus
used for the original Laurent substitution.

To omit the
incoming variables \(W\) from this cut, substitute \(D(T)^{-1}W\)
in the transformed polynomial. A monomial of degree \(d_W\) accumulates \(d_W\) inverse weights, each either zero or \(-1\). Its Laurent order in this substitution is therefore
\[
 -\sum_{a=1}^{d_W}\varepsilon_a\ge-d_W,\qquad \varepsilon_a\in\{0,1\}.
\]
Therefore
\[
 \operatorname{ord}_{\rm prefix\ without\ incoming}(A_j)
 \ge \operatorname{ord}_{\rm original\ prefix}(A_j)-d_W .
\]
Coefficient extraction in the now untransformed \(W\)
preserves this bound for each remaining polynomial. Under collapse those incoming monomials go to frame
\(j-1\), outside this true link. Here \(J\ge2\) follows from the
frame span bound \(G'\) and \(3G'\le n'\), since one frame would
recur only after \(n'\) positions. All factors from other
segments and from \(b\) that meet this link use vectors there.
Projection to the required multidegree preserves the lower bound.

Only an incoming degree \(d_m^W\) with
\[
 P_{j-1,m}d_m^W\le (d_{\rm pl}(j-1,m)-|b_{j-1,m}|)_+
\]

can contribute. At the row \(y_m\) of that preceding frame put
\(t_m=z_{j-1}-y_m\), \(g=z_j-z_{j-1}\). Removing the quotient
of the external monomial leaves at most the original generator
degree \(2a\) at true scale. Indeed after multiplying the
remaining quotient by \(p^{t_m}\) the degree contributed by
the segment products is bounded by their prescribed true degree,
which is at most the flat degree. The \(t_m\) are distinct and
nonnegative. Consequently
\[
 d_W\le\sum_m 2a p^{-g-t_m}\le 3a p^{-g}.
\]

If \(y\) is the oldest even predecessor of the current outgoing
row, there are at least \(r\) intervening even tops before it is
tested. The first intervening top is more than \(L_1\) after
\(z_{j-1}\). Hence \(y\ge z_{j-1}+L_1+1\), and
\(a\le B_u p^{z_j-y}\). The host spacing gives
\(p^{L_1+1}\ge48Yk^2\); the size bound gives
\(B_u/(Yk^2)\ge1\), since \(B_u\ge a/Q\).
The permitted incoming degree is therefore bounded by
\[
 d_W\le \frac{3B_u}{p^{L_1+1}}
     \le \frac{B_u}{16Yk^2}\le s_u/8 .
\]
The last step uses \(B_u/(Yk^2)\le2s_u\).
Subtracting this degree from the full prefix order \(T_u\)
gives the initial link bound.

\proofstep{Admitting row operators in the order estimate.}
It remains to allow row operators. For one coordinate of low radix
\(P\), write its true power as \(v+Pq\), \(0\le v<P\).
Write \(x^v F(X)\) with \(X=x^P\). The two derivative
cases follow from the chain rule:
\[
 \partial_x\!\bigl(x^vF(x^P)\bigr)=
 \begin{cases}
   x^{v-1}(v+P X\partial_X)F(X),&v>0,\\
   x^{P-1}P\partial_X F(X),&v=0 .
 \end{cases}
\]
For multiplication the two cases are
\[
 x\,x^v F(X)=
 \begin{cases}
 x^{v+1}F(X),&v+1<P,\\
 X F(X),&v=P-1 .
 \end{cases}
\]
At \(P=1\) the second case applies to \(v=0\). The same
formulas at the second digit extraction act on the incoming
\(W\) quotient, or on the block coordinate at radix one. On an exponent \(v+Pq\), the Euler operator \(x\partial_x\) acts by \(v+Pq\). Thus on the quotient coordinate \(y\) in the residue class \(v\) its action is
\[
 v+P\,y\partial_y.
\]

Thus each multiplication or differentiation before extraction costs
at most one in a lower cut order that meets the letter. Under
conjugation the torus on a coordinate has weights \(0,1\) on a
vector and \(0,-1\) on a covector, with opposite signs on the
derivatives. The Euler operator also costs at most one, since one
of its two factors has nonnegative least weight. On a disjoint cut
these operations preserve the order.

For a complement operation, insert the joint decomposition of the
input before the operation and project the output. At each prefix \(I\), the upper-order estimate for an
operator at \(x\) allows a transition from \(k_I\) only to
\[
                          k'_I\le k_I+\mathbf1_{x\in I}.
\]
Complement coefficients commute with the plane cuts. Thus
this transition changes the maximum defining \(T_u\) by
at most one for a complement operation. It can
occur only for a top at or before \(x_u\) in that segment.
A plane operation can affect the shared \(W\) cut even if its
letter belongs to the previous frame, but that row occurs in the
incoming part of the segment, also before \(x_u\).
This accounts for both occurrences of a shared row, including
across the cyclic endpoint.

\proofstep{The exact operator span.}
If \(a_u\) operators in the word act at such tops up to the tested position \(u\), induction therefore bounds their cumulative loss by 
\[
 \underbrace{2+\cdots+2}_{a_u\,\text{operators}}=2a_u.
\]
 The
cut and type estimates for a single half-operation apply in all
input digit classes and degrees. The induction therefore survives
both quotient-degree changes and the sum over input types. For one cut, at most \(N_1c_0\) operators can touch its letters.
The host inequalities give
\[
 a\ge16N_1c_0Yk^2Q,
\]
and
\[
 B_u\ge a/Q
\]
implies
\[
 s_u\ge16N_1c_0.
\]

For a word let \(m_I\) be the number of operators meeting
the cut. The induction allows at most \(2m_I\) in the
chain and at most \(4m_I\) in the link, with both occurrences
of a shared row included. The budget \(m_I\le N_1c_0\) and
\(16N_1c_0\le s_u\) in particular give
\[
    \operatorname{ord}_{\rm chain}\ge T_u-s_u/8,\qquad
    \operatorname{ord}_{\rm link}\ge T_u-s_u/8-s_u/4 .
\]
The first loss is charged from the starting chain threshold
and the second from the initial link threshold.
These link thresholds are positive since \(T_u\ge B_u/2\)
and \(s_u\le B_u/(Yk^2)\). Linear combinations over the fraction field
preserve these inequalities. This proves the membership used for
the selected plane polynomial in Section~\ref{par:plane-flat-test}.
 \label{cor:plane-orbit-1}
\end{proof}

The preceding order and Gram results first give a test at
the flat weight. We use them to select a nonzero plane
component, realize its symbol, and restore the weight. This
supplies the detecting value used in
Proposition~\ref{prop:residue-plane-neighbor-test}.
\paragraph{Choosing the plane component.}
\label{par:plane-flat-test}
\hypertarget{source-plane-flat-test}{}
Let a flat source polynomial \(H\) in the full \(k\) coordinates
in the exact operator span of Section~\ref{par:plane-orbit} have
integral coefficients in the chosen edge frame after scaling.
Take \(H\) primitive there. At each letter use the radix
specified by its frame position. Fix the labels \(\rho\)
and the degrees of the plane and complement quotients using a
monomial of \(H\) with unit coefficient. In this class, write
\(H_{\alpha,\beta}\) for the coefficient at plane quotient
\(\alpha\) and complement quotient \(\beta\). Write
\((\alpha_0,\beta_0)\) for the two quotients of the selected
monomial. All the letter positions are distinct. Coordinatewise, radix \(P\), digit \(\rho\), and quotient \(q\) give
\[
 m=\rho+Pq,\qquad 0\le\rho<P.
\]

Taking \(q\) from the plane or complement quotient as appropriate therefore recovers every true exponent. Denote the assembly of a plane polynomial
\(A\) and a complement polynomial \(N\) by
\(\operatorname{Asm}_\rho(A,N)\): multiply the fixed label
monomial by the two polynomials with variables raised to
their letter radices. Their degrees are fixed by the chosen
class. The factorial identity of Section~\ref{par:plane-digits}
gives the exact formula
\begin{equation}\label{eq:plane-slice-factor}
 \begin{split}
 [H,\operatorname{Asm}_\rho(A,N)]
 &= b_\rho
     \sum_{\alpha,\beta} H_{\alpha,\beta}\,
       A_\alpha\,u_{\rm pl}(\alpha)\alpha!\,
       N_\beta\,u_{\rm no}(\beta)\beta!.
 \end{split}
\end{equation}
Here \(b_\rho\) depends on the labels and the quotient
multidegrees, and \(u_{\rm pl},u_{\rm no}\) are the products
of the factorial units in the two sets of coordinates.
The unit \(u_{\rm pl}\) is the twist \(\gamma\) in the
plane Gram estimate. Both signs have the required uniform
approximations. The formula is obtained monomial by
monomial, so it holds for the whole class.

\begin{lemma}[A nonzero plane polynomial and a complementary test]\label{lem:plane-split-polynomial}
Let \(H\) be an integral primitive flat polynomial in the
operator span of Proposition~\ref{prop:plane-operator-prefix-orders}.
Use the digits and quotient degrees of one unit-coefficient
monomial to determine the label class in \eqref{eq:plane-slice-factor}.
Then there
are a nonzero plane polynomial \(F\), an integral complementary
test \(N_0\), and a scalar \(k_0\ne0\) such that
\eqref{eq:plane-slice-value} and \eqref{eq:plane-slice-depth}
hold for every integral plane polynomial \(A\) at the selected
quotient degrees. In particular the scalar in \eqref{eq:plane-slice-depth} satisfies
\[
 v(k_0)\le V_d-n_0+D_n.
\]
The indices of the complementary types and their order losses are chosen before \(A\).
\end{lemma}

\begin{proof}
\proofstep{The projections onto the complementary factors.}
On the complement take the product projections onto one
isotypic type at each of its chain cuts. These cuts are
nested within a frame and disjoint across frames.
The products over all choices of types sum to the identity.
For \(Y=\sum_\beta H_{\alpha_0,\beta}x^\beta\), there is a
product \(E\) with \(\nu(EY)\le\nu(Y)\). Otherwise
\(\nu(\sum_E EY)>\nu(Y)\) by the ultrametric inequality,
contrary to the sum being \(Y\). In particular \(EY\ne0\).
For this \(E\), choose an exponent \(\tau\) attaining
\(\nu(EY)=q_1\), and put \(u=v(\tau!)/2\). The plane
polynomial \(P\) has coefficients
\[
 P_\alpha=\left(E\sum_\beta H_{\alpha,\beta}x^\beta\right)_\tau .
\]

The selected plane coefficient has \(v(P_{\alpha_0})=q_1-u\), and
\[
 q_1\le\nu(Y)\le v(\beta_0!)/2,\qquad
 q_3:=\nu(P)\le q_1-u+v(\alpha_0!)/2 .
\]
The second bound for \(q_1\) uses the unit coefficient.
Choose \(z\) with \(v(z)=q_3\), possible because the
coefficients have valuations in \(\frac12\mathbb Z\)
and \(w^2=p\). The normalized plane polynomial \(F=P/z\)
then has \(\nu(F)=0\). The orbit calculation gives the signed chain and link
orders for precisely this coefficient \(P\) of the projected
complement. Scalar multiplication by \(z^{-1}\) preserves them.

\proofstep{Scalar orders in the selected plane.}
Let
\[
 n_{0,{\rm no}}=\tfrac12v(\prod_i d_{{\rm no},i}!)
\]

for the complement degrees, and let \(D_{\rm no}\) be an
integer loss bound in the factorial norm for \(E\).
It is the sum of the cut projector losses, rounded up.
There is a scalar \(s\) of valuation
\[
 n_{0,{\rm no}}+D_{\rm no}-u.
\]

The complement polynomial \(s E x^\tau\) is integral
by the weighted projection bound. It has the upper torus
weight bound of the chosen type at each cut. Apply a binomial
approximation to \(u_{\rm no}^{-1}\) with integral coefficients
of degree \(\le 2(c_P-1)\) modulo \(p^{c_P}\).
This gives an integral complement test \(N_0\). In the upper
weight convention a diagonal binomial of degree \(a\) raises
the bound by at most \(a\), by applying the Laurent estimate
with the reciprocal torus variable. Thus the upper bounds of the chosen types increase by at most the permitted allowance \(4c_P\).

If one uses the exact inverse twist in place of the
approximation, Fock self-adjointness of \(E\) in
\eqref{eq:plane-slice-factor} gives exactly
\(k_0[F,\gamma A]\) with
\[
 k_0=b_\rho s\,\tau!z.
\]

Let \(V_d=v(\prod_i d_i!)\) be the full factorial valuation.
For each exponent in full multidegree let \(\kappa_{\rm fl}\)
bound its factorial deficiency from \(V_d\).
For example \(\sum_i k\lfloor\log_p d_i\rfloor\) works,
with a zero term if \(d_i=0\), by counting carries.
When \(A\) is integral the complement approximation gives
\begin{align}
 v\bigl([H,\operatorname{Asm}_\rho(A,N_0)]-k_0[F,\gamma A]\bigr)
       &\ge c_P+V_d-\kappa_{\rm fl}, \label{eq:plane-slice-value}\\
 v(k_0)&\le V_d-n_0+D_n, \label{eq:plane-slice-depth}
\end{align}
where \(n_0\) belongs to the selected plane degrees and
\[
 D_n=D_{\rm no}+
        n_0-\tfrac12v(\alpha_0!)+
        n_{0,{\rm no}}-\tfrac12v(\beta_0!).
\]
Indeed the factorial of the selected monomial has
valuation \(v(b_\rho)+v(\alpha_0!)+v(\beta_0!)\le V_d\).
Using \(v(z)=q_3\le v(\alpha_0!)/2+v(\beta_0!)/2-u\)
in the expression for \(v(k_0)\) gives
\eqref{eq:plane-slice-depth}. Assembly preserves coefficient
valuation bounds, giving \eqref{eq:plane-slice-value}
from the full factorial lower bound.

\end{proof}

\paragraph{Using a realised flat symbol.}
Apply the plane Gram estimate with conductor \(c_P\), which
is an integer here. Its chain and link losses \(12c_P\) and
\(30c_P\) also satisfy the looser allowances \(24c_P\) and
\(60c_P\) of the flat construction. In the signed-chain
hypothesis the slack \(26rc_P\) suffices for the \(13rc_P\)
in \eqref{eq:plane-gram-reserves}. Let \(A=\phi Z_*\)
and write \(D_b=D_m+D_l+3K_d\). The construction before
the last simplification in the Gram estimate gives
\[
 v([F,\gamma A])\le n_0+D_l+1+\lceil D_m+K_d/2\rceil
                 \le n_0+D_b .
\]
Here the integral conductor data have \(K_d\ge1\).

If \(c_P>D_n+D_b+\kappa_{\rm fl}\), the error in
\eqref{eq:plane-slice-value} is above the depth of the
main term. We therefore obtain
\[
 0\ \le\
 v\left(\frac{p^{\kappa_{\rm fl}}}
              {\prod_i d_i!}[H,\operatorname{Asm}_\rho(A,N_0)]\right)
 \ \le\ D_n+D_b+\kappa_{\rm fl}.
\]
The left inequality uses coefficient integrality and the
factorial deficiency termwise. No error polynomial is
needed for this application of the Gram estimate. The bounds and decomposition below apply uniformly
to all choices of the primitive source.

The detecting value must have a depth bound common to all
primitive source combinations. We choose this bound from the
weight data before choosing a detecting polynomial. Lemma~\ref{lem:lift-plane-symbol}
then lifts the resulting flat symbol to the integral polynomial lattice.
\paragraph{A depth fixed before the choice of a source.}
\label{par:plane-uniform}
\hypertarget{source-plane-return}{}
Let
\[
 L=Nk^2\log_p a.
\]
Here \(a\) denotes the lower bound for
the flat degrees on the tops, as in Section~\ref{par:plane-rank}.
By the inequality on the selected parameters \(a\ge\max(p,4N,k,2)\), and the full
degrees are at most \(2a\). There are at most \(N\) true letters
in the signed model: their positions in the cycle are distinct tops.
Each quotient degree is at most \(2a\) divided by its low radix.
These statements hold before the digits, projections or unit
monomial of the source have been chosen.

\begin{proposition}[An integral lift of the detecting symbol at uniform depth and order]\label{prop:plane-uniform-symbol-return}
Under the hypotheses on the frames and the parameter sizes, let \(X\) be the normal degree bound in the proof and put
\[
 D_X=\lceil49Nk^2\log_p a\rceil,\qquad t_F=\lfloor X\rfloor+Nc_0.
\]
 For each primitive polynomial
in the exact model span there is a detecting flat symbol
of order at most \(t_F\) at depth \(D_X\), realizable
at the selected plane with competitor conductor \(c_P\).
Both \(D_X\) and \(t_F\) are fixed by the weight and size
data before choosing the primitive polynomial.
\end{proposition}

\begin{proof}
\proofstep{The common bounds before choosing a source.}
In \eqref{eq:plane-slice-depth} take for \(D_{\rm no}\) the sum of
the integral projector exponents. At a complement cut of total
degree \(D\), the trace projector bound gives
\(\binom K2\lfloor\log_p(2D+K-1)\rfloor\), interpreted as zero
for a scalar action. The estimate also holds in the factorial
norm by self-adjointness. There are at most \(2N\) such cuts,
since the segment with \(h+r\) listed tops has \(h+r+1\) cuts.
Now
\[
 2D+K\le4Na+k\le a^3;
\]
hence \(D_{\rm no}\le6L\).
Counting carries at each true letter bounds the plane and
complement factorial deficiencies by \(3Nr\log_p a\) and
\(3NK\log_p a\), respectively. Their halves and the
projector loss thus give \(D_n\le9L\).
The full factorial loss \(\kappa_{\rm fl}\) is at most
\(Nk\log_p(2a)\le 2Nk\log_p a\le20L\).

Similarly \(e_d\le4k^2\log_p a\) and
\(K_d\le4Nk\log_p a\), since \(2\sum_i d_i+r\le a^3\).
There are at most \(N\) links and at most \(N\) chain cuts,
so \[
 D_b=D_m+D_l+3K_d\le20L.
\] We may take the integer
\[
                         D_X=\lceil49L\rceil\le50L .
\]
Choose the absolute constant \(C\) in \(D_p\ge CL\) before
taking the host sufficiently large. The conductor hypothesis gives
\[
 c_P\ge c_\star\ge8(D_p+1)>\max\{2D_X,4D_b\},
 \qquad D_X\le50L,\quad D_b\le20L .
\]

For the second reserve in \eqref{eq:plane-gram-reserves}, write
\(D=z_j-a_{j,1}\) and \(P=p^{z_{j+1}-z_j}\). Intervening letters
have distinct distances \(m\ge D+L_1+1\) from \(z_j\).
On a chain cut in segment \(j+1\) the oldest even predecessor
lies at or after \(a_{j,1}\). These two facts give
\[
 b_j\le 2a\sum_{m\ge D+L_1+1}p^{-m}
     \le 3a p^{-D-L_1-1},
 \qquad B_u\ge \frac{a}{p^D P}.
\]
Set \(x=B_u/(Yk^2)\), so \(s_u=\lfloor x\rfloor\).

At the tested letters \(B_u\ge a/Q\) and the host size bound
give \(s_u\ge100r^2 c_P\); in particular \(x\le2s_u\).
The spacing bound \(p^{L_1+1}\ge48rYk^2\) follows from the
host spacing with \(r\le Y\) and \(N\ge1\). Hence
\begin{equation}\label{eq:plane-cut-budget}
 \frac{rb_j}{P}\le \frac{3rB_u}{p^{L_1+1}}
       \le\frac{x}{16}\le\frac{s_u}{8},\qquad
 26r c_P\le\frac{s_u}{4},\qquad 128c_P\le s_u .
\end{equation}
The last two inequalities use \(r\ge2\). Adding the first two costs in \eqref{eq:plane-cut-budget} gives
\[
 rb_j/P+26r c_P\le3s_u/8
\]
on the same cut.

\proofstep{An integral symbol from the model polynomial.}
\label{par:plane-reassembly}
Use the integral model polynomial \(Z_*\) given by the Gram
estimate, and the complement test \(N_0\). Give each monomial
in the model the segment positions of its letters. A block letter
has its intervening position, \(L_{jm}\) and \(T_{jm}\) both have
the outgoing position of that shared row, and \(W_{jm}\) has
its incoming position in the next segment. Keep segment indices
on variables even when their positions in the cycle agree. Place
the complement variables only at the intervening and outgoing
positions. These renamings commute with the signed action on a
prefix and its inverse image. This follows on the variables;
the frame rows are all even and translation by \(n'\) preserves
parity. Coincidence of the label and tangent positions causes
no loss for this assertion.

\proofstep{Prefix orders in the model plane.}
For the renamed plane polynomial the prefix before \(x_u\)
has order at least
\[
                        t_u=T_u-s_u/2-60c_P .
\]
For \(u\le h\) the inverse image is the chain cut. The initial lower bound for the order is \(T_u-s_u/8\).
The two losses in the Gram construction are at most \(3s_u/8\) and
\(24c_P\), respectively. The resulting order is at least
\[
 (T_u-s_u/8)-(3s_u/8+24c_P)=T_u-s_u/2-24c_P\ge t_u .
\]

For \(u>h\) the inverse
image consists of the model link and the outgoing label letters
at earlier rows. These extra letters are vectors. Enlarging the
torus action to them multiplies the previous transform by a
commuting substitution with nonnegative powers, so the lower
bound
\[
 T_u-3s_u/8-60c_P
\]
suffices. For the prefix through an
intervening letter the chain cut at \(u+1\) gives \(t_{u+1}\).
The renamed complement polynomial has upper orders at most
\(k_u+4c_P\) and
\[
 k_{u+1}+4c_P
\]
before and through \(x_u\).

These bounds hold on the full polynomials over \(O\). To use them
segment by segment, take homogeneous components in each
segment letter on the plane side. The complement letter degrees
are already fixed. In one component all the permitted coefficients
in a particular segment, after fixing monomials in the other
segments, belong to the same integral submodule given by the
prefix inequalities. In fact taking such coefficients commutes
with the group of the specified segment. This submodule is
saturated in the finite free module of the given degrees:
its inequalities are vanishing conditions on Taylor or Laurent
coefficients over the fraction field. The quotient is torsion free,
hence free over the DVR.

For the fixed segment degrees, let \(M_j\) be the full
polynomial lattice and \(S_j\subset M_j\) the saturated
submodule cut out by its prefix bounds. Choose an integral
retraction
\[
 r_j:M_j\longrightarrow S_j,\qquad r_j|_{S_j}=\mathrm{id}_{S_j},
\]
possible since the quotient
is free. Fixing all monomials outside \(j\) lands in \(S_j\),
so the endomorphism \(r_j\) in factor \(j\) fixes the component
\(A\). These endomorphisms commute and hence
\[
                   A=(\bigotimes_j r_j)A
                           \in\bigotimes_j S_j.
\]
Applying the same map to a monomial expansion produces
a finite integral sum of products of permitted factors.
The identical argument on the complement introduces no
denominators there either.

\proofstep{Extracting factors with the true degrees.}
Write \(A_j,G_j\) for the factors of a resulting pair. Both have
the designated letter degrees. For an even intervening \(x_u\) take \(M=k_{u+1}+4c_P\),
the upper order of \(G_j\) through the letter. The plane
order before it is at least \(t_u\), and
\[
 t_u-M\ \ge\ m_*+s_u/2-64c_P\ \ge\ m_* .
\]

For an odd intervening \(x_u\) take \(M=k_u+4c_P\)
before the letter. The plane order through it is at least
\(t_{u+1}\) and
\(t_{u+1}-M\ge m_*+s_{u+1}/2-64c_P\ge m_*\).
The lower bounds follow from the respective inequalities
for \(T_u,T_{u+1}\).
The absolute lower bound before each letter is
\(B_u/2-s_u\), since \(T_u\ge B_u/2\).
These are exactly the two parity comparisons and absolute
bound used in the determinant extraction.

For clarity, consider the degree conditions in that extraction.
At an intervening position \(y\), let \(d_{\rm low}(y)\) be the
outer digit degree and \(d_A(y),d_G(y)\) the degrees of these two
factors. The class equation gives
\[
  d_{\rm low}(y)+p^{z_j-y}\bigl(d_A(y)+d_G(y)\bigr)\le\deg(y).
\]
Thus the separate order contributions can be paid out of their sum. At an incoming position
the only letter is \(W_{j-1,m}\), of degree \(b\), and the
complementary degree is zero.

Let \(R_{\rm low}\) be the original low radix and \(R_{\rm seg}\) the incoming segment radix at this letter. Substitution first multiplies the exponent by \(P_{j-1,m}\); passing to true degree then gives
\[
 b\ \longmapsto\ P_{j-1,m}b
      \ \longmapsto\ R_{\rm low}P_{j-1,m}b=R_{\rm seg}b .
\]
This is the contribution from the incoming letter alone. Other terms at the shared letter have nonnegative degree, so if \(d_{\rm coll}\) is its total collapsed plane degree and \(d_{\rm true}\) its true degree, then
\[
 R_{\rm seg}b\le R_{\rm low}d_{\rm coll}\le d_{\rm true}.
\] This argument applies
to every model monomial used, each of which has the prescribed
collapsed degrees. After collecting terms and projecting to a
segment multidegree the same upper bound holds. In particular
the upper allowance
\[
 1.01\deg(y)/p^{z_j-y}
\]
holds in both
places. The extraction proof in Section~\ref{sec:plane-lifts}
only uses this upper inequality when estimating the earlier rows.
It therefore applies also to components which have lost much of
their degree in the incoming frame. This is the one-sided form
of the realization lemma.

\proofstep{Recovering the original low coordinates.}
Multiply each product of segment factors by the original low
label after applying the segment power substitutions. The sum
is \(\operatorname{Asm}_\rho(\phi Z_*,N_0)\): the extra power
on an incoming row is exactly \(P_{j-1,m}\) and the label has
zero support off the tops. Within each fixed component the
result has the required full degrees.

In each true letter the substitutions multiply a segment degree by its radix. Summing over the contributing parts \(b\) therefore gives 
\[
 d_{\rm collapsed}=\sum_b P_b d_b,
\]
 where \(P_b,d_b\) are the radix and degree of the intervening, incoming or outgoing part. Restoring the label then gives the original degree
by the class equation. Its normal degree is also constant: at
an even top add its normal low digits to the low radix times the
complement quotient degree; at an odd top use its plane digits
and plane quotient degree. Sum these numbers over the tops.
If \(S\) is the component of order \(o\), the one-sided realization gives an integral \(P\) with
\[
 P|_{\rm chosen}\equiv w^o S\pmod{w^{o+c_P}},\qquad P|_{\rm competitor}\equiv0\pmod{w^{o+c_P}}
\]
at every competing plane. This proves the symbol membership and its realization on the stated one-sided range.

\proofstep{Keeping one order for all tested generators.}
Finally fix an order before choosing \(H\). The counted generators
have normal degree bounded by
\[
 X=N(\mathcal L a/(pk)+10a p^{-L_1}+kQ)
\]
by
Section~\ref{par:plane-rank}. The common factor has normal degree zero. For a counted generator \(A\) and a product of \(j\le Nc_0\) row operators \(D_1,\ldots,D_j\), the per-operator increment bound gives
\[
 \deg_{\rm normal}(D_1\cdots D_j A)
     \le \lfloor X\rfloor+j\le t_F,\qquad t_F=\lfloor X\rfloor+Nc_0,
\]
whenever the result is nonzero. Taking the span preserves this upper bound. The Fock form pairs only equal exponent tuples.

If \(A=\sum_{o\le t_F}A_o\) is a tested input and
\(Z=\sum_o Z_o\) is the assembled symbol, then the diagonal
Fock pairing gives
\[
                  [A,Z]=\sum_{o\le t_F}[A_o,Z_o] .
\]
The discarded terms have disjoint exponent support from \(A\).
After this deletion the ultrametric inequality selects a
summand with pairing valuation no greater than that of the sum. We have proved the flat symbol estimate
at the same \(D_X,t_F,c_P\) for every primitive polynomial in
the exact model span. \label{cor:plane-return-1}
\end{proof}


Since the evaluation on the sum has valuation at most \(D_X\),
one of its realizable summands of order \(o\le t_F\) has
valuation at most \(D_X\).

\paragraph{Returning to the restored weight.}

The flat symbol has now been lifted with uniform depth. We must compare its detecting covector with one for the original highest weight.

Write \(\lambda\) for the full weight, \(\mu\) for its
flat part and \(A_\eta=A_r(\eta)\) for the normalisation
exponent of a weight \(\eta\). In the notation of the
edge pairing, \(R^\eta_e=w^{-A_\eta}H^\eta_e\),
\(R^{\prime\eta}_e=w^{-A_\eta}(H^\eta_e)^*\), and the
normalised evaluation is \(b^\eta_e(u,T)=w^{2A_\eta}T(u)\).
\begin{established}[Transport maps and the uniform primitive-vector bound]\label{hyp:plane-full-weight-transport}
The operator construction needed for transport has the
following properties. Its maps \(\rho_\beta:V_\lambda\to
V_\mu\) send \(R^\lambda_e\) to \(R^\mu_e\). At every
other edge into the same target they lose at most
\[
 s=rNc_0+2rn'\mathrm{cpp}
\]
powers of \(w\) on the corresponding lattices. On the
target vertex lattices they map \(H^\lambda_+\) into
\(w^{2(A_\lambda-A_\mu)-s}H^\mu_+\). They send products
of a flat generator and a remaining-weight vector into
the allowed span of operator words on that generator.
Finally, for every primitive \(u\in R^\lambda_e\), some
\(\rho_\beta u\) lies outside \(w^{2\kappa+1}R^\mu_e\),
where \(\kappa\) is bounded uniformly before \(u\) is chosen.
\end{established}
\begin{proof}
Take the maps in Proposition~\ref{prop:plane-bounded-weight-operators}. Its divided-power construction at each tensor position preserves the selected edge lattice. The sum of the losses from the other positions is \(rNc_0+2rn'\mathrm{cpp}\). The integral dual estimate there finds an image of a primitive vector outside \(w^{2\kappa+1}R_e^\mu\); its bound on \(\kappa\) depends only on the weight spreads. The construction acts by words on the remaining tensor factors and hence has the required action on the span.
\end{proof}

\begin{proposition}[A flat covector transported to the full weight]\label{prop:plane-restore-covector}
Assume the lattice maps and primitive-vector bound of
Lemma~\ref{hyp:plane-full-weight-transport}. Suppose every
primitive flat vector in the prescribed operator span has a
flat covector of depth \(D_F\), endpoint order \(t_F\), and
competitor conductor \(c_P=c_\star+s\). Then every primitive
\(u\in R^\lambda_e\) in the span of the restored model family
has a full-weight covector whose depth \(D\), endpoint order \(o\) and competitor conductor \(c\) satisfy
\[
 D\le\kappa+D_F,\qquad o\le t_F+s,\qquad c=c_\star.
\]

\end{proposition}

\begin{proof}
\proofstep{Testing a primitive combination in the restored lattice.}
For a chosen primitive combination \(u\), select such a
\(\beta\) and write \(\rho_\beta u=w^a u_1\) with \(u_1\)
primitive and \(0\le a\le2\kappa\). The operator-word
span is over the fraction field, so it contains \(u_1\).
Transport a flat covector \(T\) by
\(\Lambda T=w^{2(A_\mu-A_\lambda)}\rho_\beta^* T\).
For example \(H^\lambda_e\) maps by \(\rho_\beta\)
into \(w^{A_\lambda-A_\mu}H^\mu_e\). Pulling back the
flat lattice \(w^{-A_\mu}(H^\mu_e)^*\) and inserting
the scalar \(w^{2(A_\mu-A_\lambda)}\) gives precisely
the required inclusion in \(R^{\prime\lambda}_e\).

The calculation on the competing edge lattices loses \(s\)
powers of \(w\); the stated vertex inclusion also gives
loss \(s\) at the target vertex. The normalised pairings obey
\[
       b^\lambda_e(u,\Lambda T)=w^a b^\mu_e(u_1,T).
\]
The pairing gains valuation \(a/2\le\kappa\), since
\(w^2=p\). At the vertex and the competing edges the lattice
bounds lose at most \(s\) powers of \(w\), while the chosen
edge retains integrality. Thus the transported test has
depth at most \(\kappa+D_F\), order at most \(t_F+s\)
and conductor at least \(c_P-s=c_\star\).  The model-to-clean
congruence at these depths was checked in
Section~\ref{par:plane-middle-weights}.

\end{proof}

The final step at a single edge is to place the detecting
functional in the target lattice of
Proposition~\ref{prop:residue-plane-neighbor-test}.
The dual-plane substitution accounts for the vanishing at
the other incoming edges.
\paragraph{Dual planes at the target.}
\label{par:plane-dual-edge}
\hypertarget{source-plane-edge-covector}{}
For clarity write the distinguished edge in adapted coordinates as
\(L_-=b W^k\), \(L_+=b D_p^{-1}W^k\), where
\(D_t=\operatorname{diag}(tI_r,I_{k-r})\). Its dual edge goes
from \(L_+^\vee=b^{-t}D_p W^k\) to
\(L_-^\vee=b^{-t}W^k\).

At each Frobenius position the middle
frame on this dual edge is the inverse transpose of the original
middle frame: in the original frame use \(bD_w^{-1}\), and in
the dual use
\[
 b^{-t}D_p D_w^{-1}=b^{-t}D_w,
\]
with the
indicated Frobenius action on \(b\). Distinct incoming edges at
\(L_+\) give distinct outgoing edges at \(L_+^\vee\) by double
duality. There they have distinct residue planes, since a lattice
between a lattice and its \(p^{-1}\)-multiple is determined by
its image in that quotient. For the standard target,
\(bD_p^{-1}\in\mathrm{GL}_k(W)\); it gives the Witt frame
whose inverse transpose is \(b^{-t}D_p\). Thus a single choice
of distinguished frame in the one-sided realization handles all
incoming competitors.

\begin{lemma}[A one-sided plane symbol gives an integral dual-edge test]\label{lem:plane-dual-edge-test}
For the edge from \(bW^k\) to \(bD_p^{-1}W^k\), use
\(bD_w^{-1}\) and \(b^{-t}D_w\) as the respective middle frames.
Let a detecting
symbol of order \(o\) be realized to precision \(o+c_P\)
at the distinguished dual edge and at least that order
at its competitors. If \(D_X<c_P/2\), the resulting flat covector on the original edge has order \(o'\), conductor \(c'\) and depth \(D'\) satisfying
\[
 o'\le o,\qquad c'\ge c_P,\qquad D'\le D_X.
\]

\end{lemma}

\begin{proof}
\proofstep{The dual pairing on the target plane.}
In flat polynomial coordinates use the bilinear form
\(p^{\kappa_{\rm fl}}[P,Q]/\prod_i d_i!\).
It is integral on coefficient lattices in inverse transpose
frames. Indeed the factorial deficiency gives this in the
standard coordinates; the substitution identity
\[
 [P(MX),Q(X)]=[P(X),Q(M^tX)]
\]
follows by writing the
pairing as a constant-coefficient differentiation at zero.
The identity applies to each letter, also when it carries the
inverse transpose representation. The scalar adjustments to the two edge substitutions give 
\[
 w^{A_\mu}w^{A_\mu}=w^{2A_\mu}=p^{A_\mu},
\]
 precisely
the normalized edge pairing for the flat weight \(\mu\), by
Section~\ref{par:plane-middle-weights}.

Take the integral polynomial realizing a detecting symbol
\(\sigma\) of order \(o\), and view it as a flat section in the
dual vertex frame. Divide it by \(w^o\), and pair by this form to
obtain \(T\). At the target its order is at most \(o\), by
integrality before division. In normalized coordinates on the
dual edge its polynomial after division is \(\sigma+w^{c_P}E\)
with integral \(E\); on any other dual edge at the same source
it belongs to \(w^{c_P}O[\mathbf X]\). The inverse transpose
calculation proves \(T\in R_e^{\prime\mu}\) and
\[
 T\in w^{c_P}R_{e'}^{\prime\mu}
\]
on every competing incoming
edge. For any normalized integral input, the model error pairs into
\(w^{c_P}O\), because the factorial rescaling is integral on
the paired edge lattices. Its valuation is at least
\[
 c_P/2,
\]

strictly larger than the detected valuation bounded by \(D_X\).
The valuation of the actual pairing therefore equals that
of the model detecting value.

\end{proof}

\paragraph{Operators for the bounded remaining weights.}
\label{par:plane-rho}
First retain only the top positions in the full weight \(\lambda\).
Write \(\lambda=\mu+\eta\) there, with \(\mu\) the flat weight
and \(\eta\) the remaining weight. Its spread at each top is at
most \(c_0\). We need two facts about the Schur lattices
of Section~\ref{sec:edge-contraction}.

First, for each of the weights \(\xi=\lambda,\mu\) there is a
contravariant identification \(\Psi_\xi:V_\xi^*\longrightarrow
V_\xi\) satisfying, simultaneously in every vertex frame \(g\),
\[
 p^{h_\xi} H_\xi(g^{-t})\ \subseteq \Psi_\xi\bigl(H_\xi(g)^*\bigr)\ \subseteq\ H_\xi(g^{-t}).
\]
In this display \(g\) denotes the tuple of the component frames
and \(H^*\) the integral dual. The same inclusions hold with
\(R_e^{\prime\xi},R_{e^\vee}^\xi\) in the middle and outside,
since the two normalizations multiply the unscaled lattices by
the same \(w^{-A_\xi}\).
Here \(h_\xi\) is the sum over positions of
\(\binom k2\lfloor\log_p(l+k-1)\rfloor\) at positive spreads
\(l\), with zero at spread zero.

\begin{proposition}[Uniform operators for the remaining weights]\label{prop:plane-bounded-weight-operators}
With the bounded top and off-top spreads, the maps
\(\rho_\beta\) of Proposition~\ref{prop:plane-restore-covector}
satisfy its three lattice conditions with
\(s=rNc_0+2rn'\mathrm{cpp}\). Every primitive vector
has an image outside \(w^{2\kappa+1}R_e^\mu\) for some
such map. On the restored model family in the prescribed range of parameters, the target depth \(D\), order \(o\) and conductor \(c\) satisfy
\[
 D\le D_X+\kappa\le54Nk^2\log_p a,\qquad o\le t_F+s,\qquad c=c_\star.
\]

\end{proposition}

\begin{proof}
\proofstep{Normalizing one position in the tensor.}
To verify the normalization at one position, remove the determinant
twist. Take the matrix-polynomial Fock form divided by the
squared Fock norm of the highest polynomial \(v\), the product
of the initial principal minors prescribed by the shape.
This norm is a positive integer. Pairing with \(v\) gives evaluation
at the identity: both vanish outside its weight by row and column
degrees, and both take value one on \(v\). The highest weight
space has dimension one. For a divided-power word applied to
\(v\), transpose the word in the pairing. It follows that its
pairing with any integral polynomial in the Schur lattice is
integral.

Conversely, integrality against the Schur lattice forces
\[
 Q\in O[\mathbf X]
\]
in the matrix coordinates.
Indeed substitute a generic matrix \(M\) in the highest polynomial
by the row action. The coefficients of the resulting polynomial
in the entries of \(M\) are integral elements of the Schur
lattice, by the minor expansion. Pair it with \(Q\) and transpose
the row action. The result is \(Q\) evaluated on the generic
transposed matrix, by the identity evaluation just proved.
Hence every coefficient of \(Q\) is integral, and saturation
of the standard monomial lattice gives the assertion.
The same argument proves nondegeneracy. The divided-power span of
\(v\) contains
\[
 p^{\binom k2\lfloor\log_p(l+k-1)\rfloor}
\]

times the Schur lattice by the trace argument of
Section~\ref{sec:edge-contraction}. Thus the integral dual for the
normalized form lies between the two asserted lattices. Restoring
the determinant twists on inverse transpose frames cancels their
two scalar factors. Tensoring the forms and integral dual bases
adds the exponents and gives \(\Psi_\xi\). Its inverse pairing
is a nonzero scalar times the matrix-polynomial Fock form.
Here \(l+k-1\le a^2\) for both weights at a sufficiently late
host: the flat degree is at most \(2a\), the residual spread
is at most \(c_0\), and the host size bound gives
\(2a+c_0+k\le a^2\). Consequently
\[
 \kappa=h_\lambda+h_\mu\le4L.
\]

Second, for residual sections \(s,\beta\) consider the flat
endomorphism \(U_{s,\beta}\) defined by
\[
  [\theta(U_{s,\beta}F),\theta H]_{\rm mat}
       =[\theta F\cdot s,\theta H\cdot\beta]_{\rm mat}.
\]
The matrix-polynomial form restricted to the image of the flat
model is nondegenerate, by the preceding argument. The flat
insertion \(\theta\) uses the first column or the cofactors.
The maps \(U_{s,\beta}\) are linear in either residual input.
At a single position their span is stable under commutators
with the row operators \(E_{ij}=X_i\partial_j\).
In fact transpose adjunction and the product rule move the
commutator to the two terms applying row operators to \(s\) and
to \(\beta\), with opposite signs and transposed indices.

For cofactors the induced row action on the flat coordinates is
the negative transpose action plus the scalar trace term, by
Jacobi's identity. If the induced operator is \(A+\lambda I\), its commutator on
an endomorphism \(U\) is
\[
 (A+\lambda I)U-U(A+\lambda I)=AU-UA .
\]
Thus the determinant twist introduces no extra row weight.
The residual space is generated by products of flag minors with
at most \(c_0\) columns. Each row occurs at most once in a
column. Thus the commutator weight of a generator
\(U_{s,\beta}\) has each coordinate between \(-c_0\) and \(c_0\):
it is the difference, in one of the two orders, of the two
residual row weights. The determinant twists cancel in this difference.

\proofstep{Operators from the diagonal action.}
For degree \(d\) the endomorphism representation of
\(\operatorname{Sym}^d\) in characteristic zero has one summand for each of the highest weights
\[
 \lambda^{(i)}=(i,0,\ldots,0,-i),\qquad 0\le i\le d.
\]
Indeed \(\operatorname{Sym}^d\otimes
(\operatorname{Sym}^d)^*\) is the bidegree \((d,d)\)
space decomposed by the powers of the contraction in
Section~\ref{sec:edge-transfer}; its harmonic quotients are
the irreducible incidence spaces described there. For summand
\(i\) the row word \(E_{1k}^i\) has that highest weight and
is killed under commutator by every raising operator. Its action on a single monomial gives
\[
 E_{1k}^i X_k^d
 =\frac{d!}{(d-i)!}X_1^i X_k^{\,d-i}\ne0
 \qquad(0\le i\le d),
\]
including the identity at \(i=0\). Words of length at
most \(i\) form a commutator-stable span by the product rule and
the matrix-unit commutator. They therefore contain the entire
summand.

Complete reducibility now implies that the span of the
\(U_{s,\beta}\) lies in words of length at most \(c_0\). Indeed
any summand present has its highest weight present, and the
weight bound just proved gives \(i\le c_0\).
At several positions expand the residual sections in pure
tensors and use factorization of the Fock form. Each term gives
at most \(c_0\) row operators at each top. Changes of flat frame
preserve this word space by conjugation. The word bound also holds when the flat degree is smaller than the allowed word length.

For \(\beta\in R_{e^\vee}^{\eta}\), write \(m_\beta\) for
multiplication by that residual section. Define
\[
   \Lambda_\beta T
       =p^{h_\lambda}\Psi_\lambda^{-1}
                      m_\beta\Psi_\mu T,\qquad
   T(\rho_\beta u)=p^{A_\eta}(\Lambda_\beta T)(u).
\]
The second equality defines a linear map by finite-dimensional
duality. The edge normalizations give
\[
 b_e^\mu(\rho_\beta u,T)=b_e^\lambda(u,\Lambda_\beta T).
\]
For \(T\in R_e^{\prime\mu}\), multiplication on the dual edge and
the inclusions for \(\Psi_\xi\) give
\[
 \Psi_\mu T\in R_{e^\vee}^\mu,\qquad
 m_\beta\Psi_\mu T\in R_{e^\vee}^\lambda,\qquad
 p^{h_\lambda}R_{e^\vee}^\lambda\subseteq
                              \Psi_\lambda(R_e^{\prime\lambda}).
\]
Thus \(\Lambda_\beta T\in R_e^{\prime\lambda}\).
In unscaled dual bases the two lattice factors \(w^{-A}\) cancel
the pairing factor \(w^{2A}\). The preceding adjoint identity
and this perfectness imply
\[
 \rho_\beta R_e^\lambda\subseteq R_e^\mu.
\]

On an input
\(u=\theta F\cdot s\), the definition of the two forms shows
that its flat output is a scalar multiple of
\(U_{s,\beta}F\) in polynomial coordinates. This proves
membership in the exact operator span.

\proofstep{Reservations for the remaining weights.}
Put \(s_T=rNc_0\). Multiplying any \(\beta\) from the chosen
dual edge by \(w^{s_T}\) puts it in the dual-edge source vertex
lattice. To see this, at each position the edge normalization
subtracts the top \(r\)-row weight sum. The largest possible
difference between that sum and a sum on any
\[
 r
\]
rows is at
most \(r\) times the weight spread, by the tableau weight bound.
This proves the claim in an adapted diagonal frame, including
determinant twists, and hence in every frame. Every outgoing normalized dual edge contains the vertex lattice.
Hence the cost outside the distinguished edge is at most \(s_T\).
For the target vertex, test against \(T\in(H_+^\mu)^*\). The vertex
inclusions and \(p=w^2\) give
\begin{equation}\label{eq:edge-target-adjoint}
 \begin{aligned}
 w^{s_T}\Lambda_\beta T &\in(H_+^\lambda)^*,\\
 T(\rho_\beta u)&=w^{2A_\eta-s_T}
          (w^{s_T}\Lambda_\beta T)(u)\in w^{2A_\eta-s_T}O
                      \qquad(u\in H_+^\lambda).
 \end{aligned}
\end{equation}
Testing every integral dual functional in \eqref{eq:edge-target-adjoint} yields
\(\rho_\beta H_+^\lambda\subseteq w^{2A_\eta-s_T}H_+^\mu\).

There is also a lower bound for this family of maps. The full
lattice on the dual edge is the integral span of products
\(y\beta\) from the flat and residual edge lattices. In a standard
tableau product, the column heights of the sum shape are the
union, with multiplicities, of those of the two shapes. Split
the columns accordingly and take a tensor of these factorizations
over the positions. The frame action, determinant twist and
edge powers all multiply. Thus for every full dual integral
functional \(G\), \(\Psi_\lambda G\) is an integral sum of these
products.

If
\[
 \Psi_\lambda G=\sum_i y_i\beta_i,
\]
take
\(T_i=\Psi_\mu^{-1}(p^{h_\mu}y_i)\), which is flat
dual integral. The multiplier in \(\Lambda_\beta\) then gives
\[
 p^\kappa G=\sum_i\Lambda_{\beta_i} T_i
\]
with
\(\kappa=h_\lambda+h_\mu\). For the normalized pairings,
\[
 p^\kappa b_e^\lambda(u,G)
       =\sum_i b_e^\mu(\rho_{\beta_i}u,T_i).
\]
Suppose \(\rho_\beta u\in w^{2\kappa+1}R_e^\mu\) for every
integral choice of \(\beta\). Testing the displayed sum against the
integral flat duals then gives, for every \(G\in R_e^{\prime\lambda}\),
\[
 p^\kappa b_e^\lambda(u,G)\in w^{2\kappa+1}O
 \quad\Longrightarrow\quad b_e^\lambda(u,G)\in wO
 \quad(p^\kappa=w^{2\kappa}).
\]
Perfectness gives \(u\in w R_e^\lambda\), a contradiction for a
primitive \(u\). Thus at least one residual test detects it.

\proofstep{Restoring the non-top tensor factor.}
Restore the non-top tensor factor, of weight \(\xi\), by taking
all contractions on it with \(p^{A_\xi}G\),
\(G\in R_e^{\prime\xi}\), before applying the top maps.
In a pair of normalized dual bases these contractions are
ordinary coordinate extractions. Some extraction of a primitive
tensor remains primitive. Thus this step adds no depth. The
spread off the tops is at most \(2\mathrm{cpp}\), so
\[
 w^{2rn'\mathrm{cpp}}G
\]
belongs to the target dual vertex
lattice and hence to all competing incoming dual edge lattices.

For this inclusion, the passage from an intermediate to a target
frame expands the first \(r\) directions once more by \(w^{-1}\);
a coordinate functional then loses at most the difference of two
\(r\)-row weight sums. On the target itself the contraction
carries the additional scalar \(p^{A_\xi}\). Adding the competing-edge loss from these positions to the
top contribution gives the budget
\[
 s=s_T+2rn'\mathrm{cpp},\qquad
 2A_\eta+2A_\xi-s=2(A_\lambda-A_\mu)-s .
\]
Here the right-hand exponent is the resulting target-vertex exponent.
The flat
degree is zero at every omitted position, so this contraction
retains the top word allowance.

Apply the adjoint calculation to the primitive restored model
combination. The three resulting bounds are
\begin{equation}\label{eq:restored-covector-budget}
 \begin{aligned}
 \text{endpoint order}&\le t_F+s\le\lceil t_{\rm II}\rceil,\\
 \text{depth}&\le D_X+\kappa\le54Nk^2\log_p a\le D_{\rm II},\\
 \text{competing conductor}&\ge c_P-s=c_\star,
                    \qquad c_\star=8(D_{\rm II}+1).
 \end{aligned}
\end{equation}
The last depth comparison and the endpoint bound in
\eqref{eq:restored-covector-budget} are the host choices of Section~\ref{sec:host-edge-tests},
with its absolute coefficient taken sufficiently large.
On the source side cleaning
the normal-degree components of the original flat generators
costs at most \(\lfloor X\rfloor\) at the source, by the
one-sided lift and division by the order of each component.
Multiplication by a remaining-weight section from the chosen
edge costs at most \(s\) at that vertex, by the same spread
estimate. Thus the source order also fits. With \(c'=c_P\)
one has \[
 c'\ge2(D_X+\kappa+1),\qquad c'\ge c_\star+s.
\]
Section~\ref{par:plane-middle-weights} transfers the pairing
to the clean combination and gives its competing source
conductors. This completes the single-edge estimate with the
order and depth fixed on the host and the stated quantifier over
primitive combinations. \label{cor:plane-edge-covector-1}
\end{proof}



\paragraph{From one plane to a full target.}
\begin{proof}[Conclusion of the proof of Proposition~\ref{prop:residue-plane-neighbor-test}]

Write \(d=\dim V\), and put \(a_0=a-c_G\).
The preceding construction on every edge into the standard
vertex gives \(m_e\) clean source vectors, with
\[
 m_e\ \ge\
 (1-CNr^2 k/Y_2)(1-k c_0/a_0)^N
                d\, q^{-r(k-r)} .
\]
For each combination within this block having a unit
coefficient, the covector of Section~\ref{par:plane-rho} evaluates
with valuation at most \(D_{\rm II}\) and satisfies the conductor
condition at every competing incoming edge. Both endpoint orders
are bounded by \(\lceil t_{\rm II}\rceil\).

Take incoming edges for which the codimension-\(r\)
subspace of \(L_+/pL_+\) is a graph over the last \(k-r\)
coordinates. All \(q^{r(k-r)}\) graphs give distinct lattices.
An arbitrary target is reached by a basis change over \(K_0\)
and its Frobenius twists. For a primitive combination of the
vectors from all these edges, choose a block containing a unit
coefficient and its detecting covector. The other blocks contribute
terms of valuation at least
\[
 c/2>D_{\rm II},
\]
by integrality
of the edge pairings and \(c=8(D_{\rm II}+1)\). They preserve
the valuation of the detected value.

\proofstep{Combining the independent directions at the target.}
If
\[
 M=\sum_e m_e,
\]
the evaluation image of the target lattice
in \(O^M\) contains \(p^{D_{\rm II}}O^M\). Indeed in Smith
coordinates any missing direction or any elementary divisor of
larger valuation gives a primitive linear combination taking
only values of larger valuation, contrary to what was just
proved. Choosing a preimage of each scaled basis vector gives
the integral right inverse at depth \(D_{\rm II}\). Finally
\[
 (1-kc_0/a_0)^N\ge 1-Nkc_0/a_0
\]
and \(c_0/a_0=O(1/X')\).
The resulting deficit is at most \(\min(1,C'Nr^2 k/Y_2)\)
after increasing the constant, including the range where this
bound is one. At \(Y_2=x^{26}\) this deficit is \(O(Nx^{-12})\)
with the host bounds of Section~\ref{sec:host-edge-tests}.
The finite extension argument there preserves this estimate
on the selected good cycles. This proves the smaller-prime input
at the same host as the larger-prime estimate.
 \label{cor:plane-lifts-1}
\end{proof}

